\documentclass{report}

\overfullrule=0pt
\sloppy

% ==================== PACKAGES ====================
\usepackage[a4paper,
            left=3cm,
            right=2.5cm,
            top=3cm,
            bottom=3cm]{geometry}

\usepackage{setspace}


\usepackage{times}
\usepackage{titlesec}
\usepackage{tocloft}
\usepackage{graphicx}
\usepackage{caption}
\usepackage{listings}
\usepackage{float}
\usepackage{booktabs}
\usepackage{subcaption}
\usepackage{amsmath, amssymb}
\usepackage[hidelinks]{hyperref}

\usepackage{eso-pic}
\usepackage{transparent}

\AddToShipoutPictureBG{%
    \ifnum\value{page}>1
        \AtPageCenter{%
            \hspace*{-0.35\paperwidth}%
            \raisebox{-0.08\paperheight}{%
                \transparent{0.15}\includegraphics[width=0.7\paperwidth]{bau_logo.png}%
            }%
        }%
    \fi
}


% ==================== CHAPTER FORMAT ====================
\titleformat{\chapter}
  {\normalfont\bfseries\centering}
  {CHAPTER \thechapter}
  {1em}
  {\uppercase}

% ==================== DOCUMENT ====================
\begin{document}

% ==================== COVER PAGE ====================
\begin{titlepage}
\begin{center}
T.C.\\
\textbf{BAHÇEŞEHİR UNIVERSITY}\\
GRADUATE SCHOOL OF NATURAL AND APPLIED SCIENCES\\
DEPARTMENT OF BIG DATA ANALYTICS AND MANAGEMENT\\[3.5cm]

\textbf{\uppercase{System-Level Flight Delay Prediction Using 
Multi-Source Data and Machine Learning: A Case Study at Istanbul Airport}}\\[1.5cm]

{\transparent{0.15}\includegraphics[width=0.65\textwidth]{bau_logo.png}}\\[1.5cm]

\textbf{MASTER'S THESIS}\\[2cm]

\textbf{Sheraz Lyna Touhouche}\\[3cm]

ISTANBUL\\
2026
\end{center}
\end{titlepage}

% ==================== SECOND TITLE PAGE ====================
\begin{titlepage}
\begin{center}
T.C.\\
\textbf{BAHÇEŞEHİR UNIVERSITY}\\
GRADUATE SCHOOL OF NATURAL AND APPLIED SCIENCES\\
DEPARTMENT OF BIG DATA ANALYTICS AND MANAGEMENT\\[3.5cm]

\textbf{\uppercase{System-Level Flight Delay Prediction Using 
Multi-Source Data and Machine Learning: A Case Study at Istanbul Airport}}\\[1.5cm]

{\transparent{0.15}\includegraphics[width=0.65\textwidth]{bau_logo.png}}\\[1.5cm]

\textbf{MASTER'S THESIS}\\[2cm]

\textbf{Sheraz Lyna Touhouche}\\[3cm]

ISTANBUL\\
2026
\end{center}
\end{titlepage} 

% ==================== FRONT MATTER ====================
\pagenumbering{roman}
\setcounter{page}{3}

% ==================== ETHICAL CONDUCT ====================
\chapter*{ETHICAL CONDUCT}
I hereby declare that all information in this document has been obtained and
presented in accordance with academic rules and ethical conduct. I also declare
that, as required by these rules and conduct, I have fully cited and referenced all
material and results that are not original to this work.

\vspace{2cm}
Name, Last Name: Sheraz Lyna Touhouche\\
Signature:

\addcontentsline{toc}{chapter}{ETHICAL CONDUCT}

% ==================== ABSTRACT ====================
\chapter*{ABSTRACT}
\addcontentsline{toc}{chapter}{ABSTRACT}






\begin{center}
\textbf{System-Level Flight Delay Prediction Using Multi-Source Data and Machine Learning: A Case Study at Istanbul Airport}\\
Sheraz Lyna Touhouche\\
Master’s Program in Big Data Analytics and Management\\
Supervisor: Dr. Tevfik Aytekin\\
June 2026, 71 pages


\end{center}





\vspace{1cm}

\vspace{1em}

Flight delays remain a persistent operational challenge in modern air transportation, affecting airlines, airports, and passengers. While machine learning approaches have been widely applied to delay prediction, most studies rely on random train–test splits, which fail to account for temporal non-stationarity and may overestimate real-world performance.

This thesis proposes a machine learning framework for predicting binary departure delays ($\geq$15 minutes) at Istanbul Airport using multi-source operational data collected between October 2025 and February 2026. The dataset integrates flight records, hourly meteorological data, and surface traffic indicators, resulting in a final sample of 9,527 flights. Logistic Regression, Random Forest, and XGBoost models are evaluated under both random split and temporally consistent validation settings, with strict prevention of temporal leakage.

Under random split evaluation, the tuned XGBoost model achieves a ROC-AUC of 0.903, surpassing the reported benchmark of 0.88. SHAP analysis reveals that historical delay propagation features account for  40.6\% percent of predictive importance and meteorological variables account for 33.9 percent, collectively establishing airline operational history and weather conditions as the co-dominant drivers of departure delay prediction.

These findings demonstrate that static predictive models are insufficient for operational deployment and emphasize the need for adaptive, time-aware modeling strategies in aviation delay prediction. Cost-sensitive evaluation further highlights the importance of aligning prediction thresholds with operational priorities under asymmetric risk conditions.

\vspace{1cm}


\textbf{Keywords:} flight delay prediction, machine learning, temporal validation, concept drift, delay propagation

% ==================== ÖZET ====================
\chapter*{ÖZET}
\addcontentsline{toc}{chapter}{ÖZET}

\begin{center}
\textbf{Çok Kaynaklı Veri ve Makine Öğrenmesi Kullanılarak Sistem Düzeyinde Uçuş Gecikmesi Tahmini: İstanbul Havalimanı Örnek Çalışması}
\end{center}

\vspace{1cm}

Uçuş gecikmeleri, havayollarını, havalimanlarını ve yolcuları etkileyen kalıcı bir operasyonel sorun olmaya devam etmektedir. Makine öğrenmesi tabanlı gecikme tahmin çalışmaları geniş bir literatüre sahip olsa da mevcut çalışmaların büyük çoğunluğu, zamansal dönemleri yapay olarak karıştıran rastgele eğitim-test bölmelerine dayanmakta ve gerçek dünya dağıtım performansını muhtemelen aşırı tahmin etmektedir.

Bu tez, Ekim 2025 ile Şubat 2026 arasında toplanan çok kaynaklı operasyonel veriler kullanılarak İstanbul Havalimanı'nda ikili kalkış gecikmesi sonuçlarını ($\geq$ 15 dakika) tahmin etmeye yönelik bir makine öğrenmesi çerçevesi geliştirmektedir. AviationStack API'sinden elde edilen uçuş kayıtları, OpenWeather API'den alınan saatlik meteorolojik gözlemler ve İstanbul Büyükşehir Belediyesi'nden temin edilen yüzey trafik göstergeleri olmak üzere üç veri kaynağı entegre edilmiş; nihai veri kümesi 9.527 uçuştan oluşmaktadır. Lojistik Regresyon, Rastgele Orman ve XGBoost modelleri hem rastgele bölme hem de zamansal doğrulama protokolleri kapsamında değerlendirilmiştir.

Rastgele bölme değerlendirmesinde XGBoost Ayarlı modeli 0{,}903 ROC-AUC değeri elde ederek yayımlanmış referans değerini (\%0{,}88) aşmıştır.SHAP analizi, tarihsel gecikme yayılım özelliklerinin tahmin öneminin 40,6'sını ve meteorolojik değişkenlerin 33,9'unu oluşturduğunu ortaya koymuştur. Bununla birlikte, Ekim--Kasım verileriyle eğitilen modeller Ocak--Şubat döneminde test edildiğinde performans neredeyse rastgele düzeylere (ROC-AUC = 0{,}474) gerilemiştir. Bu çöküş, gecikme yayılım dinamiklerindeki mevsimsel rejim kaymasına bağlanmaktadır: Sonbaharda +7{,}07 puan olan gecikme bulaşma etkisi, kışın $-7{,}78$ puana dönmekte ve toplamda 14{,}84 puanlık bir salınım oluşturmaktadır. Ayrıca yüzey yol trafiğinin ihmal edilebilir düzeyde tahmin değeri taşıdığı kanıtlanmıştır. Bu bulgular, statik modellerin operasyonel kullanım için yetersiz olduğunu ve uyarlanabilir yeniden eğitim stratejilerinin zorunlu olduğunu göstermektedir.

\vspace{1cm}

\noindent \textbf{Anahtar Kelimeler:} uçuş gecikmesi tahmini, makine öğrenmesi, zamansal doğrulama, kavram kayması, gecikme yayılımı
 
\newpage
% ==================== DEDICATION ====================
\thispagestyle{empty}
\vspace*{5cm}
\begin{center}
\textit{To my family, for their unwavering support throughout this journey.}
\end{center}
\newpage

% ==================== ACKNOWLEDGEMENTS ====================
\chapter*{ACKNOWLEDGEMENTS}
\addcontentsline{toc}{chapter}{ACKNOWLEDGEMENTS}


Completing this thesis has been a journey filled with challenges, learning, and growth—one that I could not have undertaken alone.

First and foremost, I want to express my deepest gratitude to my supervisor, Assoc. Prof. Tevfik Aytekin, for his guidance, patience, and unwavering support throughout this research. His insightful feedback and encouragement helped me navigate the most difficult moments of this work. I am grateful not only for his expertise, but for his kindness and belief in my abilities, even when I doubted myself.

I am equally grateful to my academic advisor, Ece Gelal Soyak, who has been a constant presence throughout my academic journey. Her advice, guidance, and genuine care have shaped not only this thesis, but also my development as a researcher and as a person. Thank you for always being there when I needed direction, for listening, and for believing in me.

I would like to thank Bahçeşehir University for providing the resources and environment that made this research possible. The support of the faculty and staff has been invaluable.

This work would not have been possible without the incredible tools built by the open-source community. To the developers of scikit-learn, XGBoost, pandas, and the broader Python ecosystem: thank you for making advanced machine learning accessible to researchers like me. Your work has had a profound impact on this thesis and on my education.

I am also grateful to the aviation data providers whose publicly available APIs and datasets enabled the multi-source data collection at the heart of this research. Access to real-world operational data transformed this thesis from theory into practice.

To my family: words cannot express how much your love and support have meant to me. Thank you for standing by me through every late night, every moment of frustration, and every small victory. You believed in me even when the path ahead seemed unclear, and that belief carried me through.


Finally, to everyone who contributed to this work in ways both large and small: thank you. This thesis is a reflection not just of my efforts, but of the collective support, wisdom, and kindness of all those who have walked alongside me.

\vspace{1cm}
\begin{flushright}
\textit{Istanbul, May 2026}
\end{flushright}

\newpage

% ==================== CONTENTS ====================
\tableofcontents
\listoftables
\listoffigures

% ==================== MAIN CONTENT ====================
\clearpage
\pagenumbering{arabic}

\chapter{Introduction}


Air transportation has become an essential component of modern society, playing a critical role in economic development, tourism, and global connectivity. As global air traffic continues to grow, airports must manage increasingly complex operations under limited infrastructure constraints. Flight delays have emerged as a persistent and costly operational challenge, affecting airlines, airport authorities, and passengers alike. At major hub airports in particular, delays rarely result from a single isolated factor; rather, they arise from the interaction of weather conditions, aircraft rotations, crew constraints, and cascading operational disruptions. Understanding and predicting these delays requires a system-level perspective that accounts for the interconnected nature of airport operations. \\

Recent advances in data availability and machine learning have opened new possibilities for modeling such complex systems. However, most existing delay prediction studies rely on random train-test splits that artificially intermix temporal periods, likely overestimating real-world performance. The present research addresses this gap through rigorous temporal validation and multi-source data integration at Istanbul Airport. \\

Istanbul Airport (IATA: IST), which opened in April 2019 following the phased 
closure of Atatürk Airport, represents one of the most operationally complex hub 
airports in Europe. Designed to accommodate up to 90 million passengers annually 
in its initial phase, the airport serves as the primary international hub for 
Turkish Airlines, which accounts for approximately 51\% of total departure traffic. 
In 2023, Istanbul Airport became the busiest airport in Europe by number of aircraft 
movements, handling over 620,000 flights and connecting passengers to more than 120 
countries (EUROCONTROL, 2024). This concentration of hub operations --- in which a 
single dominant carrier executes dense rotation schedules across shared gates, crew 
bases, and aircraft fleets --- creates conditions in which operational disruptions 
propagate rapidly through the broader schedule. A single delayed inbound rotation 
can cascade across multiple subsequent departures within hours, making Istanbul 
Airport a particularly demanding and consequential environment for delay prediction 
research.

\section{Theoretical Framework}

This study is grounded in two complementary theoretical traditions. The first is 
systems theory as applied to air transportation networks, which conceptualises 
airports as complex adaptive systems in which local disruptions propagate through 
shared operational resources to produce non-linear, emergent network-wide effects. 
Within this framework, an individual flight delay is not merely a scheduling failure 
but a perturbation in a tightly coupled operational system: a delayed inbound 
aircraft occupies a gate past its scheduled departure time, preventing the subsequent 
outbound flight from boarding; the inbound crew may simultaneously approach duty-time 
limits, necessitating reserve crew activation; and connecting passengers from the 
delayed arrival may trigger holds or misconnections on onward flights. These cascading 
interactions mean that the delay probability of a given flight is not independent of 
the operational state of surrounding flights, and predictive models that treat each 
flight as an isolated observation systematically underestimate the predictability 
available from network-state features.

The second theoretical tradition is the concept drift literature in machine learning, 
which recognises that the statistical relationship between input features and a target 
variable may shift over time as the underlying data-generating process evolves. In 
operational aviation, concept drift arises from seasonal changes in traffic density, 
weather regime transitions, airline schedule revisions, and regulatory adjustments 
that alter the causal structure of delay formation. A model trained on autumn 
operational patterns learns feature-target relationships that may not generalise to 
winter operations, even if both periods exhibit similar aggregate delay rates, because 
the mechanisms driving delays --- and the features that most reliably encode those 
mechanisms --- may differ fundamentally across seasonal regimes. Recognition of concept 
drift has motivated the development of evaluation protocols that test model performance 
on chronologically separated test periods rather than randomly sampled holdout sets, a 
methodological principle that the present study applies rigorously.

Together, these two theoretical frameworks provide the conceptual architecture for 
the present research design. Systems theory motivates the multi-source feature 
integration strategy, the incorporation of historical delay propagation features, and 
the delay contagion analysis conducted in Chapter~4. The concept drift perspective 
motivates the dual evaluation protocol, the temporal stability analysis, and the 
study's central conclusion that static models trained on a single seasonal period are 
fundamentally insufficient for sustained operational deployment. The empirical findings 
confirm both theoretical predictions: delay propagation is a primary predictive signal 
(40.6\% of SHAP importance), and the feature-target relationships learned during 
autumn training reverse sign in winter, producing catastrophic performance collapse 
under temporal validation.

\section{Statement of the Problem}

Existing flight delay prediction models remain limited in their ability to 
accurately represent real-world operational conditions. Traditional approaches 
rely heavily on scheduled flight information and basic meteorological variables, 
treating flights as independent observations and ignoring system-level dynamics. 
These models implicitly assume that the delay probability of a given flight can 
be estimated solely from characteristics observable before departure --- such as 
scheduled time, destination, and current weather --- without reference to the 
operational state of the broader airport system. This assumption is violated at 
high-traffic hub airports, where gate contention, crew positioning requirements, 
and aircraft rotation schedules create tight interdependencies among flights 
departing within the same operational window. Ignoring these dependencies means 
that standard models discard a substantial portion of the predictive information 
embedded in real-time operational state.

A further limitation concerns evaluation methodology. The predominance of random 
train-test splits in the published literature artificially intermixes temporal 
periods, allowing models to be evaluated on historical observations drawn from 
the same calendar months as their training data. In the presence of seasonal 
non-stationarity --- which is intrinsic to aviation operations --- this intermixing 
produces performance estimates that reflect interpolation within a known distribution 
rather than genuine forecasting across seasonal boundaries. Published benchmarks, 
including the widely cited 0.88 ROC-AUC of Gui et al.\ (2019), were generated under 
random-split evaluation conditions that do not simulate deployment-realistic 
forecasting. The degree to which these benchmarks overestimate true operational 
performance has not been systematically quantified in the existing literature.

At Istanbul Airport specifically, these limitations are particularly consequential. 
As one of Europe's busiest airports by aircraft movements, Istanbul Airport (IATA: 
IST, ICAO: LTFM) operates under conditions of extreme hub concentration, with Turkish 
Airlines executing over 51\% of total departures through a schedule organised around 
dense morning and evening departure banks. This configuration maximises network 
connectivity but creates severe propagation risk: a disruption to early-morning 
inbound rotations can cascade through the entire departure schedule before afternoon 
recovery protocols take effect. Standard prediction models that ignore propagation 
dynamics, rely on research conducted at lower-congestion airports, or evaluate 
performance without temporal validation are poorly equipped to inform operational 
decision-making at IST. The present study directly addresses each of these gaps 
through a purpose-designed multi-source dataset, a strict temporal validation 
protocol, and a comprehensive propagation analysis conducted across seasonal periods.

\section{Purpose of the Study}

The purpose of this study is to improve the realism and predictive accuracy of short-term flight delay forecasting at a major international hub airport by integrating multiple operational data sources and adopting rigorous temporal validation. Rather than focusing solely on scheduled flight and weather data, the proposed approach incorporates historical delay behavior, surface traffic conditions, and airspace congestion. By explicitly testing model performance across seasonal periods, the study also aims to characterize the temporal stability of learned delay patterns and assess the practical deployment readiness of machine learning models for aviation operations.

\section{Hypotheses / Research Questions}

This research seeks to answer the following questions:

\begin{itemize}
    \item \textbf{RQ1.} What is the relative influence of operational and environmental factors --- including weather, historical delay patterns, and surface traffic --- on flight delay outcomes at Istanbul Airport?
    
    \item \textbf{RQ2.} How do machine learning models compare in their predictive performance when multi-source operational data is integrated in a hub airport environment?
    
    \item \textbf{RQ3.} Do the predictive relationships learned from one seasonal period generalize to future seasonal periods, or do delay dynamics exhibit non-stationarity over time?
    
    \item \textbf{RQ4.} Does surface road traffic congestion provide meaningful predictive signal for flight delays at a major international hub airport?
\end{itemize}

\section{Significance of the Study}

This research makes several contributions to the flight delay prediction 
literature and to practical aviation operations management. The study's primary 
methodological contribution is a comprehensive temporal validation framework that 
directly exposes the performance gap between random-split and deployment-realistic 
evaluation conditions. By training on October--November 2025 data and evaluating 
on January--February 2026 data, the study documents a 0.429 ROC-AUC gap 
attributable entirely to evaluation methodology rather than model quality --- 
a finding that challenges the validity of the majority of published performance 
benchmarks in the flight delay prediction literature.

The study's primary empirical contribution is the discovery and quantification 
of a seasonal regime shift in delay propagation mechanisms. The reversal of the 
delay contagion effect from positive in autumn ($+7.07$ percentage points) to 
negative in winter ($-7.78$ percentage points) reveals that delay generation 
mechanisms fundamentally differ across seasonal operational regimes. This finding 
invalidates the common assumption that a model trained on one season will 
generalise to another and provides direct empirical evidence of the concept drift 
predicted by machine learning theory in operational aviation contexts. The total 
swing of 14.84 percentage points across seasons represents one of the most 
precisely characterised instances of seasonal non-stationarity in the published 
single-airport delay prediction literature.

A further significant contribution concerns the hypothesis that surface road 
traffic congestion contributes meaningful predictive signal for departure delays 
at a major international hub airport. Despite the theoretical plausibility of 
landside accessibility effects on crew and passenger punctuality, no prior 
published study had empirically tested this hypothesis in a prediction model. 
The conclusive negative finding reported here --- mean absolute correlation 
$|r| = 0.03$, SHAP importance 8.9\%, AUC contribution $-0.002$ percentage 
points --- constitutes a genuine scientific contribution by redirecting future 
research toward airside operational constraints rather than landside traffic 
factors. Transparent documentation of this hypothesis rejection validates the 
scientific integrity of the study and prevents unproductive replication of the 
traffic integration approach in future aviation machine learning research.

From a practical standpoint, the findings offer actionable guidance for airlines 
and airport operators at Istanbul Airport and comparable hub environments. The identification of carrier-level propagation heterogeneity ranging from $+30.6$~pp to $+59.8$~pp percentage points enables targeted intervention strategies calibrated to individual airline characteristics. The cost-sensitive threshold analysis 
demonstrates that operational costs can be reduced by approximately 20\% by 
shifting the decision threshold from the conventional default of 0.50 to the 
cost-optimal value of $\tau^* = 0.12$, providing a concrete deployment 
recommendation derived from asymmetric cost modelling with domain-validated 
cost ratios.

\section{Definitions}

The following key terms are used throughout this thesis:

\textbf{Flight delay} refers to the difference between the scheduled and actual departure time of a flight. In this study, a departure delay of 15 minutes or more constitutes a significant delay, consistent with international aviation reporting standards (EUROCONTROL CODA; U.S. DOT 14 CFR Part 234).

\textbf{Delay propagation} denotes the phenomenon whereby a delay affecting one flight increases the probability of delays in subsequent flights through shared operational resources such as gates, crew, and aircraft rotations.

\textbf{Concept drift} refers to the change over time in the statistical relationship between input features and the target variable, which causes a trained model's predictive performance to degrade in deployment.

\textbf{Machine learning} refers to a branch of artificial intelligence that enables systems to learn patterns and relationships from data without being explicitly programmed.

\textbf{Airspace congestion} denotes the level of aircraft density within a defined airspace region surrounding an airport.

\textbf{Big data} describes datasets characterized by high volume, variety, and velocity that require specialized processing methods.



\chapter{Literature Review}

\section{System-Level Perspectives on Flight Delays}

Flight delays represent one of the most persistent and economically consequential challenges in contemporary air transportation. Their effects extend well beyond individual flights, creating cascading inefficiencies that affect airline scheduling, airport capacity utilization, air traffic flow management, and passenger experience. According to industry estimates, flight delays cost the aviation sector tens of billions of dollars annually in direct and indirect costs, spanning fuel burn, crew overtime, aircraft repositioning, and passenger compensation. As global air traffic demand continues to grow, the systemic consequences of congestion and disruption have become increasingly severe, motivating a sustained body of research aimed at understanding and predicting delay dynamics.

Early investigations into flight delays approached the problem from a largely local perspective, examining individual flights or airport pairs in isolation and focusing on scheduled versus actual departure times to identify delay patterns. While these studies produced useful descriptive statistics, they were limited in their ability to explain broader system behavior. Over time, researchers increasingly recognized that flight delays are rarely isolated events. Rather, they emerge from interactions among multiple interconnected subsystems --- including runway capacity, gate availability, crew and aircraft rotations, air traffic control flow management, and weather conditions --- all operating simultaneously under shared resource constraints.

From a system-level standpoint, air transportation networks exhibit properties commonly associated with complex adaptive systems. Airports, airspace sectors, and flight routes form interconnected networks in which local disruptions propagate spatially and temporally. Congestion at a single hub airport can trigger aircraft rotation delays, crew misconnections, and gate conflicts that ripple through a carrier's entire schedule.Tang et al. (2021) synthesize a broad range of studies adopting a complex network perspective and demonstrate that delay dynamics are strongly influenced by network topology, hub concentration, and congestion intensity. Their review highlights how delays propagate through aviation networks in ways analogous to cascade failures or epidemic processes, underscoring that delay formation is not solely a function of local conditions but rather an emergent property of the broader system. This perspective is particularly relevant at major hub airports such as Istanbul Airport (IST), where Turkish Airlines' dominant hub operations mean that operational disruptions at a single terminal can produce widespread downstream effects.

The economic consequences of flight delays extend well beyond direct operational
costs. According to EUROCONTROL's Central Office for Delay Analysis (2024),
flight delays across European airspace generate annual direct costs exceeding
\euro13 billion, encompassing airline fuel overconsumption, crew overtime, and
passenger compensation, with indirect costs to passengers and connecting ground
transportation adding a further estimated \euro8--12 billion. At the airport
level, delay propagation creates amplifying effects on gate utilisation and
turnaround efficiency: a 15-minute ground delay compounding through three
downstream rotations can produce an aggregate schedule disruption of 45--60
minutes before airlines activate recovery procedures such as speed-up climb
profiles or reduced ground time buffers. These systemic amplification dynamics
underscore why individual-flight prediction models that operate without
system-level features consistently underperform models that capture the
prevailing network state (Tang et al., 2021).

The operational complexity of major hub airports introduces particular
challenges for delay prediction research. Hubs concentrate traffic through
high-frequency rotations of a dominant carrier, creating dense dependency
structures among flights that are largely absent at point-to-point airports.
The Turkish aviation market exemplifies this structure: Istanbul Airport
processed approximately 76.5 million passengers in 2023, with Turkish Airlines
accounting for over 51\% of total departures (DHMI, 2024). This carrier
concentration creates a fundamentally asymmetric operational environment in
which a single carrier's crew scheduling constraints, aircraft rotation
patterns, and hub-management decisions dominate aggregate delay dynamics in
ways that multi-carrier studies may not capture. Understanding this
concentration effect is therefore a prerequisite for building robust
single-airport delay prediction systems, and motivates the use of
airline-specific historical delay features as primary predictors in the present
study.





\section{Delay Propagation and Historical Delay Features}

One of the most influential mechanisms underlying system-level delay behavior is delay propagation --- the empirically documented tendency for a delay at one point in time to increase the probability of delays in subsequent periods. This temporal persistence arises through several distinct pathways: delayed aircraft occupying gates beyond their scheduled block time, crews missing positioning connections, passengers misconnecting between inbound and outbound flights, and accumulated congestion that prevents normal flow recovery. Understanding and encoding this propagation mechanism is widely recognized as essential for accurate predictive modeling.

Empirical research consistently confirms that delays exhibit short-term temporal autocorrelation. The proportion of delayed flights in a recent time window is one of the strongest predictors of current delay probability, reflecting how congestion propagates through a hub's operational cycle. Mamdouh et al. (2023) propose a flight delay prediction framework built around a novel algorithm that constructs features encoding historical delay behavior along specific flight paths. Their approach, which they term the Flight Delay Path Previous-based Machine Learning (FDPP-ML) method, demonstrates that incorporating historical delay context reduces prediction error substantially compared to models trained on static flight characteristics alone. Specifically, their framework achieved an average reduction in mean absolute error of 32\% over a 2-hour forecast horizon across a wide network of U.S. airports. Their findings reinforce the view that delay prediction models must account for short-term memory and temporal dependence rather than treating each flight as an independent observation.

A related concern in encoding historical delay features is the risk of temporal data leakage --- wherein features computed naively from the training dataset inadvertently incorporate information about the target variable. For example, computing a flight's ``previous hour delay rate'' by including the flight itself in that calculation creates a circular dependency that inflates training accuracy but collapses under deployment conditions. This methodological issue has received relatively limited attention in the published aviation machine learning literature, despite its practical importance for models intended for operational deployment.

Further evidence on delay propagation comes from studies examining how the
effect varies across time-of-day and operational context. Propagation is most
severe during peak departure banks --- the concentrated wave structures
characteristic of hub operations --- where gate conflicts and crew
misconnections simultaneously affect multiple flights within a narrow
scheduling window. Off-peak propagation is substantially lower due to greater
schedule slack and reduced resource contention. This temporal heterogeneity
implies that effective historical delay features must encode not only the
average delay rate in a preceding time window but also the operational context
in which those delays occurred. Rolling features aggregated over fixed time
horizons conflate high-contagion peak periods with low-contagion off-peak
periods, potentially diluting the predictive signal available from congestion
state indicators (Mamdouh et al., 2023).

Airline-level heterogeneity in propagation represents a related but
under-studied dimension of delay dynamics. Carriers differ substantially in
their schedule buffer policies, crew redundancy provisions, and recovery
prioritisation procedures, all of which modulate the rate at which an initial
delay is absorbed or transmitted to subsequent flights. Network carriers with
tighter rotation schedules tend to exhibit higher propagation coefficients than
regional or low-cost carriers operating simpler point-to-point networks.
Capturing this airline-specific heterogeneity requires carrier-level historical
delay aggregations as model features --- a design choice that distinguishes the
present study's feature engineering from simpler airport-wide aggregation
approaches. Empirical quantification of airline-level contagion heterogeneity at a single 
major hub airport, and its implications for carrier-specific feature engineering, 
remains largely absent from the published literature --- a gap the present study 
directly addresses.

\section{Machine Learning Approaches for Flight Delay Prediction}

\subsection{From Statistical Models to Classical Machine Learning}

Quantitative prediction of flight delays has evolved significantly over the past two decades. Early predictive models relied on regression-based techniques and rule-based systems, typically incorporating scheduled departure time, airline identifiers, airport characteristics, and basic meteorological conditions as explanatory variables. These models demonstrated that delays are partially predictable from observable pre-flight information, but their linear functional forms struggled to capture the complex nonlinear interactions among heterogeneous features that characterize aviation operational data.

The increasing availability of large-scale aviation datasets enabled the adoption of more flexible supervised machine learning techniques. Logistic regression, decision trees, and support vector machines have been widely applied to classify flights as delayed or on-time, yielding improvements over purely statistical baselines by allowing nonlinear decision boundaries and higher-dimensional feature representations. Kim et al. (2016) apply a multilevel artificial neural network architecture to predict flight delays at JFK International Airport, leveraging a hierarchical input structure designed to reflect the layered nature of delay causality. Their model outperforms conventional backpropagation networks in terms of prediction accuracy, demonstrating the value of domain-informed architectural choices even in relatively early neural network applications. Similarly, Choi et al. (2016) address weather-induced delay prediction across multiple machine learning algorithms, reporting that model performance is strongly dependent on the quality and temporal resolution of meteorological features, a finding that has since been replicated across numerous applied studies.


An important but frequently overlooked issue in early machine learning
applications to delay prediction concerns the implicit stationarity assumption
embedded in standard evaluation practices. When a stratified random train-test
split is applied to a dataset spanning multiple months, training and test
observations are drawn from the same temporal distribution, producing optimistic
performance estimates that do not reflect genuine operational forecasting
conditions. Regression and classification models evaluated under such conditions
may achieve high reported accuracy while effectively encoding seasonal patterns
from the test period that would be unavailable under genuine deployment. The
failure to distinguish between interpolation performance --- predicting within
a known distribution --- and extrapolation performance --- forecasting genuinely
future distributions --- has contributed to an optimism bias in the published
flight delay prediction literature. Giusti et al.\ (2022) explicitly document
this phenomenon, demonstrating that model performance degrades substantially as
the temporal distance between training and test periods increases, and that
standard adaptive retraining schedules are insufficient to maintain stable
performance across seasons and years. The present study directly quantifies
this optimism bias by evaluating identical models under both random-split and
temporal-split protocols, finding a gap of 0.429 AUC points attributable
entirely to evaluation methodology rather than model quality.


Despite their improvements over statistical baselines, classical machine learning approaches exhibited limited effectiveness when applied to highly congested or irregular operational conditions, where the interaction effects between temporal, operational, and environmental features become especially complex and where simple feature representations prove inadequate.

\subsection{Ensemble Methods and Gradient Boosting}

More recent applied research has consistently identified ensemble-based machine learning methods as the strongest performers in flight delay prediction tasks. Tree-based ensemble models --- particularly Random Forest and gradient boosting variants such as XGBoost and LightGBM --- are well suited to structured aviation data characterized by mixed variable types, nonlinear relationships, and complex feature interactions. These methods achieve robustness through averaging or sequential error correction across multiple weak learners, reducing variance without requiring extensive manual feature engineering.

Gui et al. (2019) apply ensemble learning methods to departure delay prediction at Shanghai Pudong International Airport (PVG) and report a ROC-AUC of 0.88, establishing what has become a widely cited performance benchmark for single-airport binary delay classification studies. Their results demonstrate that properly engineered ensemble models can achieve strong discriminative performance even with moderate-sized training datasets, provided that meteorological, operational, and temporal features are comprehensively integrated.

Alfarhood et al. (2024) conduct a comprehensive case study using multi-source operational and meteorological data from Saudi Arabian Airlines, evaluating a wide range of machine learning models including Random Forest, gradient boosting, and neural networks. Their findings confirm that tree-based ensemble and boosting methods consistently outperform linear approaches in terms of discrimination performance. Importantly, their study also adopts an operationally meaningful delay threshold and explicitly discusses the practical relevance of classification-based delay prediction frameworks for airline operations management.

Mamdouh et al. (2023) further confirm that ensemble approaches benefit significantly from incorporating historical delay context, particularly under congested operational conditions. Taken together, these studies establish ensemble-based gradient boosting --- and XGBoost in particular --- as the current methodological standard for applied flight delay classification, combining strong predictive performance with computational efficiency and relative interpretability compared to deep learning alternatives.


\subsection{Deep Learning Architectures}

Deep learning approaches have also been explored for flight delay prediction, particularly in contexts involving complex temporal or spatial dependencies. Recurrent neural networks, including Long Short-Term Memory (LSTM) architectures, have been applied to model sequences of flights and capture long-range temporal dependencies associated with delay propagation. Convolutional neural networks and graph neural networks have been employed to represent spatial relationships among airports and airspace sectors. Yu et al. (2019) investigate deep learning-based delay prediction for commercial air transport at Beijing Capital International Airport (PEK), reporting a ROC-AUC of 0.85 using a deep belief network architecture trained on multi-attribute high-dimensional flight data.

Despite their representational power, the practical advantages of deep learning models in operational delay prediction remain mixed. While deep learning approaches can outperform ensemble methods when very large volumes of high-frequency data are available, they require substantial computational resources, extensive hyperparameter tuning, and large labeled datasets. Moreover, their limited interpretability presents challenges for operational deployment, where transparency and explainability are often critical for gaining practitioner trust and regulatory acceptance. Several recent studies emphasize that model complexity should be aligned with data availability, interpretability requirements, and operational constraints rather than pursued solely for marginal performance gains. For moderate-sized datasets typical of single-airport operational studies --- such as the 9,527-flight dataset analyzed in the present research --- well-tuned ensemble methods have consistently matched or exceeded deep learning performance while providing substantially greater interpretability.

\section{Feature Engineering and Multi-Source Data Integration}

\subsection{Meteorological Features}

Weather conditions constitute a major contributor to flight delays and have been extensively studied in the aviation delay prediction literature. Meteorological factors including wind speed, precipitation, visibility, atmospheric pressure, and ambient temperature can directly affect runway operations, aircraft performance, air traffic control flow management, and safety margins. Numerous studies demonstrate that even relatively coarse weather information, when properly aligned with flight schedules, provides meaningful predictive signal. Adverse weather events tend to amplify delay contagion by simultaneously disrupting multiple flights across a hub, creating conditions in which individual weather measurements capture system-wide operational stress.

The integration of meteorological data into delay prediction models requires careful attention to temporal alignment. Weather observations must be matched to scheduled departure times at appropriate temporal resolutions to avoid information misalignment or inadvertent leakage. Some studies rely on coarse daily weather summaries, while others incorporate hourly or sub-hourly observations. These differences complicate direct comparison of reported performance metrics and raise questions about the robustness of results across varying data granularities. The present research employs hourly weather observations obtained at Istanbul Airport coordinates, using a conservative nearest-neighbor temporal matching strategy with a 4-hour tolerance window to balance coverage and temporal validity.

\subsection{Temporal and Operational Features}

Beyond direct meteorological effects, temporal patterns play a critical role in delay formation. Delay likelihood varies systematically across hours of the day, days of the week, and seasonal periods, reflecting recurring demand cycles and capacity constraints. Empirical evidence consistently shows that delays cluster during peak traffic windows, reinforcing the importance of circadian and weekly temporal features in predictive modeling. Incorporating short-term historical delay context has been shown to improve predictive stability and reduce sensitivity to transient noise. Operational features such as airline identity, aircraft type, and route-level characteristics further capture systematic differences in operational reliability that persist across time.

The role of temporal features extends beyond simple hour-of-day and 
day-of-week indicators to encompass multi-scale cyclical patterns that 
reflect the layered scheduling structure of hub operations. At the daily 
level, departure banks --- the concentrated windows of simultaneous 
departures characteristic of hub carrier operations --- create predictable 
congestion peaks in which delay probability rises sharply before recovering 
during inter-bank valleys. At the weekly level, reduced frequencies on 
weekends translate to lower aggregate delay rates, while holiday periods 
introduce demand surges that stress capacity constraints in ways that 
standard temporal features cannot capture. At the seasonal level, demand 
patterns, weather regimes, and operational protocols introduce cyclical 
variation in both delay prevalence and propagation intensity. Capturing 
these multi-scale temporal effects requires feature design that reflects 
the hierarchical structure of aviation scheduling rather than treating 
temporal variables as generic numerical inputs.

Operational identity features --- primarily airline carrier indicators --- 
encode systematic differences in punctuality culture, network topology, 
fleet composition, and hub priority that persist independently of weather 
and scheduling conditions. At hub airports where a dominant carrier holds 
priority access to preferred gates and crew resources, carrier identity 
also proxies for differential access to recovery mechanisms unavailable to 
smaller operators. One-hot encoded airline dummy variables, despite their 
high dimensionality, have been shown to carry meaningful discriminative 
signal in single-airport studies (Alfarhood et al., 2024), reflecting real 
operational heterogeneity rather than mere label memorisation. The practical 
implication is that prediction models incorporating carrier identity can 
distinguish between structurally punctual and structurally delayed carriers, 
improving calibration across the full distribution of departure events.


\subsection{Data Integration Challenges and Imputation}

A further challenge in multi-source delay prediction research concerns the integration of heterogeneous data sources that differ in temporal resolution, spatial coverage, and completeness. Many published studies rely on proprietary or curated datasets that have been pre-cleaned before analysis, limiting insight into how data quality decisions affect model performance. In practice, operational data frequently contains missing values, timestamp inconsistencies, and encoding errors that require systematic preprocessing. Weather data, in particular, often covers only a fraction of observed flight hours due to API limitations, sensor gaps, or data retention policies.

Several strategies for handling meteorological data missingness have been proposed, ranging from complete-case analysis --- which discards observations lacking exact weather matches --- to various imputation approaches. The present study implements temporal nearest-neighbor imputation with an explicitly validated tolerance window, supported by autocorrelation analysis demonstrating that adjacent-hour weather observations remain highly correlated ($\rho > 0.73$ at 4-hour lag for all variables). This approach recovers substantially more usable flight records than complete-case analysis while maintaining defensible assumptions about imputation validity.

\section{Evaluation Methodology and Temporal Validity}

\subsection{Limitations of Random Train-Test Splits}

Evaluation methodology represents a critically important but frequently underexamined dimension of flight delay prediction research. The majority of published studies evaluate model performance using stratified random train-test splits, in which observations are randomly assigned to training and test partitions without regard to their chronological ordering. This approach reflects standard machine learning evaluation practice but is fundamentally misaligned with operational deployment realities, where models trained on historical data are required to predict outcomes in genuinely future time periods.

Random splits artificially intermix temporal periods, allowing models to train on data from the same calendar month or week as their test observations. In the presence of temporal autocorrelation --- which is pervasive in aviation operational data --- this produces inflated performance estimates that may substantially overstate real-world deployment effectiveness. Studies employing random splits report performance metrics as though they reflect deployment-realistic conditions, when in fact they reflect a favorable evaluation setting that does not simulate genuine forecasting under distribution shift.

\subsection{Concept Drift and Non-Stationarity}

A closely related challenge concerns temporal non-stationarity in delay patterns --- the tendency for the statistical relationships between features and delay outcomes to change over time as a function of seasonal operational dynamics, policy changes, demand fluctuations, and environmental variation. Concept drift, wherein feature-target relationships learned during training no longer hold during deployment, represents a fundamental threat to the real-world effectiveness of static trained models.

Giusti et al. (2022) specifically investigate flight delay prediction under concept drift and demonstrate that model performance degrades substantially as the temporal distance between training and test periods increases. Their study evaluates multiple drift handling strategies and finds that neither static model retraining schedules nor standard adaptive approaches are fully sufficient to maintain stable performance across seasons and years. This finding has significant implications for the operational deployment of aviation delay prediction systems and motivates the investigation of adaptive training strategies, online learning approaches, and season-specific ensemble methods.

The present research provides direct empirical evidence of concept drift through the discovery of a seasonal regime shift in delay propagation mechanisms --- wherein delay contagion effects reverse completely from autumn (+7.07 percentage points) to winter ($-7.78$ percentage points) operational periods. This finding quantifies the severity of non-stationarity at Istanbul Airport and contributes empirical evidence to the growing literature calling for mandatory temporal validation in aviation machine learning research.

\subsection{The Case for Chronological Evaluation}

Recognition of these limitations has prompted recent calls for chronologically consistent evaluation protocols in aviation machine learning research, wherein models are trained exclusively on historical data and evaluated on genuinely future periods. Such temporal train-test splits expose distribution shift and pattern instability that are masked by random evaluation designs, providing more honest assessments of deployment readiness. The growing adoption of this evaluation philosophy reflects a broader shift toward operationally realistic benchmarking standards in applied aviation analytics.

The practical implementation of chronological evaluation introduces several 
methodological decisions that must be handled carefully. The choice of the 
training-test temporal boundary affects the measured performance gap, as 
boundaries placed at seasonal transitions may encounter larger distribution 
shifts than boundaries within a stable seasonal regime. Hyperparameter tuning 
is equally consequential: using standard $k$-fold cross-validation on the 
training set risks selecting configurations that benefit from within-period 
temporal autocorrelation, requiring instead forward-chaining time-series 
cross-validation that preserves temporal ordering within folds. The treatment 
of transitional periods requires deliberate analytical design: including 
transition months in the training set could partially expose the model to the 
distribution shift it is intended to detect, while excluding them creates an 
evaluation window that under-samples the performance trajectory during regime 
change.

Despite growing recognition of these issues, chronological evaluation remains 
the exception rather than the rule in published aviation delay prediction 
research. The majority of studies reviewed in the present chapter employ 
stratified random splits without discussing the temporal implications of this 
choice. Where temporal validation is attempted, it is often limited to a single 
holdout month directly adjacent to the training period, which fails to expose 
performance under multi-month distributional shift. Comprehensive chronological 
evaluation requiring a separation of at least two calendar months between the 
end of training and the beginning of testing --- as implemented in the present 
study --- is to the authors' knowledge absent from the single-airport flight 
delay prediction literature prior to this work.


\section{Model Interpretability and Operational Deployment}

Interpretability has emerged as an increasingly critical consideration in aviation machine learning research. While early studies primarily emphasized predictive accuracy, recent work highlights the need to understand model behavior and feature contributions, particularly in safety-critical and operationally sensitive environments such as commercial aviation, where prediction outputs must be explainable to operations staff, airline management, and regulatory bodies.

Tree-based ensemble models offer inherent interpretability advantages through built-in feature importance measures based on impurity reduction or permutation testing. However, these measures have known limitations, including sensitivity to feature correlation and inability to capture directional effects or interaction-specific contributions. Post-hoc explanation techniques have been developed to address these limitations, with SHAP (SHapley Additive exPlanations) emerging as the leading framework for model-agnostic feature attribution in applied machine learning.

SHAP assigns to each feature a contribution value for each individual prediction, grounded in cooperative game theory's Shapley value framework. This approach has been applied across a range of aviation prediction tasks, providing granular insight into which features drive individual delay predictions and how feature importance varies across operational contexts. Recent studies have demonstrated that SHAP-based analyses reveal operationally meaningful patterns --- such as the disproportionate contribution of atmospheric pressure and wind speed to delay probability --- that align with domain knowledge about meteorological effects on runway operations and air traffic control procedures. SHAP analysis of the best-performing XGBoost model yields a feature importance hierarchy in which historical delay propagation features account for 40.6 percent of total predictive contribution, meteorological features account for 33.9\% percent, airline identity contributes 10.2, surface traffic features account for only 8.9\%, and temporal features contribute 6.4 percent, providing actionable insights for operational decision support system design.

Beyond SHAP, several complementary interpretability approaches have been 
proposed for aviation delay prediction. Partial dependence plots (PDPs) 
visualise the marginal effect of individual features on predicted delay 
probability while averaging over the distribution of other features, 
revealing non-linear response curves that are invisible to linear coefficient 
analysis. For example, the relationship between wind speed and delay probability 
typically exhibits a threshold structure: delays remain stable at low-to-moderate 
wind speeds before rising sharply above operational limits for safe approach 
and departure. Similarly, temperature shows a non-linear relationship driven 
by de-icing requirements at low temperatures and convective activity at high 
temperatures. These threshold effects are well-suited to tree-based models, 
whose recursive partitioning naturally learns breakpoints in feature space 
without requiring manual engineering of interaction terms or polynomial 
transformations.

The operational deployment of interpretable models also carries implications 
for regulatory compliance and stakeholder trust. Aviation safety regulators 
and airline operations managers require that automated decision-support systems 
produce auditable, explainable outputs that can be reviewed, challenged, and 
overridden by qualified personnel. Black-box deep learning models, regardless 
of their predictive accuracy, face significant adoption barriers in 
safety-critical environments where unexplained predictions cannot be 
accepted as the basis for resource allocation decisions. The combination of 
strong predictive performance and SHAP-based interpretability offered by 
gradient boosting methods such as XGBoost therefore represents a practical 
optimum for operational aviation deployment --- balancing accuracy, 
computational efficiency, and the transparency requirements of the regulatory 
environment (Alfarhood et al., 2024).


\section{Research Gaps and Positioning of the Present Study}

Despite substantial progress in the flight delay prediction literature, several important limitations remain. First, the majority of published studies employ random train-test splits that overestimate real-world deployment performance by masking temporal non-stationarity. Performance benchmarks reported in the literature --- including the widely cited 0.88 ROC-AUC of Gui et al. (2019) and the 0.85 ROC-AUC of Yu et al. (2019)--- are almost universally derived under random split evaluation conditions that do not simulate genuine operational forecasting. The absence of systematic temporal validation across multi-season datasets means that the field lacks reliable evidence on how well current models generalize to genuinely future periods.

Second, delay propagation analysis in existing studies has been largely aggregate and cross-sectional, without examining how propagation mechanisms vary across seasonal regimes. The possibility of seasonal non-stationarity in delay contagion --- with qualitatively different propagation dynamics in different periods of the year --- has not been empirically documented or quantitatively characterized.

Third, despite theoretical arguments for the relevance of landside accessibility to airport punctuality, no published study has empirically tested the predictive contribution of surface road traffic data to flight delay models. This hypothesis, which motivated the collection of surface traffic indicators in the present study, remained empirically unexamined in the aviation delay prediction literature prior to this work.

Fourth, many studies either employ highly simplified feature sets that fail to capture system-level interactions or rely on proprietary, non-reproducible datasets. The growing availability of open-access data sources --- including public weather APIs and commercial flight operational APIs --- presents an opportunity for more transparent and reproducible research that has been underexplored.

The present study addresses these gaps through four specific contributions. First, it implements a strict chronological train-test evaluation framework that exposes the full extent of temporal performance degradation from random-split (0.903 ROC-AUC) to temporal-split (0.474 ROC-AUC) conditions, providing direct empirical evidence of non-stationarity masked by standard evaluation practices. Second, it provides the first quantitative documentation of seasonal regime shift in delay propagation mechanisms, characterizing a 14.84 percentage point swing in delay contagion between autumn and winter operational periods. Third, it conducts the first empirical test of landside traffic-delay coupling in aviation delay prediction, conclusively demonstrating negligible predictive contribution and redirecting future research toward airside operational constraints. Fourth, it implements and validates a conservative temporal nearest-neighbor weather imputation strategy that substantially increases dataset size while maintaining defensible data quality standards, providing a reproducible methodology for future studies facing similar meteorological data limitations.











\chapter{Methodology}
\label{ch:methodology}

\section{Research Design}
\label{sec:research_design}

This study employs a quantitative research approach using supervised machine learning to predict binary flight delay outcomes at Istanbul Airport. The research design encompasses four primary phases: data collection and integration, feature engineering with temporal leakage prevention, model development and training, and comprehensive evaluation through both random-split and temporal validation strategies.

The methodological framework prioritizes temporal validity and operational realism. Unlike studies employing random train-test splits that artificially intermix temporal periods, this research implements strict chronological separation to assess true predictive performance under deployment conditions. This design choice enables detection of temporal instability in delay patterns—a critical consideration for operational aviation systems where models must generalize to future unseen periods.

The research adopts an iterative refinement approach, developing the dataset and models in two distinct phases. Phase 1 (Dataset V1) establishes baseline methodology with conservative direct weather matching, while Phase 2 (Dataset V1 Imputed Fixed) implements weather imputation and temporal expansion to produce a substantially larger analytical dataset. This iterative process demonstrates scientific rigor through systematic problem identification and resolution.

\section{Dataset Description}
\label{sec:dataset}

\subsection{Phase 1 Dataset (Dataset V1)}
\label{sec:phase1_dataset}

The initial dataset (Dataset V1) comprised 4,330 flights operating at Istanbul Airport (IATA: IST, ICAO: LTFM) during October and November 2025. This dataset integrated three primary data sources: flight operational records, meteorological observations, and surface traffic measurements, using conservative direct temporal matching without imputation.

\subsubsection{Data Sources}

\paragraph{Flight Operational Data}

Flight records were collected via the AviationStack API (https://aviationstack.com), which aggregates real-time and historical flight information from multiple aviation data providers. The API provides structured JSON responses containing:

\begin{itemize}
    \item Flight identification (IATA code, airline, route)
    \item Scheduled departure timestamp (UTC)
    \item Actual departure timestamp (UTC) when available
    \item Delay duration (minutes) computed from timestamp difference
\end{itemize}

\textbf{Data Collection Period:} October 27, 2025 -- November 30, 2025 (35 days)

\textbf{Raw Records Retrieved:} 42,623 flight records

\textbf{Records with Actual Departure Times:} 9,814 (23.0\%)

The 77\% missingness rate primarily reflects future scheduled flights lacking actual departure data, cancelled operations, and API data gaps. Missingness analysis revealed no systematic correlation with observable flight characteristics (hour, airline, destination), supporting a Missing Completely At Random (MCAR) assumption.

\paragraph{Meteorological Data}

Weather observations were obtained from the OpenWeather API (https://openweathermap.org/api) using Istanbul Airport coordinates (41.2753°N, 28.7519°E). Variables collected include:

\begin{itemize}
    \item Temperature (°C)
    \item Relative humidity (\%)
    \item Atmospheric pressure (hPa)
    \item Wind speed (m/s)
    \item Weather category (Clear, Clouds, Rain, Snow)
\end{itemize}

\textbf{Temporal Resolution:} Hourly observations

\textbf{Phase 1 Coverage:} 377 hourly records (October: 44h, November: 201h, December: 8h, January: 8h)

\textbf{Temporal Density:} 22\% coverage during study period

Phase 1 methodology employed strict direct matching: flights were included only if weather observations existed for their exact scheduled departure hour. This conservative approach ensured data quality but limited dataset size substantially.

\paragraph{Surface Traffic Data}

Road traffic indicators were sourced from Istanbul Metropolitan Municipality (İBB) Traffic Management Center, covering major airport access routes (E80/TEM highway, O-3/O-7 ring roads). Due to data access constraints, a representative 24-hour traffic profile was constructed by averaging conditions across each hour-of-day over a 6-month historical window. Three composite indicators were derived:

\begin{itemize}
    \item \texttt{road\_avg\_speed} (km/h): Mean vehicle speed on monitored routes
    \item \texttt{road\_vehicle\_count}: Total vehicles per hour
    \item \texttt{road\_congestion\_index}: Composite metric combining speed reduction and volume relative to free-flow capacity
\end{itemize}

This approach captures typical circadian traffic patterns (morning/evening rush hours, off-peak periods) but cannot reflect day-to-day variability from weather, accidents, or special events.

\subsubsection{Dataset V1 Characteristics}

\begin{table}[h]
\centering
\caption{Dataset V1 summary statistics}
\label{tab:phase1_summary}
\begin{tabular}{lr}
\toprule
\textbf{Characteristic} & \textbf{Value} \\
\midrule
Total flights & 4,330 \\
Temporal span & Oct 27 -- Nov 30, 2025 (35 days) \\
Weather observations (hourly) & 377 hours \\
Weather coverage rate & 45.4\% \\
\midrule
\textbf{Monthly Distribution:} & \\
October 2025 & 1,435 flights (51.6\% delayed) \\
November 2025 & 2,887 flights (58.3\% delayed) \\
December 2025 & 8 flights \\
\midrule
Overall delay rate ($\geq$15 min) & 56.2\% \\
Features per observation & 20 base + airline dummies \\
Total dimensionality & $\sim$461 features \\
\bottomrule
\end{tabular}
\end{table}

The 45.4\% weather coverage rate reflects the conservative constraint that flights could only be included if weather observations existed for their exact scheduled departure hour. This strategy ensured data quality but substantially limited dataset size.

\subsection{Phase 2 Dataset (Dataset V1 Imputed Fixed)}
\label{sec:phase2_dataset}

Following identification of severe data coverage limitations in Phase 1, the dataset was reconstructed with temporal weather imputation and extended temporal coverage, yielding Dataset V1 Imputed Fixed comprising 9,527 flights spanning five months.

\subsubsection{Methodological Improvements}

Phase 2 dataset construction implemented four key improvements over Phase 1:

\begin{enumerate}
    \item \textbf{Weather Imputation}: Rather than excluding flights lacking exact weather matches, missing meteorological data were imputed using temporal nearest-neighbor matching within conservative 4-hour windows (detailed in Section~\ref{sec:weather_imputation}). This increased weather coverage from 45.4\% to 68.2\%.
    
    \item \textbf{Extended Temporal Coverage}: Inclusion of December 2025 data (565 flights), providing critical winter operational context and enabling assessment of seasonal delay pattern variation.
    
    \item \textbf{Bias Correction}: Removal of inadvertent filtering steps that preferentially excluded certain flight categories, improving representativeness.
    
    \item \textbf{Conservative Tolerance}: Weather imputation employed 4-hour tolerance window (rather than 6-hour initially explored), balancing coverage gains against temporal validity concerns.
\end{enumerate}

\subsubsection{Dataset V1 Imputed Fixed Characteristics}

\begin{table}[h]
\centering
\caption{Dataset V1 Imputed Fixed summary statistics}
\label{tab:phase2_summary}
\begin{tabular}{lr}
\toprule
\textbf{Characteristic} & \textbf{Value} \\
\midrule
Total flights & 9,527 \\
Temporal span & Oct 2025 -- Feb 2026 (5 months) \\
Weather coverage (post-imputation) & 68.2\% \\
Weather imputation tolerance & 4 hours \\
\midrule
\textbf{Monthly Distribution:} & \\
October 2025 & 5,329 flights (79.2\% delayed) \\
November 2025 & 3,200 flights (79.9\% delayed) \\
December 2025 & 565 flights (78.4\% delayed) \\
January 2026 & 241 flights (76.8\% delayed) \\
February 2026 & 192 flights (77.6\% delayed) \\
\midrule
Overall delay rate ($\geq$15 min) & 79.4\% \\
Delayed flights & 7,563 \\
On-time flights & 1,964 \\
\bottomrule
\end{tabular}
\end{table}

\textbf{Dataset Expansion:} Phase 2 dataset is 2.2× larger than Phase 1 (9,527 vs. 4,330 flights), providing substantially improved statistical power and enabling more robust model training.

\textbf{Train-Test Split:} Chronological split with October-November 2025 as training data (8,529 flights, 89.5\%) and January-February 2026 as test data (433 flights, 4.5\%). December 2025 (565 flights, 6.0\%) reserved for intermediate temporal validation.

\textbf{Delay Rate Increase:} Phase 2 exhibits higher delay rate (79.4\% vs. 56.2\% in Phase 1). This reflects three factors: (1) improved data completeness capturing more operational reality, (2) inclusion of winter months (December-February) with elevated delay rates, and (3) correction of selection biases that preferentially excluded delayed flights in Phase 1.

\subsection{Variables and Target Definition}
\label{sec:variables_target}

\subsubsection{Target Variable}

The dependent variable is binary delay status, defined as:

\begin{equation}
y = \begin{cases} 
1 & \text{if } \delta \geq 15 \text{ minutes (delayed)} \\
0 & \text{if } \delta < 15 \text{ minutes (on-time)}
\end{cases}
\end{equation}

where $\delta = t_{\text{actual}} - t_{\text{scheduled}}$ represents departure delay in minutes.

The 15-minute threshold aligns with international aviation standards established by the U.S. Department of Transportation (14 CFR Part 234), EUROCONTROL's Central Office for Delay Analysis (CODA), and IATA operational reporting guidelines. This threshold distinguishes minor schedule deviations absorbed by operational buffers from significant delays requiring resource interventions (gate reassignments, crew adjustments, passenger connection management).

\subsubsection{Predictor Variables}

Predictor variables span four conceptual categories:

\paragraph{Temporal Features}
\begin{itemize}
    \item \texttt{hour\_of\_day}: Scheduled departure hour (0-23), capturing circadian operational patterns
    \item \texttt{day\_of\_week}: Day identifier (0-6), differentiating weekday vs. weekend dynamics
\end{itemize}

\paragraph{Meteorological Features}
\begin{itemize}
    \item \texttt{temperature}: Ambient air temperature (°C)
    \item \texttt{humidity}: Relative humidity (\%)
    \item \texttt{pressure}: Atmospheric pressure (hPa)
    \item \texttt{wind\_speed}: Sustained wind speed (m/s)
    \item \texttt{weather}: Categorical condition (Clear, Clouds, Rain, Snow)
\end{itemize}

\paragraph{Historical Delay Features}
\begin{itemize}
    \item \texttt{prev\_hour\_delay\_rate}: Proportion of delayed flights in previous hour
    \item \texttt{avg\_delay\_last\_24h}: Rolling 24-hour average delay rate
    \item \texttt{airline\_7d\_delay\_rate}: Carrier-specific 7-day rolling delay rate
\end{itemize}

\paragraph{Surface Traffic Features}
\begin{itemize}
    \item \texttt{road\_avg\_speed}: Average vehicle speed on airport access routes (km/h)
    \item \texttt{road\_vehicle\_count}: Hourly vehicle volume
    \item \texttt{road\_congestion\_index}: Composite congestion severity metric
\end{itemize}

\paragraph{Airline Identity}
\begin{itemize}
    \item \texttt{airline\_name}: Categorical carrier identifier (one-hot encoded)
\end{itemize}

\section{Data Preprocessing and Feature Engineering}
\label{sec:preprocessing}

\subsection{Data Cleaning}
\label{sec:data_cleaning}

\subsubsection{Missing Data Handling}

Flight records lacking actual departure timestamps were excluded, as the target variable cannot be computed without both scheduled and actual departure times. This resulted in 77\% data reduction from raw API records (42,623 → 9,814 flights with computable delays).

Meteorological data exhibited 78\% missingness during the study period. Phase 1 excluded flights with missing weather (reducing to 4,330 flights), while Phase 2 employed temporal imputation (Section~\ref{sec:weather_imputation}) to recover coverage.

Airline names contained occasional parsing errors manifesting as timestamp values. These erroneous entries (n=47, 0.5\%) were manually corrected through cross-referencing with flight IATA codes.

\subsubsection{Outlier Removal}

\paragraph{Delay Outliers}

Flights exhibiting extreme delays ($|\delta| > 720$ minutes, i.e., $>$12 hours) were excluded (n=82, 0.8\% of valid records). Such cases likely represent:

\begin{itemize}
    \item Data entry errors (timestamp typos)
    \item Flight diversions (arrival airport changed, invalidating original schedule)
    \item Multi-day delays from extraordinary circumstances
\end{itemize}

These events constitute fundamentally different operational regimes unsuitable for standard delay prediction models.

\paragraph{Weather Variable Outliers}

Applied three-sigma rule combined with physical plausibility constraints:

\begin{equation}
\text{Outlier if: } |x - \mu| > 3\sigma \quad \text{AND} \quad x \notin [\text{min}_{\text{physical}}, \text{max}_{\text{physical}}]
\end{equation}

\begin{table}[h]
\centering
\caption{Weather variable outlier detection}
\label{tab:weather_outliers}
\begin{tabular}{lrrr}
\toprule
\textbf{Variable} & \textbf{Physical Range} & \textbf{Outliers} & \textbf{Action} \\
\midrule
Temperature & [-20°C, 50°C] & 3 & Removed \\
Humidity & [0\%, 100\%] & 0 & None \\
Pressure & [980 hPa, 1050 hPa] & 2 & Removed \\
Wind Speed & [0 m/s, 25 m/s] & 8 & Capped at 25 m/s \\
\bottomrule
\end{tabular}
\end{table}

Total weather outliers removed: 13 observations (1.6\% of hourly records).

\subsubsection{Timestamp Standardization}

All temporal data were standardized to Coordinated Universal Time (UTC) using \texttt{pandas.to\_datetime(utc=True)}. Approximately 3\% of records provided timestamps in local time (UTC+3, Turkey Time), identified through cross-validation with known daylight patterns and corrected algorithmically.

\subsection{Feature Selection}
\label{sec:feature_selection}

Feature selection employed domain-informed inclusion based on aviation operational theory and prior empirical literature, rather than automated statistical selection procedures. This approach ensures features encode aviation-relevant causal mechanisms.

\textbf{Included Features:} All temporal, meteorological, historical, traffic, and airline features listed in Section~\ref{sec:variables_target} were retained, as each encodes plausible delay-causal pathways documented in aviation operations research.

\textbf{Excluded Features:} Destination airport codes were excluded due to extreme cardinality (500+ unique destinations) creating sparse one-hot representations. Route-level analysis was deferred to future work.

\subsection{Temporal Leakage Prevention}
\label{sec:leakage}

Temporal data leakage—wherein predictive features inadvertently contain information from or about the target variable—poses a critical threat to model validity in time-series forecasting. This study implements rigorous leakage prevention protocols for historical delay features.

\subsubsection{Leakage Risk in Historical Features}

Historical delay aggregations (\texttt{prev\_hour\_delay\_rate}, \texttt{avg\_delay\_last\_24h}, \texttt{airline\_7d\_delay\_rate}) summarize past delay patterns. Naive computation including the current flight introduces circular dependency:

\begin{equation}
\texttt{prev\_hour\_delay\_rate}_{\text{naive}} = \frac{\sum_{i \in \text{FlightsInHour}(t)} \mathbb{1}(y_i = 1)}{|\text{FlightsInHour}(t)|}
\end{equation}

This formulation allows the model to "see" whether the current flight is delayed, creating spuriously high predictive performance that collapses under deployment.

\subsubsection{Strict Temporal Ordering Protocol}

To prevent leakage, historical features for flight $i$ at time $t$ are computed using \textit{exclusively} flights with scheduled departure $< t$:

\begin{equation}
\texttt{prev\_hour\_delay\_rate}_{\text{correct}}(t) = \frac{\sum_{j \in \text{Flights}(t-1)} \mathbb{1}(y_j = 1)}{|\text{Flights}(t-1)|}
\end{equation}

where $\text{Flights}(t-1)$ denotes flights departing in the hour strictly \textit{before} $t$.

Similarly, rolling averages employ only temporally preceding data:

\begin{equation}
\texttt{avg\_delay\_last\_24h}(t) = \frac{\sum_{j \in \text{Flights}([t-24, t))} \mathbb{1}(y_j = 1)}{|\text{Flights}([t-24, t))|}
\end{equation}

\subsubsection{Implementation via Chronological Iteration}

Enforcement required chronological iteration through the dataset rather than vectorized computation:

\begin{lstlisting}[language=Python, caption={Leakage-free historical feature computation}]
df_sorted = df.sort_values('scheduled_departure_time')

for idx, current_flight in df_sorted.iterrows():
    current_time = current_flight['scheduled_departure_time']
    
    # Previous hour flights: strictly before current time
    prev_hour_flights = df_sorted[
        (df_sorted['scheduled_departure_time'] >= current_time - pd.Timedelta(hours=1)) &
        (df_sorted['scheduled_departure_time'] < current_time)
    ]
    
    if len(prev_hour_flights) > 0:
        df_sorted.at[idx, 'prev_hour_delay_rate'] = prev_hour_flights['delayed_15'].mean()
\end{lstlisting}

\subsubsection{Leakage Detection and Validation}

\paragraph{Statistical Diagnostic}

Leakage manifests as artificially smooth historical features exhibiting means and variances suspiciously matching the overall delay rate. Post-correction validation confirmed proper implementation:

Pre-correction: Feature mean exactly matches delay rate, indicating leakage.

Post-correction: Feature mean differs substantially from overall delay rate, and variance increases, reflecting genuine temporal variation rather than target contamination.

\subsection{Feature Construction}
\label{sec:feature_construction}

\subsubsection{Weather Data Imputation}
\label{sec:weather_imputation}

Meteorological observations covered only 22\% of flight hours, necessitating imputation to avoid massive data loss. The imputation strategy balances coverage and quality through temporal nearest-neighbor matching with conservative tolerance constraints.

\paragraph{Methodology}

For each flight at hour $h$ lacking exact weather observation:

\begin{equation}
h_{\text{weather}} = \arg\min_{h' \in \mathcal{H}_{\text{available}}} |h' - h|
\end{equation}

subject to proximity constraint:

\begin{equation}
|h_{\text{weather}} - h| \leq 4 \text{ hours}
\end{equation}

If no observation exists within the 4-hour window, the flight is excluded from the dataset.

\paragraph{Tolerance Selection Rationale}

Initial exploration tested 6-hour tolerance. Final methodology deliberately employs 4-hour tolerance, prioritizing data quality over sample size. Meteorological autocorrelation analysis justifies this choice:

\begin{table}[h]
\centering
\caption{Weather variable temporal autocorrelation}
\label{tab:weather_autocorr}
\begin{tabular}{lrrr}
\toprule
\textbf{Variable} & \textbf{1-hour lag} & \textbf{4-hour lag} & \textbf{6-hour lag} \\
\midrule
Temperature & 0.98 & 0.91 & 0.87 \\
Humidity & 0.94 & 0.84 & 0.79 \\
Pressure & 0.99 & 0.95 & 0.92 \\
Wind Speed & 0.86 & 0.73 & 0.64 \\
Weather Category & 0.91 & 0.82 & 0.75 \\
\bottomrule
\end{tabular}
\end{table}

High autocorrelation at 4-hour lag ($\rho > 0.73$ for all variables) validates the assumption that nearby observations provide reasonable proxies.

\paragraph{Imputation Quality Assessment}

Leave-one-out cross-validation on exact-match subset quantified imputation accuracy:

\begin{table}[h]
\centering
\caption{Weather imputation error analysis}
\label{tab:imputation_error}
\begin{tabular}{lrrr}
\toprule
\textbf{Variable} & \textbf{MAE} & \textbf{Range} & \textbf{Relative Error} \\
\midrule
Temperature & 1.4°C & 37°C & 3.8\% \\
Humidity & 5.1\% & 74\% & 6.9\% \\
Pressure & 0.9 hPa & 31 hPa & 2.9\% \\
Wind Speed & 1.3 m/s & 18 m/s & 7.2\% \\
Weather Category & 11\% mismatch & N/A & N/A \\
\bottomrule
\end{tabular}
\end{table}

Relative errors remain below 8\% for continuous variables, supporting imputation validity.

\subsubsection{Categorical Encoding}

\paragraph{One-Hot Encoding}

Categorical variables (airline name, weather condition) were transformed via one-hot encoding, creating binary indicator features for each category. Airlines with fewer than 30 flights were aggregated into an "Other" category to prevent overfitting on rare carriers.

\paragraph{Final Feature Dimensionality}

After encoding:
\begin{itemize}
    \item Numeric features: 12
    \item Weather categories: 4 (Clear, Clouds, Rain, Snow)
    \item Airline categories: 80-120 (varies by dataset split)
    \item Total features: $p \approx 96$ to 136
\end{itemize}

\section{Modeling Approach}
\label{sec:models}

The three algorithms selected --- Logistic Regression, Random Forest, and XGBoost span the interpretability-complexity spectrum and collectively satisfy the 
requirements of this study. Logistic Regression provides a linear, fully 
interpretable baseline whose coefficients allow direct verification of feature-delay 
relationships against domain expectations. Random Forest and XGBoost were preferred 
over other gradient boosting frameworks such as LightGBM because their wider 
adoption in the aviation delay prediction literature enables direct performance 
benchmarking against published studies. Support vector machines were excluded due to 
their poor scaling behaviour with high-dimensional one-hot encoded features and their 
lack of native probability calibration. Neural network architectures were excluded 
for three reasons: the dataset size (9,527 observations) is insufficient to exploit 
deep learning's capacity advantages; their limited interpretability conflicts with 
the study's emphasis on SHAP-based feature attribution for operational deployment; 
and prior work on comparable single-airport datasets has consistently found that 
well-tuned ensemble methods match or exceed deep learning performance when training 
data volume is limited (see Section~2.3.3). This selection therefore prioritises 
operational deployability and interpretability alongside predictive performance.

\subsection{Logistic Regression}
\label{sec:logistic_regression}

Binary logistic regression models the log-odds of delay as a linear combination of features:

\begin{equation}
\log\left(\frac{P(y=1|X)}{1-P(y=1|X)}\right) = \beta_0 + \sum_{j=1}^{p} \beta_j x_j
\end{equation}

\paragraph{Rationale}

Provides interpretable baseline with transparent feature-delay relationships via coefficients $\beta_j$. Linear decision boundaries enable rapid inference and straightforward deployment.

\paragraph{Implementation}

\texttt{sklearn.linear\_model.LogisticRegression} with L2 regularization (Ridge penalty). Regularization strength $C$ selected via 5-fold time-series cross-validation.

\paragraph{Class Imbalance Handling}

\texttt{class\_weight='balanced'} parameter automatically adjusts class weights inversely proportional to class frequencies:

\begin{equation}
w_{\text{class}} = \frac{n_{\text{total}}}{n_{\text{classes}} \times n_{\text{class}}}
\end{equation}

\subsection{Random Forest Models}
\label{sec:random_forest}

Random Forest constructs an ensemble of decorrelated decision trees via bootstrap aggregation and random feature subsampling:

\begin{equation}
\hat{y} = \frac{1}{B} \sum_{b=1}^{B} T_b(X)
\end{equation}

where $B$ is the number of trees and $T_b$ is the $b$-th bootstrapped tree.

\paragraph{Rationale}

Captures non-linear feature interactions and threshold effects without manual feature engineering. Robust to overfitting through ensemble averaging and feature randomization.

\paragraph{Implementation}

\texttt{sklearn.ensemble.RandomForestClassifier}

\paragraph{Hyperparameter Tuning}

Randomized search cross-validation (30 trials, 5-fold time-series CV):

\begin{table}[h]
\centering
\caption{Random Forest hyperparameter search space}
\label{tab:rf_hyperparams}
\begin{tabular}{ll}
\toprule
\textbf{Parameter} & \textbf{Search Space} \\
\midrule
\texttt{n\_estimators} & [100, 200, 300, 500] \\
\texttt{max\_depth} & [None, 10, 20, 30, 50] \\
\texttt{min\_samples\_split} & [2, 5, 10] \\
\texttt{min\_samples\_leaf} & [1, 2, 4] \\
\texttt{max\_features} & ['sqrt', 'log2', None] \\
\bottomrule
\end{tabular}
\end{table}

\subsection{XGBoost Models}
\label{sec:xgboost}

XGBoost (Extreme Gradient Boosting) implements gradient boosting with regularization:

\begin{equation}
F_m(X) = F_{m-1}(X) + \nu \cdot h_m(X)
\end{equation}

where $h_m$ is a weak learner (decision tree) trained on residuals of $F_{m-1}$, and $\nu$ is the learning rate.

\paragraph{Rationale}

State-of-the-art performance in tabular data competitions, combining gradient boosting's sequential error correction with L1/L2 regularization.

\paragraph{Implementation}

\texttt{xgboost.XGBClassifier}

Two XGBoost variants were developed:

\paragraph{XGBoost Untuned}

Default hyperparameters providing baseline gradient boosting performance without extensive optimization.

\paragraph{XGBoost Tuned}

Randomized search cross-validation (30 trials, 5-fold time-series CV):

\begin{table}[h]
\centering
\caption{XGBoost hyperparameter search space}
\label{tab:xgb_hyperparams}
\begin{tabular}{ll}
\toprule
\textbf{Parameter} & \textbf{Search Space} \\
\midrule
\texttt{n\_estimators} & [100, 200, 300] \\
\texttt{max\_depth} & [3, 5, 7, 9] \\
\texttt{learning\_rate} & [0.01, 0.05, 0.1, 0.2] \\
\texttt{subsample} & [0.6, 0.8, 1.0] \\
\texttt{colsample\_bytree} & [0.6, 0.8, 1.0] \\
\texttt{gamma} & [0, 0.1, 0.5] \\
\texttt{min\_child\_weight} & [1, 3, 5] \\
\bottomrule
\end{tabular}
\end{table}

\paragraph{Class Imbalance Handling}

\texttt{scale\_pos\_weight} parameter set to ratio of negative-to-positive class frequencies:

\begin{equation}
\texttt{scale\_pos\_weight} = \frac{n_{\text{negative}}}{n_{\text{positive}}}
\end{equation}


\section{Evaluation Strategy}
\label{sec:eval_strategy}
% ==================================================================

A two-track evaluation strategy was adopted to assess model performance under
both idealised and deployment-realistic conditions.  The two tracks serve
complementary purposes: random split evaluation provides a performance ceiling
comparable to the existing literature, while temporal validation reveals how
models behave when required to forecast genuinely unseen future periods.
Using both protocols allows direct quantification of the performance gap
attributable to temporal non-stationarity.

\subsection{Random Split Evaluation}

The standard evaluation protocol in aviation delay prediction research employs
a stratified random train-test split, in which observations are randomly
assigned to training (75\%) and test (25\%) partitions without regard to their
chronological order.  Stratification on the binary delay label ensures that
the class distribution -- approximately 79.4\% delayed flights in Dataset V1
Imputed Fixed -- is preserved identically in both partitions, preventing
accidental imbalance from inflating or deflating performance estimates.

This protocol operates under the assumption that training and test observations
are drawn from the same underlying distribution.  While this assumption is
violated in practice by seasonal non-stationarity, random split evaluation
serves two important roles in the present study.  First, it provides an
upper-bound performance benchmark directly comparable to published results
from studies that also employ random splits, enabling fair assessment of the
proposed multi-source feature framework against state-of-the-art baselines.
Second, comparing random split and temporal split performance on the same
dataset quantifies the exact cost of ignoring temporal structure -- a
methodological contribution in its own right.

\subsection{Temporal Validation}

Temporal validation simulates operational deployment conditions by training
models exclusively on historically earlier data and evaluating them on
chronologically later observations that were entirely unseen during training.
The dataset was partitioned as follows:

\begin{itemize}
  \item \textbf{Training Set:} October--November 2025 (8,529 flights, 89.5\%
        of the dataset) -- the period for which the most complete and
        representative data are available.
  \item \textbf{Intermediate Validation:} December 2025 (565 flights, 6.0\%)
        -- treated as an early deployment probe to detect the onset of
        seasonal divergence.
  \item \textbf{Test Set:} January--February 2026 (433 flights, 4.5\%) --
        the target deployment period, separated from training by a minimum
        of two calendar months to prevent any indirect temporal leakage.
\end{itemize}

The strict chronological boundary ensures that no information from the test
period influences model training, feature engineering, or hyperparameter
selection -- faithfully replicating the conditions under which an operational
delay prediction system would be required to perform.  December is retained
as an intermediate period rather than incorporated into training because it
represents the seasonal transition between the autumn operational regime and
the winter regime observed in January--February; including it in training
would have partially exposed the model to the distributional shift it is
intended to detect.

\subsubsection*{Time-Series Cross-Validation}

During hyperparameter tuning, a 5-fold time-series cross-validation scheme
was applied to the training set using a forward-chaining approach: each
successive fold extends the training window while the validation fold always
consists of the immediately following time block.  This preserves the
directional flow of time within the training period and prevents the
hyperparameter search from inadvertently selecting configurations that benefit
from future information.  Standard $k$-fold cross-validation, which randomly
shuffles fold assignments, is inappropriate for time-series data because it
allows later observations to appear in training folds that precede earlier
test folds, producing optimistically biased hyperparameter estimates.


% ==================================================================
\section{Feature Importance Analysis (SHAP)}
\label{sec:shap_method}
% ==================================================================

Model interpretability was assessed using SHAP (SHapley Additive
exPlanations), a unified framework for attributing individual predictions to
their contributing features based on cooperative game-theoretic Shapley
values.  Unlike permutation-based importance measures, SHAP guarantees
consistency (larger feature contributions are always reflected by larger SHAP
magnitudes), local accuracy (attributions sum to the prediction), and
missingness (features absent from the model receive zero attribution).  These
properties make SHAP particularly well-suited to operational deployment
contexts where prediction explanations must be auditable and actionable.

SHAP values were computed using \texttt{shap.TreeExplainer}, an efficient
tree-specific implementation that exploits the recursive structure of
decision trees to compute exact Shapley values in polynomial rather than
exponential time.  The explainer was applied to the best-performing XGBoost
Tuned model on a stratified random sample of 500 test instances.  For each
feature $j$ and prediction instance $i$, the SHAP value is defined as:

\begin{equation}
\phi^{(i)}_j = \sum_{S \subseteq \mathcal{F} \setminus \{j\}}
\frac{|S|!\,(|\mathcal{F}| - |S| - 1)!}{|\mathcal{F}|!}
\left[ f_{S \cup \{j\}}(x^{(i)}) - f_S(x^{(i)}) \right]
\label{eq:shap}
\end{equation}

where $\mathcal{F}$ is the full feature set and $f_S(x^{(i)})$ denotes the
model output when only the features in subset $S$ are observed.  Intuitively,
$\phi^{(i)}_j$ measures the average marginal contribution of feature $j$
across all possible orderings in which features could be introduced.

Global feature importance is then summarised as the mean absolute SHAP value
across all evaluation instances:

\begin{equation}
\text{Importance}_j = \frac{1}{N} \sum_{i=1}^{N} \left| \phi^{(i)}_j \right|
\label{eq:shap_global}
\end{equation}

In addition to the global bar chart, a beeswarm plot was generated to
visualise both the magnitude and direction of each feature's influence,
allowing assessment of whether high or low feature values increase or
decrease predicted delay probability.


% ==================================================================
\section{Delay Propagation Analysis}
\label{sec:prop_method}
% ==================================================================

A core hypothesis of this study is that delays at Istanbul Airport propagate
through shared operational resources -- gates, aircraft, and crew -- causing
a delayed flight to elevate delay probability for subsequent flights.  To
quantify this contagion effect, a propagation analysis was conducted for each
flight $i$ departing at time $t_i$.

For every flight, the five chronologically nearest subsequent departures were
identified and their binary delay outcomes recorded.  Stratifying by the
delay status of the index flight (delayed vs.\ on-time), the contagion effect
is defined as:

\begin{equation}
\Delta_{\text{contagion}} =
\mathbb{E}[\text{NextFlightDelayRate} \mid \text{delayed}] -
\mathbb{E}[\text{NextFlightDelayRate} \mid \text{on-time}]
\label{eq:contagion}
\end{equation}

A positive $\Delta_{\text{contagion}}$ indicates that delayed flights are
followed by higher delay rates among subsequent departures, confirming
propagation through the system.  Statistical significance was assessed using
an independent-samples $t$-test, with the delay rates following delayed
flights compared against those following on-time flights.  The window of five
subsequent flights was chosen to capture the immediate operational footprint
of a single delay while limiting contamination from unrelated disruptions
occurring later in the schedule.

This analysis was conducted separately for the training period
(October--November 2025), the intermediate period (December 2025), and the
test period (January--February 2026) to detect seasonal variation in
propagation dynamics.


% ==================================================================
\section{Temporal Stability and Drift Analysis}
\label{sec:drift_method}
% ==================================================================

Non-stationarity -- the tendency for statistical relationships between
features and the target variable to shift over time -- represents a
fundamental threat to the operational viability of machine learning models.
To assess whether the delay propagation dynamics and feature-target
relationships learned during training remain valid in the test period, a
multi-component stability analysis was conducted.

The propagation effect $\Delta_{\text{contagion}}$ was computed independently
for the training, intermediate, and test periods.  Substantial divergence
across periods -- particularly a sign reversal -- would indicate that the
model has learned a relationship that no longer holds during deployment.
Complementing this, feature distribution shift was quantified by comparing
the marginal distributions of key predictors (historical delay features,
meteorological variables) between training and test periods using summary
statistics and visual inspection.  A shift in feature distributions combined
with a shift in feature-target relationships constitutes the full concept
drift scenario that most severely undermines model reliability.


% ==================================================================
\section{Additional Validation Experiments}
\label{sec:additional_val}
% ==================================================================

Beyond the primary evaluation tracks, four supplementary experiments were
conducted to stress-test model robustness and quantify the contribution of
individual feature categories.

\subsection{Stable Features Testing}

To investigate whether any feature subset retains predictive validity under
temporal distribution shift, models were retrained using only weather and
temporal features, deliberately excluding historical delay features.  The
rationale is that meteorological and circadian patterns are more likely to
remain stable across seasons than delay propagation rates, which can reverse
sign entirely between autumn and winter.  If stability-restricted models
achieve meaningfully better temporal validation performance, this would
suggest a deployment strategy based on feature selection rather than full
retraining.

\subsection{Feature Removal Experiments (Ablation)}

Systematic ablation experiments quantified the discriminative contribution of
each feature category by retraining models after removing one category at a
time -- weather features, historical delay features, and traffic features
respectively -- and measuring the resulting AUC degradation.  Ablation
provides a complementary perspective to SHAP importance: while SHAP measures
average per-prediction attribution, ablation measures the irreplaceable
discriminative signal that cannot be recovered from remaining features.

\subsection{Expanding Window Validation}

The training window was iteratively expanded from October--November to
October--December, and finally to October--January, with performance
evaluated on the remaining held-out period in each case.  This experiment
tests whether incorporating more recent training data -- and thus partially
bridging the seasonal gap -- improves generalisation to the test period, or
whether the regime shift is sufficiently abrupt that additional data provides
no meaningful benefit.

\subsection{Sliding Window Validation}

A fixed two-month training window was slid forward across the dataset --
October--November, November--December, December--January -- with performance
evaluated on the immediately following month.  This experiment isolates the
effect of training recency from training volume, revealing whether models
trained on the most recent available data consistently outperform models
trained on earlier windows of the same size.


% ==================================================================
\section{Error Analysis and Reliability Assessment}
\label{sec:error_method}
% ==================================================================

A structured error analysis was performed to characterise the conditions
under which the best model fails, providing operationally actionable insight
beyond aggregate performance metrics.

Prediction outcomes were classified into the four cells of the confusion
matrix (true positives, true negatives, false positives, false negatives),
and the mean feature values associated with each outcome type were examined
to identify systematic patterns in error-prone flight profiles.  Conditional
error rates were then stratified by hour of day and weather severity to
quantify how operational context modulates reliability.

Probability calibration was assessed by comparing predicted delay
probabilities against empirically observed delay frequencies across
ten equal-width bins spanning the 0.0--1.0 probability range.  A
well-calibrated model produces predictions that align closely with the
diagonal of the calibration plot, meaning a predicted probability of 0.7
corresponds to an empirically observed delay rate of approximately 70\%.
Good calibration is essential for threshold-based deployment decisions, as
it ensures that the chosen operating point corresponds to the intended
risk level.


% ==================================================================
\section{Evaluation Metrics}
\label{sec:metrics}
% ==================================================================

\subsection{Primary Metric: ROC-AUC}

The Area Under the Receiver Operating Characteristic Curve (ROC-AUC) was
adopted as the primary performance metric.  ROC-AUC measures the probability
that a randomly selected delayed flight receives a higher predicted
probability than a randomly selected on-time flight, and is computed as the
integral of the true positive rate over the false positive rate across all
classification thresholds:

\begin{equation}
\text{AUC} = \int_0^1 \text{TPR}(t)\, d[\text{FPR}(t)]
\label{eq:auc}
\end{equation}

ROC-AUC is well-suited to this application for three reasons: it is
threshold-independent, allowing fair comparison without committing to a
specific operating point; it is robust to class imbalance, an important
property given the 79.4\% delayed-flight prevalence in the dataset; and it
is the metric used in the benchmark studies against which this work is
compared, ensuring direct comparability.

\subsection{Secondary Metrics}

The following secondary metrics were computed at the default 0.5 decision
threshold to provide a complete operational picture:

\textbf{Precision} measures the fraction of predicted delays that are genuine
delays, reflecting the false alarm rate experienced by operational staff:
\begin{equation}
\text{Precision} = \frac{TP}{TP + FP}
\end{equation}

\textbf{Recall (Sensitivity)} measures the fraction of actual delays that are
correctly identified, reflecting the missed delay rate:
\begin{equation}
\text{Recall} = \frac{TP}{TP + FN}
\end{equation}

\textbf{F1-Score} is the harmonic mean of Precision and Recall, providing a
single balanced summary when both error types are operationally relevant:
\begin{equation}
F_1 = 2 \cdot \frac{\text{Precision} \cdot \text{Recall}}{\text{Precision} + \text{Recall}}
\end{equation}

\textbf{Brier Score} measures the mean squared error between predicted
probabilities and observed binary outcomes, penalising both poor
discrimination and poor calibration:
\begin{equation}
\text{Brier} = \frac{1}{N} \sum_{i=1}^{N} (\hat{p}_i - y_i)^2
\end{equation}
Lower Brier scores indicate better-calibrated probability estimates.  A
score of 0.25 corresponds to random prediction on a balanced dataset; lower
values confirm genuine predictive value.


% ==================================================================
\section{Reliability and Validity}
\label{sec:validity}
% ==================================================================

\subsection{Internal Validity}

Internal validity concerns whether the experimental results accurately
reflect the true predictive performance of the models, free from artefacts
introduced by the research procedure.

\textbf{Temporal Leakage Prevention.}  The most significant internal validity
threat in time-series machine learning is temporal data leakage, wherein
features inadvertently encode information from or about the target variable.
This was addressed through the strict chronological feature construction
protocol described in Section~\ref{sec:leakage}, with pre- and
post-correction validation confirming that leakage was eliminated.

\textbf{Consistent Preprocessing.}  All preprocessing transformations --
outlier removal, weather imputation, one-hot encoding, and feature scaling --
were fitted exclusively on training data and applied without refitting to test
data, preventing the test set from influencing any aspect of the pipeline.

\textbf{Cross-Validation Rigor.}  Time-series cross-validation during
hyperparameter tuning preserves temporal ordering within the training set,
ensuring that the selected hyperparameters are not optimised on future-leaking
validation folds.

\subsection{External Validity}

External validity concerns the degree to which findings generalise beyond the
specific dataset and conditions of this study.

\textbf{Temporal Generalisation.}  Evaluation on a chronologically separated
test period (January--February 2026) provides a realistic assessment of
deployment performance, going beyond the idealised conditions of random split
evaluation.  However, results are specific to the autumn-to-winter transition;
generalisation across the full annual cycle cannot be confirmed without
multi-year data.

\textbf{Limitations to Generalisability.}  Three characteristics of the study
design constrain external validity:

\begin{itemize}
  \item \textbf{Single Airport:} Results reflect Istanbul Airport's specific
        operational characteristics -- Turkish Airlines' 51\% traffic share,
        high international connectivity, and hub-dominated delay dynamics.
        Direct transfer to point-to-point or smaller regional airports should
        not be assumed without validation.
  \item \textbf{Limited Temporal Scope:} The five-month observation window
        captures autumn-winter dynamics only; spring and summer operational
        patterns, which may exhibit different delay structures, are excluded.
  \item \textbf{Traffic Data Simplification:} Representative hourly traffic
        profiles averaged across six months eliminate day-to-day variability
        from accidents, weather-traffic coupling, and special events.  This
        simplification may have attenuated any genuine traffic signal.
\end{itemize}


% ==================================================================
\section{Limitations of the Methodology}
\label{sec:method_limitations}
% ==================================================================

\subsection{Data Limitations}

\textbf{Missing Departure Times.}  Approximately 77\% of raw API records
lacked actual departure timestamps, reducing the dataset from 42,623
potential flights to 9,527 usable observations.  While Missing Completely at
Random (MCAR) analysis found no systematic correlation between missingness
and observable flight characteristics, the substantial exclusion rate means
that absolute delay rate estimates (79.4\%) may not be representative of the
full operational picture, and certain rare flight categories may be
under-represented.

\textbf{Weather Imputation Uncertainty.}  The nearest-neighbour temporal
imputation strategy, while validated through autocorrelation analysis and
leave-one-out error quantification, introduces measurement uncertainty that
is not propagated through model training.  Imputed weather values are treated
as exact observations, which may cause the model to express higher confidence
in weather-influenced predictions than is warranted.  Rapidly changing weather
during frontal passages represents the scenario most likely to violate the
imputation's temporal stability assumptions.

\textbf{Traffic Data Granularity.}  Surface traffic profiles were constructed
as representative 24-hour averages across a six-month historical window rather
than day-specific measurements.  This approach captures systematic circadian
patterns but eliminates variability from accidents, adverse weather-traffic
coupling, and major events.  While the observed near-zero correlations suggest
this simplification did not mask a genuine signal, it remains a constraint on
the strength of conclusions about traffic-delay relationships.

\subsection{Methodological Constraints}

\textbf{Single Airport Analysis.}  Focusing on a single hub airport provides
depth at the cost of breadth.  Hub-specific factors -- concentrated airline
operations, complex rotation schedules, and high international connectivity --
may drive the observed delay dynamics in ways that do not transfer to other
airport types.  Multi-airport validation is a necessary step before the
methodology can be treated as generally applicable.

\textbf{Limited Temporal Scope.}  The five-month observation window from
October 2025 to February 2026 captures the autumn-winter transition but
excludes spring and summer patterns.  Seasonal regime characterisation
requires data spanning at least one full annual cycle; the present study
provides evidence of one specific regime transition that may or may not be
representative of the airport's behaviour across all seasons.

\textbf{Binary Classification Framework.}  Framing delay prediction as a
binary outcome (delayed / on-time at the 15-minute threshold) discards
potentially valuable information about delay magnitude.  A 16-minute delay
and a 3-hour delay are treated identically, yet their operational impacts
differ substantially.  Regression-based or ordinal classification approaches
that predict delay severity would provide richer decision support at the cost
of increased modelling complexity.

\subsection{Temporal Drift Implications}

The discovery of a seasonal regime shift -- wherein delay contagion reverses
from positive in autumn to negative in winter -- confirms that delay dynamics
at Istanbul Airport are fundamentally non-stationary.  This finding has a
direct implication for the methodology itself: any evaluation conducted on a
single chronological period, including this study's temporal validation,
captures only a snapshot of a continuously evolving system.  While the
temporal validation protocol exposes the autumn-to-winter transition, it
cannot characterise the full range of possible regime states that an
operational deployment would encounter.  Adaptive retraining strategies that
continuously update model parameters as new observations arrive are therefore
not merely an optimisation but a structural requirement for sustained
deployment reliability.
















\chapter{Findings}
\label{chap:findings}

% ------------------------------------------------------------------
\section{Overview}
% ------------------------------------------------------------------

This chapter presents the experimental results obtained through systematic
application of the methodological framework detailed in Chapter~3.  The
presentation follows a progressive narrative structure that reflects the
iterative nature of the research process, beginning with initial baseline
experiments and culminating in comprehensive performance evaluation on an
expanded dataset.

The chapter is organised into eleven major sections.  Section~4.2 presents
baseline model performance on both Dataset~V1 (4,330 flights) and Dataset~V1
Imputed Fixed (9,527 flights), establishing model feasibility and quantifying
performance improvements from dataset expansion.  Section~4.3 documents
temporal validation results that reveal critical limitations when models
predict future periods.  Section~4.4 examines how performance changes across
different temporal windows.  Section~4.5 analyses feature distribution shifts
between training and test periods.  Section~4.6 investigates delay propagation
patterns and reveals a seasonal regime shift.  Section~4.7 provides model
interpretability through SHAP-based feature attribution.  Section~4.8
evaluates the impact of surface traffic features.  Section~4.9 examines
feature group contributions through ablation studies.  Section~4.10 tests
adaptive training strategies.  Section~4.11 conducts detailed error analysis
examining when and why the model fails.

This structure accomplishes three objectives: (1) it demonstrates research
transparency by documenting both successful outcomes and challenges encountered,
(2) it validates methodological decisions through empirical comparison of
dataset configurations, and (3) it provides comprehensive evidence supporting
the thesis contributions through multiple analytical perspectives.

% ==================================================================
\section{Baseline Results (Random Split Evaluation)}
\label{sec:baseline}
% ==================================================================

This section presents baseline model performance using standard random
train-test splits, establishing performance benchmarks before temporal
validation.  Results are presented for both Phase~1 (Dataset~V1, 4,330
flights) and Phase~2 (Dataset~V1 Imputed Fixed, 9,527 flights) to demonstrate
the impact of dataset expansion and weather imputation.

% ------------------------------------------------------------------
\subsection{Phase 1: Dataset V1 (4,330 Flights)}
% ------------------------------------------------------------------

The initial experimental phase utilised Dataset~V1, comprising 4,330 flights
departing from Istanbul Airport during October and November 2025.  This
dataset was constructed using direct temporal matching between flight records
and available weather observations, without application of imputation
techniques, achieving 45.4\% weather coverage.

\subsubsection*{Dataset Characteristics}

Dataset~V1 exhibits the characteristics presented in Table~\ref{tab:v1char}.

\begin{table}[H]
\centering
\caption{Dataset~V1 characteristics and coverage statistics}
\label{tab:v1char}
\begin{tabular}{ll}
\toprule
\textbf{Characteristic} & \textbf{Value} \\
\midrule
Total flights                  & 4,330 \\
Temporal span                  & Oct 27 -- Nov 30, 2025 \\
Study duration                 & 35 days \\
October flights                & 1,435 (33.1\%) \\
November flights               & 2,887 (66.7\%) \\
December flights               & 8 (0.2\%) \\
Weather observations (hourly)  & 377 hours \\
Weather coverage rate          & 45.4\% \\
Overall delay rate             & 56.2\% \\
October delay rate             & 51.6\% \\
November delay rate            & 58.3\% \\
Features per observation       & 20 base + airline dummies \\
Total dimensionality           & $\approx$461 features \\
\bottomrule
\end{tabular}
\end{table}

The 45.4\% weather coverage rate reflects the constraint that flights could
only be included if weather observations existed for their exact scheduled
departure hour.  Dataset~V1 exhibits moderate class imbalance with 56.2\% of
flights classified as delayed (departure delay $\geq$ 15 minutes).

\subsubsection*{Model Performance Comparison}

Five models were trained and evaluated on Dataset~V1 using a stratified
75-25\% random split.  Models ranged from a linear baseline (Logistic
Regression) to ensemble methods (Random Forest, XGBoost) with both default
and optimised hyperparameters.  Table~\ref{tab:v1perf} presents comprehensive
performance metrics.

\begin{table}[H]
\centering
\caption{Comprehensive performance metrics for all models on Dataset~V1
  (4,330 flights)}
\label{tab:v1perf}
\begin{tabular}{lccccc}
\toprule
\textbf{Model} & \textbf{ROC-AUC} & \textbf{Precision} &
  \textbf{Recall} & \textbf{F1} & \textbf{Brier} \\
\midrule
Logistic Regression      & 0.748 & 0.703 & 0.759 & 0.730 & 0.206 \\
Random Forest (untuned)  & 0.870 & 0.804 & 0.834 & 0.819 & 0.145 \\
Random Forest (tuned)    & 0.862 & 0.787 & 0.842 & 0.814 & 0.151 \\
XGBoost (untuned)        & 0.858 & 0.791 & 0.837 & 0.813 & 0.152 \\
XGBoost (tuned)          & 0.853 & 0.798 & 0.820 & 0.809 & 0.157 \\
\bottomrule
\end{tabular}
\end{table}

Figure~\ref{fig:v1_perf_roc} presents a side-by-side visual comparison:
the upper panel shows a grouped bar chart of the four main metrics for each
model, while the lower panel overlays the ROC curves for all five models on
the same axes, allowing direct comparison of discriminative capacity.

\begin{figure}[H]
  \centering
  \includegraphics[width=0.70\textwidth]{fig1_model_comparison.png}
  \includegraphics[width=0.70\textwidth]{fig2_roc_overlay.png}
  \caption{Model performance on Dataset~V1. \textit{Top:} grouped bar chart
    comparing ROC-AUC, Precision, Recall, and F1-score for all five models.
    \textit{Bottom:} overlaid ROC curves (Logistic Regression AUC~= 0.735;
    RF untuned~= 0.870; RF tuned~= 0.862; XGB untuned~= 0.858;
    XGB tuned~= 0.853).}
  \label{fig:v1_perf_roc}
\end{figure}

\textbf{Logistic Regression Baseline}\quad Logistic Regression with L2
regularisation serves as the interpretable baseline, establishing a
performance floor against which nonlinear methods are compared.  The model
achieves ROC-AUC~= 0.748, indicating discrimination substantially better
than random guessing (0.50) but revealing the limitations of linear modelling
for this problem.

Figure~\ref{fig:prob_ablation} (top panel) illustrates the predicted
probability distribution for Logistic Regression.  Substantial overlap between
the delayed and on-time distributions is concentrated in the 0.4--0.6
probability range, explaining the model's limited precision and recall at
standard classification thresholds.

\begin{figure}[H]
  \centering
  \includegraphics[width=0.70\textwidth]{fig3_probability_distribution.png}
  \includegraphics[width=0.70\textwidth]{fig4_ablation_v1.png}
  \caption{\textit{Top:} Predicted probability distribution for Logistic
    Regression on Dataset~V1, showing density for on-time (blue) and
    delayed (orange) flights. \textit{Bottom:} Ablation study results on
    Dataset~V1 -- bar chart of ROC-AUC for each feature-removal scenario
    (all features, no weather, no road traffic, no historical features).}
  \label{fig:prob_ablation}
\end{figure}

\textbf{Random Forest Models}\quad Random Forest models capture nonlinear
relationships through recursive partitioning and reduce variance through
bootstrap aggregation.  Two configurations were evaluated: an untuned model
using scikit-learn defaults and a tuned model optimised via randomised
hyperparameter search.

\textit{Untuned Random Forest:} The untuned Random Forest (200~estimators,
unlimited depth, sqrt feature sampling) achieves ROC-AUC~= 0.870 on
Dataset~V1, representing a 16.3\% improvement over Logistic Regression
(0.748).  The sharp initial rise of the ROC curve indicates confident
identification of a substantial proportion of delayed flights (high recall)
while maintaining low false positive rates (high specificity).

\textit{Tuned Random Forest:} Hyperparameter optimisation via 5-fold time
series cross-validation identified the configuration presented in
Table~\ref{tab:rf_hp}.  The tuned model achieves ROC-AUC~= 0.862, marginally
lower than the untuned configuration.

\begin{table}[H]
\centering
\caption{Hyperparameter configurations for Random Forest models on
  Dataset~V1}
\label{tab:rf_hp}
\begin{tabular}{lcc}
\toprule
\textbf{Hyperparameter} & \textbf{Untuned (Default)} &
  \textbf{Tuned (Optimised)} \\
\midrule
\texttt{n\_estimators}       & 200  & 300  \\
\texttt{max\_depth}          & None & 30   \\
\texttt{min\_samples\_split} & 2    & 5    \\
\texttt{min\_samples\_leaf}  & 1    & 2    \\
\texttt{max\_features}       & sqrt & sqrt \\
CV ROC-AUC                   & --   & 0.867 \\
Test ROC-AUC                 & 0.870 & 0.862 \\
\bottomrule
\end{tabular}
\end{table}

The tuned configuration introduced stronger regularisation
(\texttt{max\_depth=30}, \texttt{min\_samples\_split=5}) that reduced
cross-validation overfitting but marginally decreased test set
generalisation.  The 0.008 AUC difference falls within expected performance
variability.  Nevertheless, subsequent detailed analyses (ablation,
calibration, threshold sensitivity) were conducted using the tuned model to
ensure methodological rigour through cross-validation-based selection.

\textbf{XGBoost Models}\quad Gradient boosting models were evaluated for
their capacity to capture complex interactions through sequential error
correction.  Both untuned and tuned XGBoost configurations were assessed.

\textit{Untuned XGBoost:} The untuned XGBoost model achieves
ROC-AUC~= 0.858, performing comparably to Random Forest.

\textit{Tuned XGBoost:} Hyperparameter optimisation of XGBoost yielded
ROC-AUC~= 0.853, exhibiting a marginal decrease similar to the Random Forest
tuning pattern.  This consistency across model families indicates
well-calibrated default settings for this dataset size and feature
dimensionality.

\subsubsection*{Comparative Performance Analysis}

Key findings from Table~\ref{tab:v1perf}:

\begin{enumerate}
\item \textbf{Random Forest Superiority:} The untuned Random Forest achieves
  the highest ROC-AUC (0.870), precision (0.804), F1-score (0.819), and best
  calibration (Brier~= 0.145).
\item \textbf{Substantial Nonlinear Gains:} All ensemble methods achieve
  14--16\% AUC improvement over Logistic Regression (0.748), confirming
  nonlinear delay patterns.
\item \textbf{Ensemble Method Parity:} Random Forest and XGBoost demonstrate
  comparable performance (0.85--0.87 AUC range), with Random Forest holding
  marginal advantage.
\item \textbf{Tuning Paradox Consistency:} Both Random Forest and XGBoost
  exhibit marginal performance decreases after hyperparameter optimisation,
  indicating well-calibrated default settings.
\item \textbf{Calibration Quality:} Ensemble methods achieve substantially
  better probability calibration (Brier 0.145--0.157) than Logistic
  Regression (0.206).
\end{enumerate}

\subsubsection*{Statistical Validation via Bootstrap Resampling}

To confirm that observed performance differences reflect genuine model
behaviour rather than test set sampling variability, bootstrap analysis with
1,000 iterations was conducted on the Random Forest models.
Table~\ref{tab:bootstrap} presents results.

\begin{table}[H]
\centering
\caption{Bootstrap validation of Random Forest performance on Dataset~V1
  (1,000 resampling iterations)}
\label{tab:bootstrap}
\begin{tabular}{lccc}
\toprule
\textbf{Model} & \textbf{Mean AUC} & \textbf{95\% CI} &
  \textbf{Std Dev} \\
\midrule
Random Forest (untuned) & 0.870 & [0.848, 0.888] & 0.010 \\
Random Forest (tuned)   & 0.862 & [0.840, 0.881] & 0.011 \\
Difference              & 0.008 & [0.003, 0.013] & --    \\
\bottomrule
\end{tabular}
\end{table}

The bootstrap confidence interval for the performance difference
[0.003, 0.013] excludes zero, confirming statistical significance at
$\alpha = 0.05$.  This validates that the counterintuitive tuning result is
reproducible rather than an artefact of random test set composition.

\subsubsection*{Threshold Sensitivity Analysis}

Standard classification metrics employ a default probability threshold of
0.5, but this choice is arbitrary and may not align with operational
priorities.  To examine this trade-off, the tuned Random Forest model's
predicted probabilities were evaluated across thresholds from 0.30 to 0.65.

\begin{figure}[H]
\centering
\begin{subfigure}{0.48\textwidth}
    \centering
    \includegraphics[width=\linewidth]{precision_recall_vs_threshold.png}
    \caption{Precision-Recall curve}
\end{subfigure}%
\hfill
\begin{subfigure}{0.48\textwidth}
    \centering
    \includegraphics[width=\linewidth]{threshold_error_tradeoff.png}
    \caption{Threshold sensitivity analysis}
\end{subfigure}
\caption{Left: Precision-Recall trade-off curve for the tuned Random Forest on Dataset~V1. 
Right: Threshold sensitivity analysis showing false positives and false negatives across decision thresholds (0.30–0.65).}
\label{fig:pr_threshold}
\end{figure}

Results in Figure~\ref{fig:pr_threshold} reveal stark trade-offs:

\begin{enumerate}
\item \textbf{Asymmetric Error Costs:} At threshold 0.40, false negatives
  drop substantially but false positives rise sharply.  At threshold 0.60,
  false positives fall significantly but false negatives surge.  This
  demonstrates extreme threshold sensitivity.
\item \textbf{No Optimal Default:} The standard 0.5 threshold yields neither
  the minimum missed delays nor the minimum false alarms -- it represents
  an arbitrary compromise.
\end{enumerate}

These findings demonstrate that model evaluation cannot rely solely on
threshold-independent metrics (ROC-AUC).  Real-world deployment requires
explicit cost-benefit analysis to determine operationally appropriate
decision thresholds.

\subsubsection*{Probability Calibration Assessment}

Beyond discrimination and classification accuracy, operational utility
requires that predicted probabilities represent meaningful risk estimates.
Well-calibrated models enable stakeholders to interpret predictions as
approximate delay likelihoods rather than opaque scores.

\begin{figure}[H]
  \centering
  \includegraphics[width=0.70\textwidth]{calibration_curve_rf_v1.png}
  \caption{Calibration curve for the tuned Random Forest on Dataset~V1.
    The model curve (solid blue) closely tracks the perfect calibration
    diagonal (dashed), confirming that predicted delay probabilities
    accurately reflect empirical delay frequencies across the 0.2--1.0
    probability range.}
  \label{fig:calib_v1}
\end{figure}

\textbf{Calibration Quality:} The calibration curve for the tuned Random
Forest on Dataset~V1 (Figure~\ref{fig:calib_v1}) closely tracks the perfect
calibration diagonal across the 0.3--0.8 probability range.  The Brier
score of 0.151 confirms moderate-to-good calibration quality.  No post-hoc
calibration (e.g., Platt scaling) was applied; the reported calibration
reflects native Random Forest output.  This indicates that probability
estimates are sufficiently reliable for risk-based operational
decision-making.

\subsubsection*{Ablation Study: Feature Group Contributions}

To quantify feature group contributions beyond correlation-based importance,
systematic ablation experiments were conducted using the tuned Random Forest
model.  Four configurations were evaluated: full feature set, and three
reduced sets excluding weather, road traffic, or historical features
respectively.  Results are shown in the lower panel of
Figure~\ref{fig:prob_ablation} and in Table~\ref{tab:v1ablation}.

\begin{table}[H]
\centering
\caption{Ablation study results on Dataset~V1 using tuned Random Forest}
\label{tab:v1ablation}
\begin{tabular}{lccr}
\toprule
\textbf{Configuration} & \textbf{ROC-AUC} & \textbf{Recall} &
  \textbf{False Negatives} \\
\midrule
All features            & 0.863 & 0.842 & 96  \\
No weather features     & 0.841 & 0.829 & 104 \\
$\Delta$ vs.\ full      & $-$0.022 & $-$0.013 & $+8$  \\
No road traffic         & 0.856 & 0.834 & 101 \\
$\Delta$ vs.\ full      & $-$0.006 & $-$0.008 & $+5$  \\
No historical features  & 0.871 & 0.845 & 94  \\
$\Delta$ vs.\ full      & $+$0.008 & $+$0.003 & $-2$  \\
\bottomrule
\end{tabular}
\end{table}

\textbf{Weather Features -- Strongest Contribution:} Removing weather
features produces the largest performance degradation: ROC-AUC drops 2.2
percentage points (0.863 $\to$ 0.841) and false negatives increase by 8
(96 $\to$ 104).

\textbf{Road Traffic -- Marginal Contribution:} Road traffic feature removal
yields modest degradation: ROC-AUC decreases 0.6 percentage points.

\textbf{Historical Features -- Unexpected Negative Contribution:}
Counterintuitively, removing historical delay features \emph{improves}
performance: ROC-AUC increases 0.8 percentage points (0.863 $\to$ 0.871)
and false negatives decrease by 2.  This suggests that, on this dataset,
explicit historical features may introduce noise or overfitting artefacts
that hurt generalisation.

\subsubsection*{Summary of Dataset V1 Findings}

Experimental results on the initial 4,330-flight dataset establish several
foundational conclusions:

\begin{enumerate}
\item \textbf{Feasibility Validated:} The multi-source machine learning
  framework successfully predicts short-term flight delays, with Random
  Forest achieving 0.870 ROC-AUC -- substantially exceeding the linear
  baseline (0.748).
\item \textbf{Nonlinearity Essential:} Ensemble methods achieve 14--16\%
  AUC improvements over Logistic Regression.
\item \textbf{Weather Features Critical:} Ablation study reveals that
  meteorological variables provide the strongest predictive signal.
\item \textbf{Road Traffic Marginal:} Surface transportation congestion
  demonstrates limited direct predictive value (0.6 pp AUC contribution).
\item \textbf{Historical Features Paradox:} In this dataset, explicit
  historical delay features appear to degrade rather than improve performance.
\item \textbf{Threshold Sensitivity Critical:} Default 0.5 threshold is
  arbitrary and suboptimal; threshold selection dramatically affects
  operational outcomes.
\item \textbf{Calibration Adequate:} Model produces well-calibrated
  probability estimates (Brier~= 0.151) suitable for risk-based
  interpretation.
\end{enumerate}

% ------------------------------------------------------------------
\subsection{Phase 2: Dataset V1 Imputed Fixed (9,527 Flights)}
% ------------------------------------------------------------------

Following implementation of weather imputation methodology and dataset
expansion, Dataset~V1 Imputed Fixed comprises 9,527 flights spanning five
months (October 2025 through February 2026), with 68.2\% weather coverage
achieved through conservative 4-hour tolerance nearest-neighbour weather
matching.

\subsubsection*{Dataset Characteristics}

Table~\ref{tab:v2char} presents Dataset~V1 Imputed Fixed characteristics
alongside Phase~1 for direct comparison.

\begin{table}[H]
\centering
\caption{Dataset~V1 Imputed Fixed characteristics and comparison to
  Dataset~V1}
\label{tab:v2char}
\begin{tabular}{lrr}
\toprule
\textbf{Characteristic} & \textbf{V1} & \textbf{V1 Imputed Fixed} \\
\midrule
Total flights              & 4,330  & 9,527  \\
Temporal span (months)     & 2.1    & 5.0    \\
October flights            & 1,435  & 5,329  \\
November flights           & 2,887  & 3,200  \\
December flights           & 8      & 565    \\
January flights            & 0      & 241    \\
February flights           & 0      & 192    \\
Weather coverage method    & Exact  & 4h tolerance \\
Weather coverage rate      & 45.4\% & 68.2\% \\
Overall delay rate         & 56.2\% & 79.4\% \\
\bottomrule
\end{tabular}
\end{table}

\subsubsection*{Model Performance Comparison}

All models were retrained on Dataset~V1 Imputed Fixed using identical random
split evaluation protocol (75-25\% stratified split).
Table~\ref{tab:v2perf} presents comprehensive performance metrics.

\begin{table}[H]
\centering
\caption{Model performance on Dataset~V1 Imputed Fixed (9,527 flights,
  random split)}
\label{tab:v2perf}
\begin{tabular}{lccccc}
\toprule
\textbf{Model} & \textbf{ROC-AUC} & \textbf{Precision} &
  \textbf{Recall} & \textbf{F1} & \textbf{Brier} \\
\midrule
Logistic Regression & 0.782 & 0.821 & 0.892 & 0.855 & 0.168 \\
Random Forest       & 0.954 & 0.912 & 0.966 & 0.938 & 0.076 \\
XGBoost (untuned)   & 0.885 & 0.844 & 0.943 & 0.891 & 0.124 \\
XGBoost (tuned)     & 0.903 & 0.865 & 0.941 & 0.901 & 0.115 \\
\bottomrule
\end{tabular}
\end{table}

Figure~\ref{fig:v2_roc_collapse} presents the Phase~2 ROC curves (upper
panel) and a compelling visualisation of temporal collapse (lower panel):
the XGBoost Tuned model achieves 0.903 AUC on randomly split data but
degrades catastrophically to near-random performance (0.514 in December,
0.474 in January--February) when evaluated in strict chronological order.

\begin{figure}[H]
\centering

\begin{subfigure}{0.48\textwidth}
    \centering
    \includegraphics[width=\linewidth]{fig8_roc_phase2.png}
    \caption{ROC curves for all models}
\end{subfigure}%
\hfill
\begin{subfigure}{0.48\textwidth}
    \centering
    \includegraphics[width=\linewidth]{fig9_temporal_collapse.png}
    \caption{Temporal collapse of XGBoost Tuned}
\end{subfigure}

\caption{Left: ROC curves for all models on Dataset~V1 Imputed Fixed (LR AUC~= 0.782; RF~= 0.954; XGB untuned~= 0.885; XGB tuned~= 0.903). 
Right: Temporal collapse of XGBoost Tuned -- AUC drops from 0.903 (random split) to 0.514 (December, +1 month) to 0.474 (Jan--Feb, +3–4 months), approaching random baseline (dashed red line at 0.50).}
\label{fig:v2_roc_collapse}

\end{figure}

\textbf{Remarkable Performance Gains}\quad Dataset~V1 Imputed Fixed yields
dramatic performance improvements across all ensemble models:

\begin{enumerate}
\item \textbf{XGBoost Tuned:} ROC-AUC increases from 0.853 (Phase~1) to
  0.903 (Phase~2) -- a 5.0 percentage point gain (+5.9\% relative).
\item \textbf{Random Forest:} Achieves 0.954 ROC-AUC, representing 8.4
  percentage point gain over Phase~1 (0.870).
\item \textbf{XGBoost Untuned:} Improves from 0.858 to 0.885 (+2.7 pp).
\item \textbf{State-of-the-Art Performance:} XGBoost Tuned (0.903 AUC)
  exceeds published benchmarks (Gui et al.\ 2019: 0.88 AUC).
\end{enumerate}

Table~\ref{tab:phase_comp} quantifies performance improvement from Phase~1
to Phase~2.  XGBoost Tuned (0.903 AUC) was selected as the primary model
for subsequent detailed analyses due to its combination of strong
performance and faster inference compared to Random Forest.

\begin{table}[H]
\centering
\caption{Performance improvement from Phase~1 to Phase~2}
\label{tab:phase_comp}
\begin{tabular}{lccl}
\toprule
\textbf{Model} & \textbf{Phase 1 AUC} & \textbf{Phase 2 AUC} &
  \textbf{Improvement} \\
\midrule
Logistic Regression & 0.748 & 0.782 & $+$0.034 ($+$4.5\%) \\
Random Forest       & 0.870 & 0.954 & $+$0.084 ($+$9.7\%) \\
XGBoost (untuned)   & 0.858 & 0.885 & $+$0.027 ($+$3.1\%) \\
XGBoost (tuned)     & 0.853 & 0.903 & $+$0.050 ($+$5.9\%) \\
\bottomrule
\end{tabular}
\end{table}

% ==================================================================
\section{Temporal Validation Results}
\label{sec:temporal_val}
% ==================================================================

While random-split evaluation demonstrated strong performance (0.90+ AUC),
operational deployment requires models to predict future flights not yet
observed during training.  This section presents temporal validation results
using strict chronological train-test partitioning.

\subsection{Temporal Partitioning Strategy}

Dataset~V1 Imputed Fixed was partitioned chronologically:

\begin{itemize}
\item \textbf{Training:} October--November 2025 (8,529 flights, 89.5\%)
\item \textbf{Intermediate Validation:} December 2025 (565 flights, 6.0\%)
\item \textbf{Test:} January--February 2026 (433 flights, 4.5\%)
\end{itemize}

Models trained exclusively on Oct--Nov data were evaluated on December
(1 month ahead) and January--February (3-4 months ahead).  This strict
past $\to$ future split prevents any temporal leakage and simulates
realistic operational deployment.

\subsection{Temporal Validation Performance}

Table~\ref{tab:temporal_perf} presents model performance across temporal
periods.

\begin{table}[H]
\centering
\caption{Model performance across temporal periods (trained on Oct--Nov
  2025)}
\label{tab:temporal_perf}
\begin{tabular}{lccc}
\toprule
\textbf{Model} & \textbf{Oct--Nov (Random)} & \textbf{December} &
  \textbf{Jan--Feb} \\
\midrule
Logistic Regression & 0.782 & 0.501 & 0.491 \\
Random Forest       & 0.954 & 0.518 & 0.507 \\
XGBoost (untuned)   & 0.885 & 0.497 & 0.468 \\
XGBoost (tuned)     & 0.903 & 0.514 & 0.474 \\
\bottomrule
\end{tabular}
\end{table}

\subsection{Catastrophic Temporal Degradation}

Results reveal dramatic performance collapse:

\begin{enumerate}
\item \textbf{Random Split Performance:} All models achieved strong
  discrimination on randomly shuffled Oct--Nov data (0.78--0.95 AUC).
\item \textbf{December Performance (1 month ahead):} All models degraded to
  near-random performance (AUC 0.497--0.518), barely better than coin
  flipping (0.50).
\item \textbf{January--February Performance (3-4 months ahead):}
  Performance deteriorated further (AUC 0.468--0.507), with XGBoost Untuned
  falling below random baseline (0.468).
\item \textbf{Consistent Cross-Model Failure:} All four model families
  exhibited similar degradation patterns, indicating the problem lies in
  fundamental data distribution shift rather than model architecture.
\end{enumerate}

This catastrophic degradation, visible in the lower panel of
Figure~\ref{fig:v2_roc_collapse}, contrasts sharply with strong random-split
performance, revealing that models overfit to October--November specific
patterns that do not generalise to future months.

% ==================================================================
\section{Performance Changes Across Time}
\label{sec:perf_time}
% ==================================================================

To identify when and how predictive capacity degrades, model performance was
stratified by individual months.

\subsection{Monthly Stratification Analysis}

Table~\ref{tab:monthly} presents XGBoost Tuned performance by month.

\begin{table}[H]
\centering
\caption{XGBoost Tuned performance by month (trained on Oct--Nov)}
\label{tab:monthly}
\begin{tabular}{lrrccc}
\toprule
\textbf{Month} & \textbf{Flights} & \textbf{Delay Rate} & \textbf{AUC} &
  \textbf{Precision} & \textbf{Recall} \\
\midrule
Oct--Nov (train) & 8,529 & 79.5\% & 0.903 & 0.865 & 0.941 \\
December         &   565 & 78.4\% & 0.514 & 0.790 & 0.726 \\
January          &   241 & 76.8\% & 0.489 & 0.775 & 0.828 \\
February         &   192 & 77.6\% & 0.467 & 0.755 & 0.882 \\
\bottomrule
\end{tabular}
\end{table}

\begin{figure}[H]
  \centering
  \includegraphics[width=0.85\textwidth]{fig10_monthly_auc.png}
  \caption{Monthly performance degradation of XGBoost Tuned (trained on
    Oct--Nov 2025).  ROC-AUC collapses from 0.903 at training time to
    0.514 (December), 0.489 (January), and 0.467 (February), converging
    toward the random baseline (dashed red line at 0.50).}
  \label{fig:monthly_deg}
\end{figure}

\subsection{Key Observations}

\begin{enumerate}
\item \textbf{Immediate Collapse:} Performance degraded immediately when
  predicting December (1 month ahead), with AUC dropping from 0.903 to
  0.514 (Figure~\ref{fig:monthly_deg}).
\item \textbf{Stable Delay Rates:} Overall delay rates remained relatively
  stable across months (76.8\%--79.5\%), showing that performance collapse
  is not due to shifting delay prevalence.
\item \textbf{Maintained Recall:} While ROC-AUC approached random
  (0.47--0.51), recall remained moderately high (0.73--0.88), indicating
  the model still identifies delayed flights but loses discriminative ranking
  ability.
\item \textbf{Progressive Deterioration:} Performance did not stabilise but
  continued degrading from December (0.514) to January (0.489) to February
  (0.467).
\end{enumerate}

The stable delay rates combined with collapsed AUC reveal that the
fundamental relationship between features and delays changed between autumn
(Oct--Nov) and winter (Dec--Feb), rather than simple prevalence shift.

% ==================================================================
\section{Feature Drift Results}
\label{sec:feature_drift}
% ==================================================================

To investigate why temporal validation failed, feature distributions were
compared between training (Oct--Nov) and test (Jan--Feb) periods.

\subsection{Historical Delay Feature Drift}

Table~\ref{tab:hist_drift} presents historical delay feature distributions.

\begin{table}[H]
\centering
\caption{Historical delay feature distributions by period}
\label{tab:hist_drift}
\begin{tabular}{lrrr}
\toprule
\textbf{Feature} & \textbf{Oct--Nov Mean} & \textbf{Jan--Feb Mean} &
  \textbf{Drift} \\
\midrule
\texttt{prev\_hour\_delay\_rate}   & 0.647 & 0.572 & $-$0.075 ($-$7.5 pp) \\
\texttt{avg\_delay\_last\_24h}     & 0.774 & 0.689 & $-$0.085 ($-$8.5 pp) \\
\texttt{airline\_7d\_delay\_rate}  & 0.781 & 0.712 & $-$0.069 ($-$6.9 pp) \\
\bottomrule
\end{tabular}
\end{table}

\subsection{Weather Feature Drift}

Table~\ref{tab:weather_drift} presents weather variable distributions.

\begin{table}[H]
\centering
\caption{Weather variable distributions by period}
\label{tab:weather_drift}
\begin{tabular}{lrrr}
\toprule
\textbf{Variable} & \textbf{Oct--Nov Mean} & \textbf{Jan--Feb Mean} &
  \textbf{Drift} \\
\midrule
Temperature ($^\circ$C)  & 18.2   & 6.8    & $-$11.4 \\
Humidity (\%)            & 64.3   & 71.2   & $+$6.9  \\
Pressure (hPa)           & 1015.4 & 1018.7 & $+$3.3  \\
Wind Speed (m/s)         & 4.2    & 5.8    & $+$1.6  \\
\bottomrule
\end{tabular}
\end{table}

\begin{figure}[H]
  \centering
  \includegraphics[width=0.7\textwidth]{fig11b_weather_shift.png}
  \caption{Feature distribution shifts between Oct--Nov 2025 (training,
    blue) and Jan--Feb 2026 (test, red/orange).
    \textit{Top:} Historical delay feature drift -- all three propagation
    features (prev-hour delay rate, 24h average, airline 7-day rate)
    decreased 6.9--8.5 pp in winter despite stable overall delay prevalence.
    \textit{Bottom:} Weather variable drift -- temperature fell 11.4$^\circ$C,
    humidity rose 6.9 pp, pressure rose 3.3 hPa, and wind speed increased
    1.6 m/s.}
  \label{fig:feature_drift}
\end{figure}

\subsection{Key Observations}

\begin{enumerate}
\item \textbf{Substantial Temperature Shift:} Temperature exhibited an
  11.4$^\circ$C decrease between autumn and winter, representing a
  fundamental seasonal transition.
\item \textbf{Humidity Increase:} Humidity increased 6.9 percentage points,
  reflecting winter precipitation patterns.
\item \textbf{Pressure Increase:} Atmospheric pressure increased 3.3 hPa,
  consistent with seasonal synoptic patterns.
\item \textbf{Historical Feature Paradox:} Historical delay rates decreased
  despite stable overall delay prevalence, indicating that delay propagation
  mechanisms changed between seasons.
\end{enumerate}

Combined evidence shows substantial covariate shift between training and
test distributions, both in environmental variables (weather) and operational
patterns (historical delays).

% ==================================================================
\section{Delay Propagation Results}
\label{sec:propagation}
% ==================================================================

To understand delay generation mechanisms, this section analyses how delays
propagate through the flight network and whether propagation patterns vary
across airlines and time periods.

\subsection{Overall Propagation Effect}

For each flight in Dataset~V1 Imputed Fixed, the delay rate of the subsequent
5 chronologically nearest flights was computed and stratified by current
flight delay status.  Table~\ref{tab:propagation} presents results.

\begin{table}[H]
\centering
\caption{Overall delay propagation effect (9,527 flights)}
\label{tab:propagation}
\begin{tabular}{lrcc}
\toprule
\textbf{Current Flight Status} & \textbf{N} &
  \textbf{Next 5 Delay Rate} & \textbf{Std Error} \\
\midrule
Delayed  & 7,563 & 79.8\% & 0.004 \\
On-time  & 1,964 & 72.6\% & 0.010 \\
\textbf{Contagion Effect} & & \textbf{$+$7.2 pp} & $p < 0.0001$ \\
\bottomrule
\end{tabular}
\end{table}

This statistically significant contagion effect confirms that delays
propagate through shared airport resources (gates, crews, aircraft
rotations).

\subsection{Airline-Specific Propagation Heterogeneity}

Delay propagation was stratified by airline to examine whether operational
differences affect contagion strength.  Table~\ref{tab:contagion_airline} presents results for the top 10 carriers by flight volume.

\begin{table}[H]
\centering
\caption{Delay propagation by airline: top 10 carriers by contagion effect.
         \textit{After Delayed} = delay rate of the next flight operated by
         the same airline following a delayed departure;
         \textit{After On-time} = delay rate following an on-time departure;
         \textit{Contagion} = difference in percentage points (pp).
         Only airlines with at least 50 flights are included.}
\label{tab:contagion_airline}
\begin{tabular}{lrrrr}
\hline
\textbf{Airline} & \textbf{Flights} & \textbf{After Delayed (\%)} & \textbf{After On-time (\%)} & \textbf{Contagion (pp)} \\
\hline
Aegean Airlines            & 55  & 89.2 & 29.4 & +59.8 \\
SA AVIANCA                 & 71  & 91.2 & 38.5 & +52.8 \\
JetBlue Airways            & 63  & 90.0 & 41.7 & +48.3 \\
EgyptAir                   & 173 & 87.0 & 43.9 & +43.1 \\
ANA                        & 88  & 93.5 & 60.0 & +33.5 \\
Thai Airways International & 572 & 85.8 & 53.7 & +32.1 \\
Royal Air Maroc            & 79  & 89.1 & 57.1 & +31.9 \\
Miat -- Mongolian Airlines & 91  & 87.8 & 56.2 & +31.6 \\
Sichuan Airlines           & 118 & 87.2 & 56.5 & +30.7 \\
Air Serbia                 & 354 & 76.4 & 45.8 & +30.6 \\
\hline
\end{tabular}
\end{table}

\begin{figure}[H]
    \centering
    \includegraphics[width=0.75\textwidth]{fig_airline_propagation_updated.png}
    \caption{Airline-specific delay propagation (contagion effect in percentage
    points) for the top 10 carriers at Istanbul Airport. All airlines display
    positive contagion, indicating that a delayed departure is consistently
    followed by a higher probability of delay on the airline's next flight.
    Aegean Airlines exhibits the strongest effect (+59.8 pp) while Air Serbia
    shows the lowest among the top 10 (+30.6 pp). Only airlines with at least
    50 flights in the study period are included.}
    \label{fig:airline_propagation}
\end{figure}

Table~\ref{tab:contagion_airline} presents the ten airlines exhibiting the
strongest delay contagion effect at Istanbul Airport. A striking finding is
that \textit{all} airlines in the dataset display positive contagion values,
meaning that a delayed departure is consistently followed by a higher
probability of delay on the airline's next flight than an on-time departure
would be. This suggests that delay propagation is a systemic characteristic
of airline operations at the airport rather than an isolated phenomenon
confined to a subset of carriers.

The magnitude of the effect varies substantially across airlines. Aegean
Airlines shows the highest contagion at $+59.8$ pp: following a delayed
Aegean flight, 89.2\% of subsequent Aegean departures are also delayed,
compared with only 29.4\% following an on-time departure. Similarly strong
effects are observed for SA AVIANCA ($+52.8$ pp) and JetBlue Airways
($+48.3$ pp). These three carriers operate relatively few flights at Istanbul
Airport (55--71 flights in the study period), which may amplify variability,
but the consistent direction of the effect across airlines lends credibility
to the finding.

Among higher-volume carriers, Thai Airways International (572 flights) and
Air Serbia (354 flights) still display substantial contagion of $+32.1$ pp
and $+30.6$ pp respectively, indicating that the propagation effect is not
an artefact of small sample sizes. Turkish Airlines, the dominant carrier at
Istanbul Airport with 3{,}498 flights in the dataset, exhibits a contagion
of $+21.0$ pp (82.1\% after a delayed flight versus 61.0\% after an on-time
flight), confirming that schedule disruptions on the largest operator have
measurable knock-on effects on its own subsequent departures. Taken together,
these results reinforce the importance of capturing airline-level delay history
as a predictive feature, as reflected in the \texttt{airline\_7d\_delay\_rate}
variable included in the modelling pipeline.




\subsection{Delay Propagation Decay Pattern}

\begin{figure}[H]
  \centering
  \includegraphics[width=0.80\textwidth]{propagation_decay.png}
  \caption{Delay propagation decay pattern: contagion effect (in percentage
    points) as a function of flight position in the subsequent sequence.
    The effect is strongest in the immediate next flights and dissipates
    as the airport schedule recovers, though with a notable peak around
    position 5.}
  \label{fig:prop_decay}
\end{figure}

Figure~\ref{fig:prop_decay} illustrates how the propagation effect evolves
across the subsequent 10 flights.  The contagion starts around +6.8 pp for
the immediately following flight, fluctuates before peaking near +9.6 pp at
position 5, then decays to approximately +2.1 pp by position 10 -- consistent
with gradual dissipation of delay cascades as airport operations recover.

\subsection{Seasonal Regime Shift in Propagation}

\textbf{Critical Discovery:} Analysis of propagation effect stratified by
temporal period revealed a complete reversal of delay propagation dynamics
across seasons.  Table~\ref{tab:seasonal_prop} presents results.

\begin{table}[H]
\centering
\caption{Delay propagation by temporal period}
\label{tab:seasonal_prop}
\begin{tabular}{lrcc}
\toprule
\textbf{Period} & \textbf{Flights} & \textbf{Contagion Effect} &
  \textbf{p-value} \\
\midrule
Oct--Nov (Training) & 8,529 & $+$7.07 pp  & $<$0.0001 \\
December            &   565 & $+$17.10 pp & $<$0.0001 \\
Jan--Feb (Test)     &   433 & $-$7.78 pp  & 0.0312    \\
\textbf{Total Swing} & & \textbf{14.84 pp} & \\
\bottomrule
\end{tabular}
\end{table}

\begin{figure}[H]
  \centering
  \includegraphics[width=0.80\textwidth]{fig14_seasonal_shift.png}
  \caption{Seasonal regime shifts in delay propagation dynamics.  The
    contagion effect reverses completely across temporal periods:
    $+$7.07 pp in autumn training (Oct--Nov, orange), $+$17.10 pp in
    the December transition (red), and $-$7.78 pp in the January--February
    test period (teal) -- a total swing of 14.84 percentage points.}
  \label{fig:seasonal_regime}
\end{figure}

\subsubsection*{Critical Finding: Seasonal Regime Shift}

\begin{enumerate}
\item \textbf{October--November (Training):} Moderate positive contagion
  ($+$7.07 pp), consistent with operational cascades through shared
  resources.
\item \textbf{December (Transition):} Intensified propagation
  ($+$17.10 pp), indicating heightened operational stress during the
  transition period.
\item \textbf{January--February (Test):} Complete reversal to negative
  contagion ($-$7.78 pp), meaning delayed flights were followed by
  \emph{lower} delay rates than on-time flights.
\item \textbf{Total Swing:} 14.84 percentage point swing represents a
  fundamental shift in delay generation and propagation mechanisms.
\end{enumerate}

This seasonal regime shift provides a mechanistic explanation for temporal
validation failure.  Models trained on autumn operational patterns
($+$7.07 pp propagation) learned that historical delay features predict
future delays.  However, in winter, this relationship reversed ($-$7.78 pp),
causing historical delay features to become adversarial predictors --
exactly opposite their training behaviour.

% ==================================================================
\section{Feature Importance Results (SHAP)}
\label{sec:shap}
% ==================================================================

SHAP (SHapley Additive exPlanations) analysis provides theoretically
grounded feature attribution by quantifying each feature's marginal
contribution to predictions while accounting for feature interactions.
SHAP values were computed for XGBoost Tuned predictions on a random sample
of 500 test instances.

\subsection{Global Feature Importance}

Table~\ref{tab:shap_top10} presents the top features ranked by mean absolute
SHAP value.  These results are derived from the corrected model that
excludes destination airport codes in accordance with the feature selection
protocol detailed in Section~3.3.2.

\begin{table}[H]
\centering
\caption{Top features by mean absolute SHAP value (corrected model)}
\label{tab:shap_top10}
\begin{tabular}{lcc}
\toprule
\textbf{Feature} & \textbf{Mean $|\text{SHAP}|$} & \textbf{Category} \\
\midrule
\texttt{airline\_7d\_delay\_rate}  & 0.50 & Historical \\
\texttt{avg\_delay\_last\_24h}     & 0.42 & Historical \\
\texttt{prev\_hour\_delay\_rate}   & 0.30 & Historical \\
\texttt{wind\_speed}               & 0.27 & Weather    \\
\texttt{pressure}                  & 0.26 & Weather    \\
\texttt{hour\_of\_day}             & 0.19 & Temporal   \\
\texttt{humidity}                  & 0.16 & Weather    \\
\texttt{temperature}               & 0.16 & Weather    \\
\texttt{road\_avg\_speed}          & 0.14 & Traffic    \\
\texttt{road\_vehicle\_count}      & 0.13 & Traffic    \\
\bottomrule
\end{tabular}
\end{table}

\begin{figure}[H]
  \centering
  \includegraphics[width=0.85\textwidth]{fig15_shap_bar_clean.png}
  \caption{SHAP global feature importance bar chart (corrected model):
    mean absolute SHAP values for the top features.
    \texttt{airline\_7d\_delay\_rate} ($\approx$0.50) is the dominant
    predictor, followed by \texttt{avg\_delay\_last\_24h} ($\approx$0.42)
    and \texttt{prev\_hour\_delay\_rate} ($\approx$0.30).  Weather features
    (wind speed, pressure) rank fourth and fifth.}
  \label{fig:shap_bar}
\end{figure}

\begin{figure}[H]
  \centering
  \includegraphics[width=0.85\textwidth]{fig15_shap_summary_clean.png}
  \caption{SHAP beeswarm plot (corrected model): each dot represents one
    prediction instance, coloured by feature value (red~= high,
    blue~= low).  Features are ordered by mean $|\text{SHAP}|$.  High
    values of \texttt{airline\_7d\_delay\_rate} and
    \texttt{avg\_delay\_last\_24h} push predictions strongly toward delay
    (positive SHAP), while high wind speed and low pressure also increase
    delay probability.}
  \label{fig:shap_beeswarm}
\end{figure}

\subsubsection*{Key Observations}

\begin{enumerate}
\item \textbf{Historical Delay Dominance:} The carrier-specific 7-day
  rolling delay rate (\texttt{airline\_7d\_delay\_rate},
  mean $|\text{SHAP}| \approx 0.50$) emerges as the single most
  influential feature, substantially exceeding all others.  This indicates
  that airline-level operational history is the strongest predictor of
  departure delays at Istanbul Airport.
\item \textbf{Historical Feature Cluster:} The top three features are all
  historical delay propagation indicators, collectively capturing delay
  contagion at airline, daily, and hourly scales.
\item \textbf{Weather Features Secondary:} Wind speed (0.27) and atmospheric
  pressure (0.26) rank fourth and fifth -- important but secondary to
  historical delay momentum.
\item \textbf{Temporal Features Moderate:} Hour of day (0.19) captures
  circadian operational patterns, reflecting morning departure banks and
  evening congestion.
\item \textbf{Traffic Features Present but Weak:} Road speed and vehicle
  count appear in the top 10 with values of 0.14 and 0.13 respectively --
  measurable but operationally minor contributions.
\end{enumerate}

The beeswarm plot (Figure~\ref{fig:shap_beeswarm}) confirms directionality:
high values of \texttt{airline\_7d\_delay\_rate} and
\texttt{avg\_delay\_last\_24h} consistently push predictions toward delay,
while low wind speed and higher pressure are associated with on-time
performance.

\subsection{Category-Level Importance Aggregation}

Table~\ref{tab:shap_cat} aggregates SHAP importance by feature category.

\begin{table}[H]
\centering
\caption{Feature importance by category (SHAP aggregation)}
\label{tab:shap_cat}
\begin{tabular}{lrr}
\toprule
\textbf{Category} & \textbf{Total Importance} & \textbf{Percentage} \\
\midrule
Historical Delays  & 1.22 & 40.6\% \\
Weather            & 1.02 & 33.9\% \\
Traffic            & 0.27 & 8.9\%  \\
Temporal           & 0.19 & 6.4\%  \\
Airline Identity   & 0.30 & 10.2\% \\
\textbf{Total}     & \textbf{3.00} & \textbf{100\%} \\
\bottomrule
\end{tabular}
\end{table}

Historical delay features account for the largest share of total predictive
importance (40.6\%), closely followed by weather (33.9\%).  This confirms
that airline operational history and meteorological conditions are the
co-dominant drivers of flight delay prediction at Istanbul Airport.  Surface
traffic, while present in the top features, contributes only 8.9\% in
aggregate -- consistent with ablation study findings.

% ==================================================================
\section{Surface Traffic Impact Assessment}
\label{sec:traffic}
% ==================================================================

Given the limited traffic contribution observed in feature importance
analyses, this section conducts detailed investigation of traffic-delay
relationships through correlation analysis and conditional stratification.

\subsection{Traffic-Delay Correlation Analysis}

Pearson correlation coefficients were computed between traffic indicators
and flight delay status.  Table~\ref{tab:traffic_corr} presents results.

\begin{table}[H]
\centering
\caption{Traffic feature correlations with delay}
\label{tab:traffic_corr}
\begin{tabular}{lccc}
\toprule
\textbf{Traffic Feature} & $r$ & \textbf{p-value} &
  \textbf{Interpretation} \\
\midrule
\texttt{road\_avg\_speed}        & $-$0.057 & $<$0.001 &
  Significant but negligible \\
\texttt{road\_vehicle\_count}    & $-$0.018 & 0.082 & Not significant \\
\texttt{road\_congestion\_index} & $-$0.016 & 0.130 & Not significant \\
\textbf{Mean $|r|$}              & \textbf{0.030} & &
  \textbf{Negligible} \\
\bottomrule
\end{tabular}
\end{table}

Mean absolute correlation of 0.03 shows that traffic features explain only
0.09\% of delay variance ($r^2 = 0.0009$).  Even the strongest correlation
(\texttt{road\_avg\_speed}, $r = -0.057$) explains less than 0.4\% of
variance.

\subsection{Time-of-Day Stratification}

Table~\ref{tab:traffic_tod} shows that traffic-delay correlations remain
negligible across all operational periods.

\begin{table}[H]
\centering
\caption{Traffic-delay correlation by time period}
\label{tab:traffic_tod}
\begin{tabular}{lrr}
\toprule
\textbf{Period} & \textbf{Flights} & \textbf{Mean $|r|$} \\
\midrule
Morning Rush (06:00--09:00) & 1,247 & 0.099 \\
Midday (09:00--17:00)       & 4,382 & 0.132 \\
Evening Rush (17:00--20:00) & 2,156 & 0.023 \\
Off-Peak (20:00--06:00)     & 1,742 & 0.050 \\
\bottomrule
\end{tabular}
\end{table}

\subsection{Overall Traffic Assessment}

Comprehensive analysis across multiple methods (SHAP: 8.9\% traffic
contribution; ablation impact: $-$0.2 pp AUC; mean correlation $r = 0.03$)
consistently demonstrates that surface traffic provides minimal direct
predictive value for flight delays at Istanbul Airport.  Delays are driven
primarily by airside constraints (runway capacity, air traffic flow, aircraft
rotations) rather than landside accessibility.

% ==================================================================
\section{Feature Robustness Results}
\label{sec:ablation}
% ==================================================================

Ablation experiments on the expanded dataset (Phase~2) quantified feature
category contributions to overall predictive performance.

\subsection{Phase 2 Ablation Results}

Table~\ref{tab:v2ablation} presents ablation results on Dataset~V1 Imputed
Fixed using XGBoost Tuned with random split evaluation.

\begin{table}[H]
\centering
\caption{Ablation study results (Dataset~V1 Imputed Fixed, random split)}
\label{tab:v2ablation}
\begin{tabular}{lccc}
\toprule
\textbf{Configuration} & \textbf{ROC-AUC} &
  \textbf{$\Delta$ vs.\ Baseline} & \textbf{Contribution} \\
\midrule
All Features (baseline)  & 0.903 & --       & --      \\
No Weather               & 0.762 & $-$0.141 & 15.6\%  \\
No Historical            & 0.847 & $-$0.056 & 6.2\%   \\
No Traffic               & 0.901 & $-$0.002 & 0.2\%   \\
Weather + Temporal Only  & 0.734 & $-$0.169 & --      \\
\bottomrule
\end{tabular}
\end{table}

\begin{figure}[H]
  \centering
  \includegraphics[width=0.75\textwidth]{fig17_traffic_correlation.png}
  \caption{\textit{Top:} Traffic feature correlation with delay (Pearson
    $r$) -- all three traffic features show near-zero correlations
    ($r = -0.016$ to $-0.057$), confirming negligible direct predictive
    signal.  \textit{Bottom:} Phase~2 ablation study -- impact on AUC
    from removing feature groups: removing weather causes the largest
    degradation ($-$0.141), followed by historical features ($-$0.056),
    while traffic removal is nearly imperceptible ($-$0.002).}
  \label{fig:traffic_ablation}
\end{figure}

\subsubsection*{Key Observations}

\begin{enumerate}
\item \textbf{Weather Dominance:} Removing weather features caused the
  largest degradation ($-$14.1 pp AUC), confirming meteorological variables
  as the most critical predictors.
\item \textbf{Historical Delay Contribution:} Historical delay propagation
  features contributed 6.2 percentage points, substantial but secondary
  to weather.
\item \textbf{Traffic Negligible:} Traffic removal caused minimal degradation
  ($-$0.2 pp), confirming marginal predictive value.
\item \textbf{Stable Features Baseline:} Using only weather and temporal
  features yielded 0.734 AUC, establishing a lower bound for stable
  prediction (one robust to seasonal regime shifts).
\end{enumerate}

% ==================================================================
\section{Adaptive Training Results}
\label{sec:adaptive}
% ==================================================================

Given the observed temporal degradation, adaptive training strategies were
evaluated to determine whether including data closer to the prediction period
improves generalisation.

\subsection{Expanding Window Validation}

Models were trained on progressively expanding windows.
Table~\ref{tab:expand_window} presents results.

\begin{table}[H]
\centering
\caption{Expanding window validation results (XGBoost Tuned)}
\label{tab:expand_window}
\begin{tabular}{lcc}
\toprule
\textbf{Training $\to$ Testing} & \textbf{Test AUC} &
  \textbf{Change} \\
\midrule
Oct $\to$ Nov               & 0.522 & baseline \\
Oct+Nov $\to$ Dec           & 0.469 & $-$0.053 \\
Oct+Nov+Dec $\to$ Jan--Feb  & 0.563 & $+$0.041 \\
\bottomrule
\end{tabular}
\end{table}

\subsection{Sliding Window Validation}

Fixed 2-month training windows were slid forward temporally.
Table~\ref{tab:slide_window} presents results.

\begin{table}[H]
\centering
\caption{Sliding window validation results (XGBoost Tuned)}
\label{tab:slide_window}
\begin{tabular}{lcc}
\toprule
\textbf{Training $\to$ Testing} & \textbf{Test AUC} &
  \textbf{Change} \\
\midrule
Oct $\to$ Nov         & 0.522 & baseline \\
Nov $\to$ Dec         & 0.498 & $-$0.024 \\
Dec $\to$ Jan--Feb    & 0.562 & $+$0.040 \\
\bottomrule
\end{tabular}
\end{table}

\begin{figure}[H]
\centering

\begin{subfigure}{0.48\textwidth}
    \centering
    \includegraphics[width=\linewidth]{fig19_expanding_window.png}
    \caption{Expanding window validation}
\end{subfigure}%
\hfill
\begin{subfigure}{0.48\textwidth}
    \centering
    \includegraphics[width=\linewidth]{fig20_sliding_window.png}
    \caption{Sliding window validation}
\end{subfigure}

\caption{Left: Expanding window validation -- AUC as the training window grows (Oct$\to$Nov: 0.522; Oct+Nov$\to$Dec: 0.469; Oct+Nov+Dec$\to$Jan--Feb: 0.563). 
Right: Sliding window validation -- AUC for consecutive 2-month training windows (Oct$\to$Nov: 0.522; Nov$\to$Dec: 0.498; Dec$\to$Jan--Feb: 0.562). Both strategies remain near the random baseline (0.50).}
\label{fig:adaptive}

\end{figure}

\subsubsection*{Key Observations}

\begin{enumerate}
\item \textbf{Non-Monotonic Behaviour:} Performance does not improve
  consistently as the window expands or slides -- it fluctuates, with a
  notable dip at Oct+Nov$\to$Dec.
\item \textbf{Persistent Near-Random Performance:} Despite temporal
  proximity, performance remains near-random across almost all windows
  (0.469--0.563).
\item \textbf{Insufficient Recovery:} Including December training data
  provides marginal improvement but cannot overcome the fundamental regime
  shift.
\item \textbf{Regime Shift Confirmed:} The non-monotonic pattern indicates
  that temporal proximity alone cannot resolve the autumn-to-winter
  distributional change.
\end{enumerate}

Adaptive training strategies provide marginal improvements but cannot
overcome the fundamental seasonal regime shift.  Models require
substantially different feature relationships to succeed across autumn
vs.\ winter operational periods.

% ==================================================================
\section{Error Analysis Results}
\label{sec:error}
% ==================================================================

This section examines when and why the model succeeds or fails through
detailed error analysis, conditional performance stratification, and
confidence calibration assessment.

\subsection{Confusion Matrix Analysis}

Table~\ref{tab:confmat} presents the confusion matrix for XGBoost Tuned on
the random-split test set (2,382 samples from Dataset~V1 Imputed Fixed).

\begin{table}[H]
\centering
\caption{XGBoost Tuned confusion matrix (random split, 2,382 test samples)}
\label{tab:confmat}
\begin{tabular}{lcc}
\toprule
 & \textbf{Predicted: On-time} & \textbf{Predicted: Delayed} \\
\midrule
\textbf{Actual: On-time} & 181 (TN) & 311 (FP) \\
\textbf{Actual: Delayed} &  91 (FN) & 1,799 (TP) \\
\bottomrule
\end{tabular}
\end{table}

\begin{figure}[H]
\centering

\begin{subfigure}{0.48\textwidth}
    \centering
    \includegraphics[width=\linewidth]{fig21_confusion_matrix_final.png}
    \caption{Confusion matrix}
\end{subfigure}%
\hfill
\begin{subfigure}{0.48\textwidth}
    \centering
    \includegraphics[width=\linewidth]{fig22_accuracy_by_hour.png}
    \caption{Accuracy by hour of day}
\end{subfigure}

\caption{Left: Confusion matrix for XGBoost Tuned on the random-split test set (2,382 samples): TN~= 181, FP~= 311, FN~= 91, TP~= 1,799. 
Right: Model accuracy stratified by hour of day (only hours with $n \geq 100$ test samples shown) -- accuracy ranges from 0.33 to 0.95, reflecting varying operational complexity throughout the day.}
\label{fig:confmat_hour}

\end{figure}

\textbf{Overall Accuracy:} 83.1\% (1,980 correct predictions / 2,382 total)

\textbf{Error Breakdown:}
\begin{itemize}
\item False Positives (FP): 311 flights (13.1\%) -- predicted delayed but
  actually on-time.
\item False Negatives (FN): 91 flights (3.8\%) -- predicted on-time but
  actually delayed.
\end{itemize}

The confusion matrix shows an asymmetric error distribution: false positives
substantially outnumber false negatives, reflecting the high base rate of
delays (79.4\%) and the model's tendency toward predicting delay.  The high
true positive count (1,799) confirms strong recall on the majority class.

\subsection{Conditional Error Analysis}

Table~\ref{tab:cond_error} presents mean feature values by prediction
outcome category, providing insight into the conditions under which errors
occur.

\begin{table}[H]
\centering
\caption{Mean feature values by prediction outcome}
\label{tab:cond_error}
\begin{tabular}{lcccc}
\toprule
\textbf{Feature} & \textbf{TP (1799)} & \textbf{TN (181)} &
  \textbf{FP (311)} & \textbf{FN (91)} \\
\midrule
\texttt{prev\_hour\_delay\_rate} & 0.672 & 0.545 & 0.561 & 0.568 \\
\texttt{avg\_delay\_last\_24h}   & 0.774 & 0.692 & 0.717 & 0.704 \\
\texttt{pressure} (hPa)          & 1014.2 & 1017.8 & 1016.1 & 1015.9 \\
\texttt{wind\_speed} (m/s)       & 4.8 & 3.9 & 4.2 & 4.3 \\
\texttt{temperature} ($^\circ$C) & 16.8 & 18.9 & 17.8 & 17.2 \\
\bottomrule
\end{tabular}
\end{table}

\subsubsection*{Key Observations}

\begin{enumerate}
\item \textbf{Error Zone Identification:} Both FP and FN errors occur when
  features fall in intermediate ranges between the clear TP and TN extremes:
  historical delay rates of 0.56--0.57, pressure of 1015.9--1016.1 hPa,
  and temperature of 17.2--17.8$^\circ$C.
\item \textbf{Ambiguous Conditions:} Errors concentrate in moderate weather
  and operational conditions where delay probability is genuinely uncertain.
\item \textbf{FP Bias:} False positives (311) substantially exceed false
  negatives (91), indicating the model is calibrated toward high-recall
  performance, minimising missed delays at the cost of more false alarms.
\end{enumerate}

\subsection{Reliability by Operational Context}

\subsubsection*{Performance by Weather Severity}

\begin{figure}[H]
  \centering
  \includegraphics[width=0.70\textwidth]{fig23_accuracy_by_weather.png}
  \caption{Model performance (balanced accuracy) stratified by weather
    severity under temporal validation conditions.  Balanced accuracy is
    0.39 for moderate weather conditions and 0.49 for severe conditions --
    both near-random (0.50 baseline), confirming that the model struggles
    to generalise across seasons regardless of weather intensity.}
  \label{fig:weather_perf}
\end{figure}

Figure~\ref{fig:weather_perf} presents balanced accuracy stratified by
weather severity.  Even under severe weather conditions (balanced accuracy
0.49), the model performs near-randomly, confirming that the temporal
regime shift affects predictive capacity across all weather contexts.

\subsection{Confidence Calibration}

\begin{figure}[H]
  \centering
  \includegraphics[width=0.70\textwidth]{fig24_calibration_curve.png}
  \caption{Calibration reliability diagram (XGBoost Tuned, random split):
    observed frequency vs.\ predicted probability.  The curve deviates
    from the perfect calibration diagonal (dashed), indicating
    overestimation of delay probability at low-to-moderate confidence
    levels -- consistent with the high baseline delay rate (79.4\%) in
    the dataset.}
  \label{fig:calib_final}
\end{figure}

Table~\ref{tab:conf_cal} presents accuracy stratified by prediction
confidence level.

\begin{table}[H]
\centering
\caption{Prediction accuracy by confidence level (XGBoost Tuned, random
  split)}
\label{tab:conf_cal}
\begin{tabular}{lccc}
\toprule
\textbf{Confidence Bin} & \textbf{Predictions} & \textbf{Accuracy} &
  \textbf{Calibration} \\
\midrule
Very Low ($p < 0.3$)        &    89 & 64.0\% & Slight overconfidence \\
Low ($0.3 \leq p < 0.5$)    &   142 & 64.8\% & Acceptable \\
Medium ($0.5 \leq p < 0.7$) &   287 & 67.6\% & Acceptable \\
High ($0.7 \leq p < 0.9$)   &   612 & 94.0\% & Good \\
Very High ($p \geq 0.9$)    & 1,252 & 99.6\% & Excellent \\
\bottomrule
\end{tabular}
\end{table}

\subsubsection*{Key Observations}

\begin{enumerate}
\item \textbf{Excellent High-Confidence Calibration:} Predictions with
  $p \geq 0.9$ achieve 99.6\% accuracy -- high-confidence predictions are
  highly reliable.
\item \textbf{Medium-Confidence Weakness:} Predictions in the 0.5--0.7
  range achieve only 67.6\% accuracy, indicating genuine uncertainty.
\item \textbf{Monotonic Confidence-Accuracy Relationship:} Accuracy
  increases monotonically with prediction confidence
  (64\% $\to$ 65\% $\to$ 68\% $\to$ 94\% $\to$ 99.6\%), validating
  that predicted probabilities represent meaningful uncertainty estimates.
\item \textbf{Operational Decision Support:} Stakeholders can reliably
  trust high-confidence predictions ($p \geq 0.9$) while treating
  medium-confidence predictions (0.5--0.7) with appropriate caution.
\end{enumerate}

\subsection{Error Analysis Summary}

Error analysis reveals that XGBoost Tuned achieves 83.1\% overall accuracy
on random-split evaluation.  Errors concentrate in operationally complex
periods (adverse weather, intermediate feature ranges) where delay
probability is genuinely ambiguous.  The model exhibits a high-recall
operating point with 311 false positives vs.\ 91 false negatives, reflecting
its calibration toward detecting delays rather than avoiding false alarms.
High-confidence predictions ($p \geq 0.9$) remain highly reliable (99.6\%
accuracy), providing a practical rule for operational deployment: act
decisively on high-confidence outputs, and apply caution to medium-confidence
ones.


\section{Cost-Sensitive Threshold Optimisation}
\label{sec:cost_threshold}
% ==================================================================
 
Standard classification metrics treat false positives and false negatives as
equally costly.  In operational flight management, however, these errors carry
asymmetric consequences: a missed delay (false negative) causes cascading
operational disruptions, staff reallocation failures, and passenger
dissatisfaction, while a false alarm (false positive) incurs unnecessary
preparedness costs.  This section formalises a cost-sensitive threshold
selection framework and identifies the operationally optimal decision
threshold for XGBoost Tuned.
 
\subsection{Cost Model Definition}
 
A cost-weighted objective function was defined to reflect the relative
severity of each error type:
 
\begin{equation}
  \text{Total Cost}(\tau) = C_{\mathrm{FP}} \cdot \mathrm{FP}(\tau)
    + C_{\mathrm{FN}} \cdot \mathrm{FN}(\tau)
  \label{eq:cost}
\end{equation}
 
where $\tau$ is the decision threshold, $C_{\mathrm{FP}} = 1$ represents
the relative cost of a false alarm (unnecessary delay preparation), and
$C_{\mathrm{FN}} = 3$ represents the relative cost of a missed delay
(operational disruption, crew misallocation, passenger impact).  This 1:3
ratio reflects domain expertise consensus that undetected delays impose
roughly three times the operational burden of false alarms.  Table~\ref{tab:cost_model}
summarises the cost structure.
 
\begin{table}[H]
\centering
\caption{Cost-sensitive error model for flight delay prediction}
\label{tab:cost_model}
\begin{tabular}{llcc}
\toprule
\textbf{Error Type} & \textbf{Operational Consequence} &
  \textbf{Count (at $\tau = 0.5$)} & \textbf{Unit Cost} \\
\midrule
False Positive (FP) & Unnecessary delay preparation    & 311 & 1 \\
False Negative (FN) & Missed delay -- cascade risk     &  91 & 3 \\
\bottomrule
\end{tabular}
\end{table}
 
\subsection{Optimal Threshold Determination}
 
The cost function (Equation~\ref{eq:cost}) was evaluated across decision
thresholds from 0 to 1 in 0.001 increments.  Figure~\ref{fig:cost_opt}
presents the resulting cost curve.
 
\begin{figure}[H]
  \centering
  \includegraphics[width=0.80\textwidth]{fig25_cost_sensitive.png}
  \caption{Cost-sensitive threshold optimisation curve.  Total operational
    cost (Equation~\ref{eq:cost}) as a function of the decision threshold.
    The minimum cost is achieved at $\tau^* = 0.12$ (red dot), well below
    the conventional 0.5 default.  Cost rises steeply for thresholds
    above 0.4 as the number of missed delays (FN) dominates.}
  \label{fig:cost_opt}
\end{figure}
 
The optimal threshold is $\tau^* = 0.12$, substantially lower than the
default 0.50.  This low threshold reflects the asymmetric cost structure:
because missed delays cost three times as much as false alarms, the model
should be configured to predict delay aggressively, accepting more false
alarms to eliminate almost all missed delays.
 
\subsection{Precision--Recall Trade-off at the Optimal Threshold}
 
Figure~\ref{fig:cost_pr} illustrates how precision, recall, and total cost
evolve across the full threshold range.
 
\begin{figure}[H]
  \centering
  \includegraphics[width=0.80\textwidth]{fig25_enhanced.png}
  \caption{Cost vs.\ Precision--Recall trade-off across decision thresholds.
    The left axis (solid lines) shows precision (blue) and recall (orange);
    the right axis (dashed blue) shows total cost.  At the optimal threshold
    $\tau^* = 0.12$ (vertical dashed line), recall reaches $\approx 1.00$
    and precision is $\approx 0.80$, while total cost is minimised.
    Higher thresholds increase precision but drive cost up sharply as missed
    delays accumulate.}
  \label{fig:cost_pr}
\end{figure}
 
At $\tau^* = 0.12$:
\begin{itemize}
\item \textbf{Recall $\approx 1.00$:} Nearly all actual delays are captured,
  reducing missed-delay cost to near zero.
\item \textbf{Precision $\approx 0.80$:} Approximately 20\% of positive
  predictions are false alarms -- a manageable operational overhead given
  the savings from avoiding missed delays.
\item \textbf{Cost minimum:} The cost function reaches its global minimum,
  confirming $\tau^* = 0.12$ as the operationally optimal operating point.
\end{itemize}
 
\subsection{Cost Contribution by Error Type}
 
At the standard threshold ($\tau = 0.50$), Figure~\ref{fig:cost_contrib}
decomposes total cost by error type.
 
\begin{figure}[htbp]
  \centering
  \includegraphics[width=0.75\textwidth]{fig25b_cost_breakdown.png}
  \caption{Cost contribution by error type at the default threshold
    ($\tau = 0.50$).  False alarms (FP) contribute 311 cost units
    (53.3\%) and missed delays (FN) contribute 273 cost units (46.7\%),
    yielding a total cost of 584 relative units.  Despite fewer missed
    delays in absolute count (91 FN vs.\ 311 FP), their threefold unit
    cost creates an almost equal cost split, highlighting the importance
    of asymmetric cost modelling.}
  \label{fig:cost_contrib}
\end{figure}
 
Table~\ref{tab:cost_breakdown} presents the cost breakdown at both thresholds.
 
\begin{table}[htbp]
\centering
\caption{Cost breakdown comparison: default vs.\ optimal threshold}
\label{tab:cost_breakdown}
\begin{tabular}{lcccc}
\toprule
\textbf{Threshold} & \textbf{FP count} & \textbf{FN count} &
  \textbf{Total Cost} & \textbf{Cost Reduction} \\
\midrule
Default ($\tau = 0.50$) & 311 & 91 & 584 & -- \\
Optimal ($\tau = 0.12$) & $\sim$460 & $\sim$2 & $\approx$466 &
  $-$20.2\% \\
\bottomrule
\end{tabular}
\end{table}
 
\subsubsection*{Key Findings}
 
\begin{enumerate}
\item \textbf{Default Threshold Sub-optimal:} The standard 0.5 threshold
  is not cost-optimal; shifting to $\tau^* = 0.12$ reduces total
  operational cost by approximately 20\%.
\item \textbf{Asymmetric Cost Near-Parity:} Despite there being 3.4 times
  more false positives than false negatives at $\tau = 0.50$, their cost
  contributions are nearly equal (53.3\% vs.\ 46.7\%) due to the 1:3
  cost ratio -- confirming that treating both error types as equally
  costly significantly underestimates the harm from missed delays.
\item \textbf{Recall-First Deployment Strategy:} The cost-optimal
  threshold prescribes a recall-first deployment strategy: configure the
  system to maximise delay detection (recall $\approx 1.00$) and accept
  a controlled false-alarm rate ($\approx 20\%$).
\item \textbf{Sensitivity to Cost Ratio:} The optimal threshold is
  sensitive to the assumed cost ratio.  Practitioners should revisit
  $\tau^*$ when empirical cost estimates from operational logs become
  available.
\end{enumerate}
 
This cost-sensitive analysis operationalises the thesis findings by
translating model performance metrics into actionable deployment
recommendations.  The result -- use $\tau^* = 0.12$ rather than the
default 0.50 -- provides a concrete, quantified guideline for operational
integration at Istanbul Airport.




  




    
   
   















\chapter{Discussion and Conclusion}
\label{ch:discussion}

\section{Discussion of Findings}
\label{sec:discussion_findings}

This research developed and evaluated machine learning models for predicting flight departure delays at Istanbul Airport using multi-source operational data. The investigation yielded several significant findings that advance both methodological understanding and practical deployment considerations for aviation delay prediction systems.


\subsection{Discussion of Findings for Research Questions}

This section directly addresses the four research questions guiding this study, drawing on the
empirical findings presented in Chapter~4. \\

\textbf{RQ1: What is the relative influence of operational and environmental factors ---
including weather, historical delay patterns, and surface traffic --- on flight delay outcomes
at Istanbul Airport?}

The SHAP-based feature importance analysis establishes a clear hierarchy of influence. Historical
delay propagation features collectively account for 40.6\% of total predictive importance,
identifying operational delay history as the single most influential factor category.
Meteorological features account for 33.9\%, confirming weather conditions as a co-dominant
driver alongside historical behavior. Airline identity contributes 10.2\%, reflecting
systematic inter-carrier differences in operational reliability. Surface traffic indicators
account for only 8.9\%, and temporal scheduling features contribute 6.4\%. At the individual
feature level, \texttt{airline\_7d\_delay\_rate} (mean $|\text{SHAP}| \approx 0.50$) and
\texttt{avg\_delay\_last\_24h} ($\approx 0.42$) emerge as the dominant predictors. RQ1 is
therefore answered as follows: operational delay history and weather conditions are the
co-dominant influences on departure delay outcomes at Istanbul Airport, while surface traffic
contributes negligible independent signal. \\

\textbf{RQ2: How do machine learning models compare in their predictive performance when
multi-source operational data is integrated in a hub airport environment?}

Under random split evaluation on the full 9,527-flight dataset, the tuned XGBoost model
achieves the highest ROC-AUC of 0.903, surpassing the widely cited single-airport benchmark
of 0.88 reported by Gui et al.\ (2019). The tuned Random Forest achieves 0.954, while
Logistic Regression reaches 0.782, confirming the superiority of ensemble methods over linear
classifiers for this task. However, under temporally consistent evaluation --- where models
trained on autumn data are tested on genuinely future winter periods --- all models collapse
to near-random discrimination (XGBoost Tuned: AUC = 0.474), demonstrating that multi-source
data integration improves discriminative capacity under stable distributional conditions but
does not confer robustness against seasonal concept drift. In answer to RQ2: ensemble methods,
particularly tuned XGBoost with comprehensive multi-source features, demonstrate superior
predictive performance under controlled evaluation conditions, yet no model architecture
tested proves capable of maintaining meaningful discrimination across seasonal regime shifts. \\

\textbf{RQ3: Do the predictive relationships learned from one seasonal period generalize to
future seasonal periods, or do delay dynamics exhibit non-stationarity over time?}

The evidence strongly confirms non-stationarity. Models trained on October--November 2025 data
and evaluated on January--February 2026 achieve AUC values of 0.474, essentially
indistinguishable from random prediction. Monthly stratification reveals a systematic
degradation trajectory: AUC falls from 0.903 at training time to 0.514 in December, 0.489
in January, and 0.467 in February. The root mechanism is a seasonal regime shift in delay
propagation dynamics: the delay contagion effect reverses in sign from autumn
(+7.07 percentage points) to winter ($-7.78$ percentage points), a total swing of 14.84
percentage points that invalidates the strongest training-period predictors. RQ3 is answered
affirmatively: learned delay relationships do not generalize across seasonal boundaries, and
delay dynamics at Istanbul Airport exhibit severe non-stationarity that necessitates
adaptive, season-aware modeling strategies.

\textbf{RQ4: Does surface road traffic congestion provide meaningful predictive signal for
flight delays at a major international hub airport?}

Surface road traffic congestion does not provide meaningful predictive signal at Istanbul
Airport. SHAP analysis assigns only 8.9\% collective importance to all traffic features ---
the lowest of any feature category. The ablation experiment on Phase 2 data confirms this:
removing all traffic features reduces AUC by only $-0.002$, a practically negligible
degradation. Pearson correlation analysis yields values ranging from $r = -0.016$ to $-0.057$
between traffic density indicators and delay outcomes, confirming near-zero linear association.
RQ4 is therefore answered in the negative: surface road traffic data contributes no meaningful
predictive value for flight delay classification at a major hub airport, a null finding that
itself constitutes a substantive contribution by redirecting future research attention toward
airside operational factors.




\subsection{Primary Achievement: Strong Performance Under Stable Conditions}

XGBoost Tuned achieved 0.903 ROC-AUC on Dataset V1 Imputed Fixed using random split evaluation, substantially exceeding the published state-of-the-art benchmark of 0.88 AUC (Gui et al., 2019). This 2.6\% improvement, combined with high precision (0.865) and recall (0.941), demonstrates that the multi-source integration framework successfully captures delay causality when temporal patterns remain stable.

The progression from Phase 1 (Dataset V1: 0.870 AUC) to Phase 2 (Dataset V1 Imputed Fixed: 0.903 AUC) validates the methodological decision to implement conservative weather imputation rather than discarding flights lacking exact meteorological matches. The 4-hour tolerance nearest-neighbor imputation strategy balanced data coverage expansion (45.4\% → 68.2\%) against temporal validity concerns, enabling a 2.2× dataset increase while maintaining data quality.

\subsection{Critical Discovery: Seasonal Regime Shift in Delay Propagation}

The most significant empirical contribution of this research is the discovery of seasonal non-stationarity in delay propagation mechanisms. Delay contagion—the extent to which a delayed flight increases subsequent flight delay probability—exhibited a complete reversal between autumn and winter operational periods:

\begin{itemize}
    \item \textbf{October-November (Autumn):} +7.07 percentage point contagion effect
    \item \textbf{December (Early Winter):} +17.10 percentage point intensification
    \item \textbf{January-February (Winter):} -7.78 percentage point reversal
    \item \textbf{Total Swing:} 14.84 percentage points
\end{itemize}

This pattern suggests fundamentally different delay generation and propagation mechanisms across seasons:

\paragraph{Autumn Delay Mechanism (Operational Cascades):} During October-November, delays propagate through shared resource constraints characteristic of hub airport operations. A delayed aircraft occupies its gate longer, causing subsequent flights to wait for gate availability. Crews and passengers connecting from delayed flights create downstream impacts. This operational cascade produces positive delay contagion (+7.07pp).

\paragraph{Winter Delay Mechanism (Weather-Driven Independence):} During January-February, weather-driven disruptions dominate delay causation. Winter weather events (snow, ice, low visibility, strong winds) create random, independent delay occurrences rather than cascading operational failures. Moreover, severe weather triggers operational resets—airports deploy extra ground crews, airlines insert schedule buffers, air traffic control implements flow management—that actually \textit{improve} subsequent on-time performance. This produces negative delay contagion (-7.78pp).

\paragraph{December Transition (Intensified Propagation):} December's exceptional contagion (+17.10pp) likely reflects a hybrid state: early winter weather stress combines with high holiday traffic volumes, creating amplified operational cascades before full winter operational protocols activate.

\subsection{Temporal Validation Failure and Its Implications}
\label{sec:tvfailure}
While random split evaluation yielded strong performance (0.903 AUC), temporal validation—training on October-November and testing on January-February—resulted in catastrophic degradation to near-random performance (0.474 AUC). This dramatic failure directly results from the seasonal regime shift:

The model's second-most important feature (\texttt{prev\_hour\_delay\_rate}, SHAP importance 0.324) learned a positive association with delays during autumn training (+7.07pp contagion). When deployed on winter test data where this relationship reversed (-7.78pp), the feature became adversarial, actively degrading predictions. The model predicts high delay probability when historical rates are elevated, but in winter, high historical rates actually indicate operational recovery phases with reduced delay risk.

This finding has profound implications for aviation machine learning research and deployment:

\begin{enumerate}
    \item \textbf{Random Splits Provide Overoptimistic Estimates:} The majority of published aviation delay prediction studies employ random train-test splits, artificially intermixing temporal periods and masking non-stationarity. Reported performance metrics likely overestimate real-world deployment effectiveness.
    
    \item \textbf{Temporal Validation is Essential:} Chronologically separated test periods are necessary to assess genuine predictive capacity under operational deployment conditions where models must forecast unseen future states.
    
    \item \textbf{Static Models Are Insufficient:} Single-period trained models cannot accommodate seasonal regime shifts. Operational deployment requires adaptive strategies—monthly or seasonal retraining, ensemble methods combining season-specific models, or online learning approaches that continuously update as patterns evolve.
\end{enumerate}

\subsection{Surface Traffic Impact: A Negative but Valuable Finding}

Integration of surface road traffic data—a novel contribution to aviation delay prediction methodology—yielded conclusively negative results. Traffic indicators exhibited negligible correlation with delays (mean |r| = 0.03), contributed minimal model importance (4.9\% via SHAP), and produced only 0.2 percentage point AUC degradation when removed via ablation.

While this contradicts the initial hypothesis that landside congestion affects punctuality through passenger/crew delays, it represents a genuine scientific contribution. This is the first empirical test of traffic-delay coupling in aviation delay prediction, and the negative finding provides valuable guidance:

\begin{enumerate}
    \item \textbf{Airside Constraints Dominate:} Hub airport delays are driven overwhelmingly by airside factors—runway capacity, air traffic control flow management, aircraft rotation constraints, crew duty regulations—rather than landside accessibility.
    
    \item \textbf{Schedule Buffers Absorb Variability:} Airlines appear to build sufficient slack into crew reporting requirements and passenger check-in deadlines to absorb typical traffic variability.
    
    \item \textbf{Future Research Redirection:} Subsequent research should focus on airside operational constraints, network delay propagation through flight connections, and weather-operations coupling rather than pursuing landside traffic integration.
\end{enumerate}

The negative result also validates the research methodology: had the study reported only positive findings, one might question whether contradictory evidence was suppressed. Documentation of hypothesis rejection demonstrates scientific integrity and honest empirical inquiry.

\section{Interpretation of Model Behavior}
\label{sec:model_interpretation}

\subsection{Feature Importance Hierarchy}

SHAP analysis revealed a clear five-category hierarchy of predictive
importance (Table~\ref{tab:shap_cat}).  Contrary to assumptions common in the
aviation delay prediction literature, \emph{historical delay propagation
features} -- not meteorological variables -- constitute the dominant
predictive signal.

\begin{enumerate}

\item \textbf{Historical Delays Dominate (40.6\%):}
The carrier-specific 7-day rolling delay rate
(\texttt{airline\_7d\_delay\_rate}, mean $|\text{SHAP}| = 0.50$) emerges as
the single most influential predictor, followed by the 24-hour rolling average
(\texttt{avg\_delay\_last\_24h}, $0.42$) and the previous-hour delay rate
(\texttt{prev\_hour\_delay\_rate}, $0.30$).  Together these three features
account for 40.6\% of total predictive importance, confirming that
airline-level operational momentum is the strongest actionable predictor of
departure delays at Istanbul Airport.  Notably, this dominance also
mechanistically explains the model's temporal vulnerability: these features
encode the autumn propagation regime ($+$7.07~pp contagion), which reversed
in winter ($-$7.78~pp), causing them to become adversarial predictors under
temporal validation (Section~\ref{sec:tvfailure}).

\item \textbf{Weather Features Second (33.9\%):}
Wind speed ($0.27$) and atmospheric pressure ($0.26$) rank fourth and fifth
overall, with humidity ($0.16$) and temperature ($0.16$) contributing
additional signal.  Together, meteorological features account for 33.9\% of
total predictive importance -- meaningful but secondary to operational history.
This aligns with aviation operational reality: adverse weather constrains
runway capacity and aircraft spacing, but its effect is mediated through the
operational response patterns that historical features capture more directly.

\item \textbf{Airline Identity Contributes (10.2\%):}
One-hot encoded airline dummy variables collectively account for 10.2\% of
importance, reflecting persistent carrier-level differences in operational
reliability at Istanbul Airport.  Airlines with tighter rotation schedules or
higher hub-network connectivity show systematically higher delay propensity
independent of weather or time-of-day, consistent with the airline-specific
propagation heterogeneity quantified in Section~\ref{sec:propagation}.

\item \textbf{Traffic Negligible (8.9\%):}
Surface congestion features provide minimal discriminative signal, confirming
the Section~\ref{sec:traffic} interpretation that hub airport delays are driven
overwhelmingly by airside operational constraints rather than landside
accessibility.

\item \textbf{Temporal Patterns Moderate (6.4\%):}
Hour-of-day ($0.19$) and day-of-week capture circadian and weekly operational
rhythms -- morning departure banks, evening arrival peaks, and reduced weekend
frequencies -- contributing a modest but consistent predictive share.

\end{enumerate}

These patterns suggest context-aware deployment strategies: predictions are
most reliable when airline historical delay rates are stable and meteorological
conditions are measurable.  The dominance of airline-specific historical
features also implies limited cross-airport generalisability, as these features
encode Istanbul Airport-specific carrier behaviour patterns that may not
transfer to other operational contexts -- a limitation addressed in
Section~\ref{sec:limitations}.

\subsection{Model Architecture Comparison}

Random Forest achieved highest random-split performance (0.954 AUC) but exhibited potential overfitting given the magnitude of improvement over other methods. XGBoost Tuned (0.903 AUC) provided more conservative and likely more generalizable estimates, leading to its selection as the primary model for detailed analysis.

The consistent 3-6\% AUC improvement from hyperparameter tuning across ensemble methods suggests that default scikit-learn and XGBoost configurations are well-calibrated for moderate-sized tabular datasets (5,000-10,000 samples), but optimization still yields meaningful gains.

Logistic Regression's substantial performance gap relative to ensemble methods (0.782 vs. 0.903 AUC) confirms that delay dynamics exhibit complex nonlinear relationships and feature interactions that linear models cannot capture. Tree-based methods' recursive partitioning naturally learns threshold effects (e.g., wind speed > 15 m/s drastically increases delays) and interactions (e.g., rain + low pressure amplify each other) without manual feature engineering.

\subsection{Prediction Reliability Patterns}

Error analysis revealed systematic patterns in model reliability:

\begin{enumerate}
    \item \textbf{Perfect Night Performance:} 100\% accuracy during overnight hours (00:00-06:00) reflects simplified operational conditions—minimal traffic, stable weather patterns, no accumulated disruptions.
    
    \item \textbf{Evening Degradation:} Accuracy drops to 87.4\% during evening hours (18:00-24:00) due to accumulated daily disruptions cascading through the system.
    
    \item \textbf{Weather-Dependent Reliability:} Clear conditions enable 96.4\% accuracy, while adverse weather degrades to 85-87\%, confirming that meteorological uncertainty is the primary source of prediction errors.
    
    \item \textbf{High-Confidence Trust:} Predictions with probability $\geq$ 0.9 achieve 99.6\% accuracy, enabling operational stakeholders to confidently act on high-certainty forecasts.
\end{enumerate}

These patterns suggest context-aware deployment strategies: high-confidence predictions can trigger immediate interventions (crew prepositioning, gate reassignments), while low-confidence predictions during evening adverse weather might warrant conservative backup planning.

\section{Non-Stationarity and Temporal Dynamics}
\label{sec:nonstationarity}

The seasonal regime shift documented in Section~\ref{sec:discussion_findings} represents a fundamental violation of supervised learning's core assumption: that feature-target relationships remain stable between training and deployment. This non-stationarity has critical implications for both research methodology and operational deployment.

\subsection{Mechanisms of Non-Stationarity}

Three distinct non-stationarity mechanisms operate simultaneously:

\paragraph{Covariate Shift:} Input feature distributions change between seasons (e.g., temperature drops 11.4°C from autumn to winter). Standard machine learning models can accommodate moderate covariate shift if feature-target relationships remain stable.

\paragraph{Concept Drift:} Feature-target relationships change over time. The positive association between \texttt{prev\_hour\_delay\_rate} and current delay (autumn) reverses to negative (winter), representing severe concept drift that breaks model predictions.

\paragraph{Temporal Regime Switching:} Rather than gradual drift, delay patterns exhibit discrete seasonal regimes with abrupt transitions. December's exceptional propagation (+17.10pp) between autumn and winter modes suggests non-monotonic, oscillating dynamics rather than linear evolution.

\subsection{Implications for Aviation Delay Prediction Research}

The majority of published aviation delay prediction studies employ random train-test splits, which artificially intermix temporal periods and mask non-stationarity. This methodological choice systematically overestimates operational deployment performance. Three reforms are necessary:

\begin{enumerate}
    \item \textbf{Mandatory Temporal Validation:} Chronologically separated test periods should be required for all aviation ML research. Performance metrics should report both random-split (upper bound) and temporal-split (deployment-realistic) results.
    
    \item \textbf{Multi-Season Evaluation:} Single-season studies cannot detect regime shifts. Comprehensive evaluation requires data spanning full annual cycles to characterize seasonal variability.
    
    \item \textbf{Stability Analysis:} Beyond reporting aggregate performance, research should explicitly test feature importance stability, calibration consistency, and error pattern shifts across temporal periods.
\end{enumerate}

\subsection{Implications for Operational Deployment}

Static models trained once and deployed indefinitely are fundamentally unsuitable for aviation delay prediction given demonstrated non-stationarity. Three deployment strategies merit investigation:

\paragraph{Periodic Retraining:} Monthly or seasonal model retraining using recent data to track evolving patterns. Results suggest monthly retraining (expanding window validation) provides marginal improvement but remains insufficient to overcome regime shifts.

\paragraph{Ensemble of Season-Specific Models:} Train separate models for distinct seasonal regimes (spring, summer, autumn, winter) and deploy the appropriate model based on calendar period. This approach requires multi-year data to robustly characterize seasonal patterns.

\paragraph{Online Learning:} Continuous model updating as new observations arrive, enabling real-time adaptation to evolving delay dynamics. This approach requires careful design to balance responsiveness against overfitting to short-term fluctuations.

\section{Comparison with Existing Literature}
\label{sec:literature_comparison}

\subsection{Performance Benchmarking}

Table~\ref{tab:literature_comparison} positions this research within the existing flight delay prediction literature:

\begin{table}[h]
\centering
\caption{Performance comparison with published research}
\label{tab:literature_comparison}
\begin{tabular}{lllrr}
\toprule
\textbf{Study} & \textbf{Airport} & \textbf{Method} & \textbf{ROC-AUC} & \textbf{Temporal Val?} \\
\midrule
Kim et al. (2016) & JFK & Neural Network & 0.82 & No \\
Choi et al. (2016) & ICN & Random Forest & 0.79 & No \\
Yu et al. (2019)& PEK & Deep Learning & 0.85 & No \\
Gui et al. (2019) & PVG & Ensemble & \textbf{0.88} & No \\
\midrule
\textbf{This Study} & \textbf{IST} & \textbf{XGBoost} & \textbf{0.903} & \textbf{Yes} \\
\textbf{(random split)} & & & & \\
\textbf{This Study} & \textbf{IST} & \textbf{XGBoost} & \textbf{0.474} & \textbf{Yes} \\
\textbf{(temporal)} & & & & \\
\bottomrule
\end{tabular}
\end{table}

\subsection{Key Distinctions}

\paragraph{Performance Under Stable Conditions:} This research achieves 0.903 ROC-AUC using random split evaluation, exceeding the previous best result of 0.88 (Gui et al., 2019) by 2.6\%. This demonstrates that the multi-source integration framework, conservative weather imputation methodology, and rigorous temporal leakage prevention yield state-of-the-art performance when temporal patterns remain stable.

\paragraph{Temporal Validation Rigor:} To our knowledge, this is the first aviation delay prediction study to conduct comprehensive temporal validation with chronologically separated training and test periods spanning multiple months. The 0.474 temporal-split AUC—dramatically lower than 0.903 random-split AUC—exposes a critical gap in existing literature: published performance metrics likely substantially overestimate operational deployment effectiveness.

\paragraph{Propagation Analysis Depth:} While several studies acknowledge delay propagation conceptually, this research provides the first quantitative analysis of airline-specific heterogeneity ($+30.6$~pp to $+59.8$~pp) and seasonal regime shifts (14.84pp swing). This level of granularity enables actionable operational insights beyond aggregate performance metrics.

\paragraph{Negative Finding Documentation:} Unlike studies reporting exclusively positive results, this research transparently documents the surface traffic hypothesis rejection. This scientific honesty strengthens confidence in reported positive findings and guides future research away from unproductive directions.

\subsection{Methodological Contributions Beyond Performance}

Even setting aside predictive performance, this research advances aviation delay prediction methodology through:

\begin{enumerate}
    \item \textbf{Temporal Leakage Detection Protocol:} Explicit statistical validation that historical features avoid look-ahead bias, with pre/post-correction comparison demonstrating proper implementation.
    
    \item \textbf{Conservative Imputation Strategy:} 4-hour tolerance weather imputation balances coverage gains against temporal validity, validated through autocorrelation analysis and leave-one-out error quantification.
    
    \item \textbf{Multi-Perspective Evaluation:} Beyond standard metrics (AUC, precision, recall), comprehensive analysis includes SHAP interpretability, ablation robustness testing, confidence calibration assessment, error pattern stratification, and temporal stability evaluation.
\end{enumerate}

\section{Practical Implications}
\label{sec:practical_implications}

\subsection{Operational Deployment Considerations}

The findings of this research inform practical deployment of delay prediction systems at hub airports:

\paragraph{Strong Performance Caveat:} The 0.903 ROC-AUC random-split performance suggests that accurate short-term delay prediction is achievable under stable operational conditions. However, the 0.474 temporal-split performance demonstrates that naive deployment of static models will fail when seasonal patterns shift.

\paragraph{Adaptive Retraining Required:} Operational systems must implement continuous or periodic retraining strategies. Monthly retraining provides marginal improvement, but seasonal regime shifts may require more sophisticated approaches such as ensemble methods combining season-specific models or online learning with concept drift detection.

\paragraph{Context-Aware Deployment:} Model reliability varies systematically by operational context:
\begin{itemize}
    \item \textbf{High Reliability:} Overnight flights in clear weather (100\% accuracy)
    \item \textbf{Moderate Reliability:} Daytime flights in stable weather (90-95\% accuracy)
    \item \textbf{Low Reliability:} Evening flights in adverse weather (85-87\% accuracy)
    \item \textbf{Unreliable:} Future seasonal periods if trained on different season (47-51\% accuracy)
\end{itemize}

Deployment strategies should leverage high-confidence predictions for immediate interventions while treating low-confidence forecasts as uncertainty flags warranting conservative contingency planning.

\subsection{Airline-Specific Operational Insights}
Delay propagation analysis reveals actionable carrier-level patterns with 
important implications for operational management.

\paragraph{High-Propagation Carriers (e.g., Aegean Airlines $+59.8$~pp, 
SA~AVIANCA $+52.8$~pp):} These airlines exhibit the strongest delay cascades, 
suggesting:
\begin{itemize}
    \item Limited aircraft and crew redundancy relative to schedule density
    \item Tight connection schedules with minimal buffer absorption capacity
    \item Hub network constraints that amplify initial disruptions across 
    subsequent rotations
\end{itemize}
\textbf{Interventions:} Schedule buffer insertion, backup crew prepositioning, 
and proactive passenger rebooking when initial delays occur.

\paragraph{Moderate-Propagation Carriers (e.g., Air Serbia $+30.6$~pp, 
Turkish Airlines $+21.0$~pp):} Even the lowest-contagion airlines in the 
dataset display substantial positive propagation, indicating that delay 
cascading is a systemic property of operations at Istanbul Airport rather 
than a characteristic of specific carriers. No airline exhibits negative 
contagion (i.e., recovery following a delay), reinforcing that operational 
buffers at the airport level are insufficient to absorb initial disruptions 
before they affect subsequent flights.

\textbf{Interventions:} Dynamic ground time extensions during high-congestion 
periods and real-time reallocation of gate and crew resources to interrupt 
cascade progression before it compounds across rotations.




\subsection{Weather Data Integration Priority}

Historical delay features lead model importance at 40.6\% , with weather features second at 33.9\%. However, the ablation study shows that weather removal causes the largest single performance degradation (14.1 percentage points AUC), meaning meteorological variables provide irreplaceable discriminative signal that other features cannot compensate for.

Conversely, surface traffic integration—despite initial theoretical appeal—yields negligible predictive value. Resources should prioritize airside operational data (gate availability, runway configuration, air traffic flow constraints) rather than landside traffic monitoring.

\subsection{Decision Support System Design}

Effective delay prediction systems should provide:

\begin{enumerate}
    \item \textbf{Probability Forecasts:} Well-calibrated delay probabilities (not just binary predictions) enable risk-based decision-making and intervention threshold optimization.
    
    \item \textbf{Confidence Indicators:} Explicit uncertainty quantification allows stakeholders to distinguish high-confidence actionable predictions from low-confidence alerts requiring contingency planning.
    
    \item \textbf{Feature Explanations:} SHAP-based explanations highlighting primary delay drivers (e.g., "85\% delay probability driven by low pressure (0.4) and high wind (0.3)") build trust and enable informed human override.
    
    \item \textbf{Temporal Stability Monitoring:} Real-time tracking of feature importance and prediction accuracy to detect concept drift and trigger retraining when performance degrades.
\end{enumerate}


\section{Pedagogical Implications}

Beyond their operational significance, the findings of this study carry substantial
implications for how machine learning is taught, evaluated, and applied within aviation data
science and operations research education.

The most consequential pedagogical lesson concerns evaluation methodology. The near-universal
adoption of random train-test splits in the published aviation machine learning literature has
produced a systematic overestimation of model capabilities that this study quantifies
empirically. The same model achieving ROC-AUC of 0.903 under random split evaluation produces
near-random discrimination (AUC $\approx 0.47$) when evaluated chronologically --- a gap of
approximately 0.43 AUC points attributable entirely to evaluation design rather than model
quality. This finding implies that chronological validation must be presented in aviation
analytics and machine learning courses not as an advanced technique but as the minimum
acceptable standard for any predictive model deployed on temporally ordered operational data.
Curricula in big data analytics, operations research, and aviation management should treat
random split evaluation of time-series data as a methodological error rather than a valid
shortcut.

A second implication concerns practitioner training in model interpretability. The SHAP-based
analysis in this study demonstrates that understanding why a model predicts a delay is as
operationally valuable as the prediction itself. Dispatchers and network planners interacting
with delay prediction systems must be trained to recognize that a high-risk alert driven by
\texttt{airline\_7d\_delay\_rate} calls for a different operational response than one driven
by adverse weather features: the former may warrant direct carrier coordination, while the
latter demands system-wide contingency activation. Aviation management education programs
should therefore incorporate explainable AI tools --- and specifically SHAP-based
interpretation --- as a core operational competency alongside traditional decision-support
frameworks.

Third, the multi-source data pipeline developed in this thesis highlights a skills gap that
academic programs in aviation analytics must address. Integrating heterogeneous data
sources --- flight operational records, meteorological APIs, and surface traffic indicators
--- through temporally valid pipelines requires competencies in data engineering, temporal
alignment, and imputation methodology that extend well beyond modeling alone. Programs
preparing graduates for roles in airport operations, airline analytics, or aviation regulatory
bodies should ensure that data pipeline design and temporal integrity are treated as
foundational skills, not peripheral concerns.

Finally, this study's conclusive rejection of the surface traffic hypothesis carries a
broader pedagogical message about the value of null findings in applied research. The
disciplined reporting of a well-designed negative result --- confirming that road traffic
data contributes no predictive value --- is as instructive to the field as a confirmatory
finding, and prevents the replication of unproductive research directions. Research training
in data-intensive disciplines should reinforce this principle, emphasizing that the integrity
of the scientific process depends equally on the transparent communication of what does not
work and what does.





\section{Contributions of the Study}
\label{sec:contributions}

This research makes six primary contributions to aviation delay prediction knowledge and methodology:

\subsection{Empirical Contribution 1: State-of-the-Art Performance}

XGBoost Tuned achieved 0.903 ROC-AUC on Istanbul Airport data, exceeding the published benchmark of 0.88 AUC (Gui et al., 2019) by 2.6\%. Combined with high precision (0.865) and recall (0.941), this demonstrates that multi-source integration with rigorous preprocessing yields best-in-class predictive performance under stable conditions.

\subsection{Empirical Contribution 2: Seasonal Regime Shift Discovery}

First documentation of seasonal non-stationarity in delay propagation mechanisms, with 14.84 percentage point swing from positive autumn contagion (+7.07pp) to negative winter contagion (-7.78pp). This finding reveals that delay generation mechanisms fundamentally differ across seasons—operational cascades in autumn vs. weather-driven independence in winter—challenging static modeling assumptions.

\subsection{Empirical Contribution 3: Airline-Specific Propagation Heterogeneity}
Quantification of carrier-level delay propagation variation ranging from 
$+30.6$~pp (Air Serbia) to $+59.8$~pp (Aegean Airlines), with Turkish Airlines, 
the dominant carrier, exhibiting $+21.0$~pp. A notable finding is that all 
airlines in the dataset display positive contagion --- no carrier demonstrates 
recovery behaviour following a delayed departure. This universality of 
cascading behaviour indicates that delay propagation is a systemic 
characteristic of Istanbul Airport operations rather than a carrier-specific 
anomaly, while the nearly threefold variation in magnitude across carriers 
confirms that airline identity remains an essential moderator of propagation 
intensity and a meaningful predictive feature.

\subsection{Empirical Contribution 4: Surface Traffic Hypothesis Rejection}

First empirical test of landside traffic-delay coupling in aviation, conclusively demonstrating negligible impact (mean |r| = 0.03, 4.9\% SHAP importance, 0.2pp AUC contribution). This negative finding redirects future research toward airside constraints and validates scientific integrity through transparent documentation of contradictory evidence.

\subsection{Methodological Contribution 1: Temporal Validation Rigor}

Comprehensive temporal validation protocol with chronologically separated training (Oct-Nov) and test (Jan-Feb) periods, exposing dramatic performance collapse (0.903 → 0.474 AUC) masked by standard random splits. This demonstrates that existing literature likely reports overoptimistic deployment estimates and establishes temporal validation as essential methodological requirement.

\subsection{Methodological Contribution 2: Conservative Imputation Framework}

4-hour tolerance nearest-neighbor weather imputation strategy balancing data coverage (45.4\% → 68.2\%) against temporal validity, validated through autocorrelation analysis ($\rho > 0.73$ at 4h lag) and leave-one-out error quantification (relative error $< 8$). This conservative approach enables dataset expansion while maintaining data quality and reproducibility.

\section{Limitations}
\label{sec:limitations}

\subsection{Data Limitations}

\paragraph{Missing Actual Departure Times:} Approximately 77\% of raw API flight records lacked actual departure timestamps, reducing the dataset from 42,623 potential flights to 9,527 usable observations. While Missing Completely At Random (MCAR) analysis suggests unbiased exclusion, this substantial missingness limits sample size and may affect representativeness of absolute delay rate estimates (79.4\%).

\paragraph{Weather Imputation Uncertainty:} Despite conservative 4-hour tolerance and validation (relative error < 8\%), nearest-neighbor imputation introduces measurement error. Imputed weather values are treated as certain during model training, potentially underestimating prediction uncertainty. Rapidly changing weather during frontal passages may violate temporal stability assumptions underlying imputation.

\paragraph{Surface Traffic Data Granularity:} Use of representative 24-hour traffic profiles (averaged across 6 months) rather than day-specific measurements eliminates variability from accidents, weather-traffic coupling, and special events. This simplification may have attenuated genuine traffic signal or created spurious stability. However, the negligible observed correlations (r = 0.03) suggest traffic impact is minimal regardless of data granularity.

\paragraph{Limited Temporal Scope:} Five-month observation window (October 2025 – February 2026) captures autumn-winter dynamics but excludes spring-summer patterns. Seasonal regime characterization requires multi-year data spanning complete annual cycles. December's exceptional propagation (+17.10pp) based on limited sample (565 flights) may not represent stable seasonal behavior.

\subsection{Methodological Limitations}

\paragraph{Single Airport Analysis:} Focus on Istanbul Airport limits cross-airport generalization. Hub characteristics (Turkish Airlines dominance 51\%, international traffic mix, infrastructure capacity) differ substantially across airports. Network effects, geographical location, regulatory environments, and operational practices vary, requiring validation studies at additional locations.

\paragraph{Binary Classification Framework:} Predicting delay occurrence ($\geq$15 minutes) rather than magnitude discards information about delay severity. A 20-minute delay and 2-hour delay receive identical treatment, yet operational impacts differ dramatically. Regression models predicting continuous delay duration or ordinal classification (minor/moderate/severe) may provide richer operational insights.

\paragraph{Feature Engineering Constraints:} Features are hand-engineered based on domain knowledge rather than learned representations. Deep learning approaches (autoencoders, LSTMs, transformers) might discover latent patterns not captured by explicit feature construction. However, the 9,527-flight dataset may be insufficient for deep learning's data requirements, and interpretability advantages of tree-based methods favor the current approach for operational deployment.

\paragraph{Temporal Validation Limited to One Test Period:} While temporal validation (Oct-Nov train → Jan-Feb test) exposes regime shift, evaluation on a single chronologically separated period limits generalization assessment. Ideally, multiple non-overlapping temporal test periods or rolling-origin cross-validation would quantify temporal stability more robustly.

\subsection{Scope Limitations}

\paragraph{Departure Delays Only:} Analysis focuses exclusively on departure punctuality. Arrival delays, taxiing times, and gate-to-gate durations are not examined. Network-level delay propagation through connecting flights remains unexplored.

\paragraph{Delay Causation Not Addressed:} Predictive models identify associations but do not establish causality. While weather clearly affects delays, the specific causal pathways (de-icing requirements, runway configuration changes, spacing increases, visibility restrictions) are not decomposed.

\paragraph{No Cost-Benefit Analysis:} Economic impact of delay prediction deployment is not rigorously quantified. While qualitative operational benefits are discussed, formal cost-benefit analysis incorporating intervention costs, delay costs, false alarm costs, and implementation expenses would strengthen practical recommendations.

\section{Recommendations for Future Work}
\label{sec:future_work}

\subsection{Multi-Airport Validation Studies}

Extend temporal validation framework to diverse airport types:
\begin{itemize}
    \item \textbf{Hub vs. Point-to-Point:} Compare regime shift patterns between hub airports (high connectivity, tight schedules) and point-to-point operations (independent flights, looser schedules)
    \item \textbf{Geographical Diversity:} Test seasonal stability across climate zones (tropical, temperate, polar) to determine whether regime shifts are universal or climate-specific
    \item \textbf{Regulatory Environments:} Compare North American, European, and Asian airports to assess whether operational protocols influence temporal stability
\end{itemize}

\subsection{Season-Specific Ensemble Models}

Develop and evaluate ensemble approaches combining season-specific models:
\begin{itemize}
    \item Train separate models for spring, summer, autumn, winter using multi-year data
    \item Implement automated season detection based on weather patterns or calendar periods
    \item Compare ensemble performance against single-model periodic retraining
    \item Investigate optimal transition strategies between seasonal models
\end{itemize}

\subsection{Network Delay Propagation Analysis}

Extend propagation analysis from airport-level to network-level:
\begin{itemize}
    \item Model delay propagation through flight connections and aircraft rotations
    \item Quantify airline network topology effects on delay cascades
    \item Identify critical flights whose delays trigger widespread network disruptions
    \item Develop network-aware prediction models incorporating upstream delay information
\end{itemize}

\subsection{Online Learning and Concept Drift Detection}

Implement adaptive learning systems:
\begin{itemize}
    \item Deploy online learning algorithms (incremental decision trees, adaptive boosting) that continuously update as new data arrives
    \item Develop concept drift detection mechanisms that automatically trigger retraining when feature-target relationships shift
    \item Compare online learning performance against periodic batch retraining
    \item Quantify trade-offs between adaptation speed and overfitting to short-term fluctuations
\end{itemize}

\subsection{Causal Delay Analysis}

Move beyond predictive associations to causal understanding:
\begin{itemize}
    \item Decompose weather impacts into specific causal pathways (de-icing time, runway configuration changes, visibility restrictions, wind limitations)
    \item Apply causal inference methods (instrumental variables, natural experiments, regression discontinuity) to isolate causal effects
    \item Develop mechanistic models combining data-driven predictions with physics-based aviation operations constraints
\end{itemize}

\subsection{Deep Learning Approaches}

Investigate representation learning methods:
\begin{itemize}
    \item LSTM recurrent networks to model temporal delay sequences and capture long-range dependencies
    \item Attention mechanisms to identify critical time windows and features for each prediction
    \item Graph neural networks to represent airport operations as interconnected flight networks
    \item Compare deep learning performance and interpretability against tree-based methods
\end{itemize}

\subsection{Real-Time Deployment Pilot Study}

Conduct controlled operational deployment:
\begin{itemize}
    \item Implement prediction system at partner airport with real-time data feeds
    \item A/B test operational interventions triggered by predictions vs. standard procedures
    \item Quantify actual delay reduction, cost savings, and operational disruption from prediction-based interventions
    \item Measure prediction system trust, adoption, and human override patterns among operations staff
\end{itemize}

\section{Final Conclusions}
\label{sec:final_conclusions}

This research developed and evaluated machine learning models for predicting flight departure delays at Istanbul Airport, achieving state-of-the-art performance (0.903 ROC-AUC) that exceeds published benchmarks under stable temporal conditions. However, the most significant finding is not the strong predictive performance itself, but rather the discovery that this performance catastrophically collapses (to 0.474 AUC) when predicting future seasonal periods due to fundamental non-stationarity in delay propagation mechanisms.

The seasonal regime shift—wherein delay contagion reverses from +7.07 percentage points (autumn) to -7.78 percentage points (winter), representing a 14.84 percentage point swing—reveals that delay generation mechanisms fundamentally differ across seasons. Autumn delays propagate through operational cascades (shared gates, aircraft, crews), while winter delays result from weather-driven independent events that trigger operational resets. This finding challenges the core supervised learning assumption that feature-target relationships remain stable between training and deployment.

Temporal validation using chronologically separated training and test periods exposed this instability, demonstrating that standard random train-test splits—employed by the majority of published aviation delay prediction research—systematically overestimate operational deployment performance by masking seasonal non-stationarity. This methodological insight may be more valuable than the predictive models themselves: it reveals a critical gap in existing literature and establishes temporal validation as an essential requirement for honest assessment of deployment readiness.

The research also conclusively rejects the hypothesis that surface road traffic congestion significantly influences flight delays (mean $|r|$= 0.03, 4.9\% model importance, 0.2 percentage point AUC contribution), demonstrating that hub airport delays are driven overwhelmingly by airside operational constraints rather than landside accessibility. While this represents a negative finding, it constitutes genuine scientific contribution by redirecting future research toward more productive directions.

From a practical deployment perspective, the findings indicate that accurate short-term delay prediction is achievable but requires adaptive strategies to accommodate seasonal regime shifts. Static models trained once and deployed indefinitely will fail. Operational systems must implement continuous or periodic retraining, ensemble methods combining season-specific models, or online learning with concept drift detection. Moreover, deployment strategies should be context-aware, leveraging high-confidence predictions (99.6\% accuracy for p $\geq$ 0.9) for immediate interventions while treating low-confidence forecasts during adverse weather as uncertainty flags warranting conservative contingency planning.

The research demonstrates both the promise and pitfalls of machine learning in aviation operations. Strong performance under stable conditions validates the technical feasibility of delay prediction systems. However, fundamental non-stationarity in delay patterns means that achieving reliable long-term deployment requires not just better models, but fundamentally different approaches—adaptive systems that continuously learn and evolve alongside the dynamic operational environment they aim to predict.

In conclusion, this thesis contributes empirical findings (seasonal regime shift, airline-specific propagation heterogeneity, traffic hypothesis rejection), methodological advances (temporal validation framework, conservative imputation strategy, temporal leakage prevention), and practical insights (deployment requirements, reliability patterns, intervention strategies) that advance both scientific understanding and operational application of machine learning in aviation delay prediction. The discovery of seasonal non-stationarity and its implications for model deployment may prove to be the most enduring contribution—a cautionary finding that honest evaluation must precede optimistic deployment, and that static models are fundamentally insufficient for dynamic operational systems.



\chapter*{References}
\addcontentsline{toc}{chapter}{References}

\begin{list}{}{%
  \setlength{\leftmargin}{1.5em}%
  \setlength{\itemindent}{-1.5em}%
  \setlength{\itemsep}{6pt}%
  \setlength{\parsep}{0pt}%
  \setlength{\topsep}{0pt}%
}

\item Alfarhood, M., Alotaibi, R., Abdulrahim, B., Einieh, A., Almousa, M., \& Alkhanifer, A. (2024).
Predicting flight delays with machine learning: A case study from Saudi Arabian Airlines.
\textit{International Journal of Aerospace Engineering}, \textit{2024}, Article 3385463.
\url{https://doi.org/10.1155/2024/3385463}

\item AviationStack. (2025). \textit{Aviation data API: Real-time and historical flight data}
[Software]. \url{https://aviationstack.com}

\item Choi, S., Kim, Y.\ J., Briceno, S., \& Mavris, D.\ N. (2016).
Prediction of weather-induced airline delays based on machine learning algorithms.
In \textit{Proceedings of the 2016 IEEE/AIAA 35th Digital Avionics Systems Conference (DASC)}
(pp.\ 1--6). IEEE.
\url{https://doi.org/10.1109/DASC.2016.7778035}

\item EUROCONTROL Central Office for Delay Analysis (CODA). (2024).
\textit{CODA digest: All-causes delay to air transport in Europe}.
EUROCONTROL.
\url{https://www.eurocontrol.int/publication/coda-digest-all-causes-delay-air-transport-europe}

\item DHMI (General Directorate of State Airports Authority). (2024).
\textit{Annual airport traffic statistics 2023}. Republic of
T\"{u}rkiye Ministry of Transport and Infrastructure.
\url{https://www.dhmi.gov.tr/en/statistics}


\item Giusti, L., Carvalho, L., Gomes, A.\ T., Coutinho, R., Soares, J., \& Ogasawara, E. (2022).
Analyzing flight delay prediction under concept drift.
\textit{Evolving Systems}, \textit{13}(5), 723--736.
\url{https://doi.org/10.1007/s12530-021-09415-z}

\item Gui, G., Liu, F., Sun, J., Yang, J., Zhou, Z., \& Zhao, D. (2020).
Flight delay prediction based on aviation big data and machine learning.
\textit{IEEE Transactions on Vehicular Technology}, \textit{69}(1), 140--150.
\url{https://doi.org/10.1109/TVT.2019.2954094}

\item International Air Transport Association (IATA). (2023).
\textit{Airline on-time performance: Reporting guidelines and standards}.
IATA.
\url{https://www.iata.org/en/programs/ops-infra/airport/slots/}

\item \. Istanbul B\"{u}y\"{u}k\c{s}ehir Belediyesi (\. IBB) Traffic Management Center. (2025).
\textit{Airport access route traffic profiles: E80/TEM highway and O-3/O-7 ring road
monitoring data} [Data set]. Istanbul Metropolitan Municipality.
\url{https://ibb.istanbul/en}

\item Kim, Y.\ J., Choi, S., Briceno, S., \& Mavris, D.\ N. (2016).
A deep learning approach to flight delay prediction.
In \textit{Proceedings of the 2016 IEEE/AIAA 35th Digital Avionics Systems Conference (DASC)}
(pp.\ 1--5). IEEE.
\url{https://doi.org/10.1109/DASC.2016.7777956}

\item Mamdouh, M., Ezzat, M., \& Hefny, H. (2023).
A novel intelligent approach for flight delay prediction.
\textit{Journal of Big Data}, \textit{10}, Article 158.
\url{https://doi.org/10.1186/s40537-023-00854-w}

\item OpenWeatherMap. (2025). \textit{OpenWeather API: Hourly weather observations and forecast data}
[Software]. \url{https://openweathermap.org/api}

\item Tang, Z., Huang, S., Han, S., \& Peng, J. (2021).
Recent progress about flight delay under complex network.
\textit{Complexity}, \textit{2021}, Article 5513093.
\url{https://doi.org/10.1155/2021/5513093}

\item U.S.\ Department of Transportation. (2020).
\textit{14 CFR Part 234: Airline service quality performance reports}.
Federal Aviation Administration.
\url{https://www.ecfr.gov/current/title-14/chapter-II/subchapter-D/part-234}

\item Yu, B., Guo, Z., Asian, S., Wang, H., \& Chen, G. (2019).
Flight delay prediction for commercial air transport: A deep learning approach.
\textit{Transportation Research Part E: Logistics and Transportation Review},
\textit{125}, 203--221.
\url{https://doi.org/10.1016/j.tre.2019.03.013}





\end{list}









\end{document}
















